\newMonth{Augustus}
\setrunningtitles{Augustus}{August}

% ---- martyrology/mart08/mart0801.htm
\needspace{10\baselineskip}
\begin{paracol}{2}
\selectlanguage{latin}
\begin{center}{\color{gregoriocolor} Kaléndis Augústi. 
 Luna\dots\ }\end{center}
\switchcolumn
\selectlanguage{english}
\begin{center}{\color{gregoriocolor} The   1\textsuperscript{st} Day of 
 August. The\dots\ Day of the Moon.}\end{center}
\end{paracol}

\noindent\begin{tabularx}{\linewidth}{*{19}{>{\centering\arraybackslash}X}}
 \textcolor{gregoriocolor}{a} & \textcolor{gregoriocolor}{b} & \textcolor{gregoriocolor}{c} & \textcolor{gregoriocolor}{d} & \textcolor{gregoriocolor}{e} & \textcolor{gregoriocolor}{f} & \textcolor{gregoriocolor}{g} & \textcolor{gregoriocolor}{h} & \textcolor{gregoriocolor}{i} & \textcolor{gregoriocolor}{k} & \textcolor{gregoriocolor}{l} & \textcolor{gregoriocolor}{m} & \textcolor{gregoriocolor}{n} & \textcolor{gregoriocolor}{p} & \textcolor{gregoriocolor}{q} & \textcolor{gregoriocolor}{r} & \textcolor{gregoriocolor}{s} & \textcolor{gregoriocolor}{t} & \textcolor{gregoriocolor}{u} \\
 7 & 8 & 9 & 10 & 11 & 12 & 13 & 14 & 15 & 16 & 17 & 18 & 19 & 20 & 21 & 22 & 23 & 24 & 25 \\
\end{tabularx}
\vspace{0.5\baselineskip}
\noindent\begin{tabularx}{\linewidth}{*{12}{>{\centering\arraybackslash}X}}
 \textcolor{gregoriocolor}{A} & \textcolor{gregoriocolor}{B} & \textcolor{gregoriocolor}{C} & \textcolor{gregoriocolor}{D} & \textcolor{gregoriocolor}{E} & F & \textcolor{gregoriocolor}{F} & \textcolor{gregoriocolor}{G} & \textcolor{gregoriocolor}{H} & \textcolor{gregoriocolor}{M} & \textcolor{gregoriocolor}{N} & \textcolor{gregoriocolor}{P} \\
 26 & 27 & 28 & 29 & 1 & 2 & 1 & 2 & 3 & 4 & 5 & 6 \\
\end{tabularx}

\begin{paracol}{2}
\selectlanguage{latin}
\lettrine[lines=2]{R}{omæ,} in Exquíliis, 
 Dedicátio sancti Petri Apóstoli ad Víncula.
\switchcolumn
\selectlanguage{english}
\lettrine[lines=2]{A}{t} Rome, on the Esquiline, the 
 Dedication of the Church of St. Peter in Chains.
\switchcolumn*
\selectlanguage{latin}
Antiochíæ pássio 
 sanctórum septem fratrum Machabæórum Mártyrum, qui, cum matre sua, passi 
 sunt sub Antíocho Epíphane Rege. Eórum relíquiæ, Romam translátæ, in 
 eádem Ecclésia sancti Petri ad Víncula cónditæ fuérunt.
\switchcolumn
\selectlanguage{english}
At Antioch, the martyrdom of the 
 seven brothers, the holy Machabees, martyrs, and their mother, who suffered 
 under King Antiochus Epiphanes. Their relics were transferred to Rome, 
 and placed in the church or St. Peter in Chains.
\switchcolumn*
\selectlanguage{latin}
Vercéllis natális 
 sancti Eusébii, Epíscopi et Mártyris, qui, ob confessiónem fídei cathólicæ, 
 a Constántio Príncipe Scythópolim, in Palæstína, et inde in Cappadóciam 
 relegátus, póstmodum ad Ecclésiam suam revérsus, martyrium, persequéntibus 
 Ariánis, passus est. Ipsíus autem memória étiam décimo octávo Kaléndas 
 Januárii solémniter habétur, quo die fuit Epíscopus ordinátus; ejúsque 
 festum recólitur décimo séptimo Kaléndas Januárii.
\switchcolumn
\selectlanguage{english}
At Vercelli, St. Eusebius, bishop 
 and martyr, who, for the confession of the Catholic faith was banished to 
 Scythopolis in Palestine, and thence to Cappadocia, by Emperor Constantine. 
 Afterwards, returning to his church, he suffered martyrdom in the 
 persecution of the Arians. His memory is more especially honoured on 
 the 15th of December, when he was consecrated bishop, and his feast is kept 
 on the 16th of December.
\switchcolumn*
\selectlanguage{latin}
Nucériæ Paganórum, in 
 Campánia, item natális sancti Alfónsi-Maríæ de Ligório, Fundatóris 
 Congregatiónis a sanctíssimo Redemptóre nuncupátæ, Epíscopi sanctæ Agathæ 
 Gothórum et Confessóris, zelo animárum, scriptis, verbo et exémplo insígnis; 
 quem Summus Póntifex Gregórius Décimus sextus albo Sanctórum adscrípsit, et 
 Pius Nonus Doctórem universális Ecclésiæ declarávit, et Pius Duodécimus 
 ómnium Confessariórum ac Moralistárum cæléstem apud Deum Patrónum constítuit. 
 Ipsíus vero festívitas sequénti die celebrátur.
\switchcolumn
\selectlanguage{english}
At Nocera dei Pagani in Campani, the 
 birthday also of St. Alphonsus Maria Liguori, founder of the Congregation of 
 our most Holy Redeemer, bishop of Santa Agata dei Goti, and confessor. 
 Noted for his zeal for souls, his writings, and his example, Pope Gregory 
 XVI added him to the canon of saints, and Pius IX declared him to be a 
 doctor of the Universal Church. Pius XII established him as heavenly 
 patron of all moral theologians and of those who hear Confession. His 
 feast, however, is observed on the day following.
\switchcolumn*
\selectlanguage{latin}
Romæ, via Latína, sanctórum Mártyrum Boni Presbyteri, Fausti et Mauri, cum áliis novem; qui in 
 Actis sancti Stéphani Papæ describúntur.
\switchcolumn
\selectlanguage{english}
At Rome, on the Latin Way, the holy 
 martyrs Bonus, a priest, Faustus and Maur, with nine others, mentioned in 
 the Acts of Pope St. Stephen.
\switchcolumn*
\selectlanguage{latin}
Item Romæ pássio 
 sanctárum Vírginum Fídei, Spei et Caritátis, e sancta Sophía matre 
 progenitárum, quæ, sub Hadriáno Príncipe, martyrii corónam adéptæ sunt.
\switchcolumn
\selectlanguage{english}
Also at Rome, the holy virgins 
 Faith, Hope, and Charity, children of St. Sophia, who won the crown of 
 martyrdom under Emperor Hadrian.
\switchcolumn*
\selectlanguage{latin}
Philadelphíæ, in 
 Arábia, sanctórum Mártyrum Cyrílli, Aquilæ, Petri, Domitiáni, Rufi et 
 Menándri, una die coronatórum.
\switchcolumn
\selectlanguage{english}
At Philadelphia in Arabia, the holy 
 martyrs Cyril, Aquila, Peter, Domitian, Rufus, and Menander, crowned on the 
 same day.
\switchcolumn*
\selectlanguage{latin}
Perge, in Pamphylia, 
 sanctórum Mártyrum Leóntii, Attii, Alexándri, et aliórum sex agricolárum; 
 qui, in persecutióne Diocletiáni, sub Flaviáno Præside, decolláti sunt.
\switchcolumn
\selectlanguage{english}
At Perge in Pamphylia, the holy 
 martyrs Leontius, Attius, Alexander, and six peasants, who were beheaded in 
 the persecution of Diocletian, under the governor Flavian.
\switchcolumn*
\selectlanguage{latin}
Gerúndæ, in Hispánia, 
 natális sancti Felícis Mártyris, qui, post divérsa tormentórum génera, a Daciáno támdiu jussus est laniári, donec invíctum Christo spíritum rédderet.
\switchcolumn
\selectlanguage{english}
At Gerona in Spain, the birthday of 
 the holy martyr Felix. After enduring various torments, by order of 
 Dacian he was cut with knives until he gave his undaunted soul to Christ.
\switchcolumn*
\selectlanguage{latin}
In território Parisiénsi sancti Justíni Mártyris.
\switchcolumn
\selectlanguage{english}
In the diocese of Paris, St. Justin, 
 martyr.
\switchcolumn*
\selectlanguage{latin}
Viénnæ, in Gállia, 
 sancti Veri Epíscopi.
\switchcolumn
\selectlanguage{english}
At Vienne in France, St. Verus, 
 bishop.
\switchcolumn*
\selectlanguage{latin}
Wintóniæ, in Anglia, 
 sancti Ethelwóldi Epíscopi.
\switchcolumn
\selectlanguage{english}
At Winchester in England, St. 
 Ethelwold, bishop.
\switchcolumn*
\selectlanguage{latin}
In pago Lisuíno, in 
 Gállia, sancti Nemésii Confessóris.
\switchcolumn
\selectlanguage{english}
In the country of Lisieux, St. 
 Nemesius, confessor.
\switchcolumn*
\selectlanguage{latin}
\end{paracol}

% \continueMonth{Augustus}

% ---- martyrology/mart08/mart0802.htm
\needspace{10\baselineskip}
\begin{paracol}{2}
\selectlanguage{latin}
\begin{center}{\color{gregoriocolor} Quarto Nonas Augústi. 
 Luna\dots\ }\end{center}
\switchcolumn
\selectlanguage{english}
\begin{center}{\color{gregoriocolor} The   2\textsuperscript{nd} Day of 
 August. The\dots\ Day of the Moon.}\end{center}
\end{paracol}

\noindent\begin{tabularx}{\linewidth}{*{19}{>{\centering\arraybackslash}X}}
 \textcolor{gregoriocolor}{a} & \textcolor{gregoriocolor}{b} & \textcolor{gregoriocolor}{c} & \textcolor{gregoriocolor}{d} & \textcolor{gregoriocolor}{e} & \textcolor{gregoriocolor}{f} & \textcolor{gregoriocolor}{g} & \textcolor{gregoriocolor}{h} & \textcolor{gregoriocolor}{i} & \textcolor{gregoriocolor}{k} & \textcolor{gregoriocolor}{l} & \textcolor{gregoriocolor}{m} & \textcolor{gregoriocolor}{n} & \textcolor{gregoriocolor}{p} & \textcolor{gregoriocolor}{q} & \textcolor{gregoriocolor}{r} & \textcolor{gregoriocolor}{s} & \textcolor{gregoriocolor}{t} & \textcolor{gregoriocolor}{u} \\
 8 & 9 & 10 & 11 & 12 & 13 & 14 & 15 & 16 & 17 & 18 & 19 & 20 & 21 & 22 & 23 & 24 & 25 & 26 \\
\end{tabularx}
\vspace{0.5\baselineskip}
\noindent\begin{tabularx}{\linewidth}{*{12}{>{\centering\arraybackslash}X}}
 \textcolor{gregoriocolor}{A} & \textcolor{gregoriocolor}{B} & \textcolor{gregoriocolor}{C} & \textcolor{gregoriocolor}{D} & \textcolor{gregoriocolor}{E} & F & \textcolor{gregoriocolor}{F} & \textcolor{gregoriocolor}{G} & \textcolor{gregoriocolor}{H} & \textcolor{gregoriocolor}{M} & \textcolor{gregoriocolor}{N} & \textcolor{gregoriocolor}{P} \\
 27 & 28 & 29 & 1 & 2 & 3 & 2 & 3 & 4 & 5 & 6 & 7 \\
\end{tabularx}

\begin{paracol}{2}
\selectlanguage{latin}
\lettrine[lines=2]{S}{ancti} Alfónsi-Maríæ de 
 Ligório, Fundatóris Congregatiónis a sanctíssimo Redemptóre nuncupátæ, 
 Epíscopi sanctæ Agathæ Gothórum, Confessóris et Ecclésiæ Doctóris, qui 
 requiévit in Dómino prídie hujus diéi.
\switchcolumn
\selectlanguage{english}
\lettrine[lines=2]{S}{t.} Alphonsus Maria Liguori, founder 
 of the Congregation of our most Holy Redeemer, bishop of Santa Agata dei 
 Goti, confessor and doctor of the Church, who fell asleep in the Lord on the 
 previous day.
\switchcolumn*
\selectlanguage{latin}
Romæ, in cœmetério 
 Callísti, natális sancti Stéphani Primi, Papæ et Mártyris; qui, in 
 persecutióne Valeriáni, dum Missæ Sacrum perágeret, et, superveniéntibus 
 milítibus, ante altáre intrépidus et immóbilis cœpta mystéria perfíceret, in 
 sede sua decollátus est.
\switchcolumn
\selectlanguage{english}
At Rome, in the cemetery of 
 Callistus, the birthday of St. Stephen I, pope and martyr. In the 
 persecution of Valerian, the soldiers suddenly entered while he was saying 
 Mass, but remaining before the altar, fearless and unmoved, he concluded the 
 sacred mysteries, and was beheaded on his throne.
\switchcolumn*
\selectlanguage{latin}
Nicææ, in Bithynia, 
 pássio sanctæ Theódotæ, cum tribus fíliis suis. Ex his primogénitum, 
 nómine Evódium, cum Christum fiduciáliter confiterétur, fecit primo Nicétius, 
 Consuláris Bithyniæ, fústibus cædi; deínde matrem, cum ómnibus fíliis, igne 
 consúmi.
\switchcolumn
\selectlanguage{english}
At Nicaea in Bithynia, the martyrdom 
 of St. Theodota with her three sons. The eldest named Evodius, 
 confessing Christ with confidence, was first beaten with rods by order of 
 Nicetius, exconsul of Bithynia, and then the mother with all her sons, was 
 consumed by fire.
\switchcolumn*
\selectlanguage{latin}
In Africa sancti 
 Rutílii Mártyris, qui cum sæpius, de loco in locum fúgiens, persecutiónem 
 declinásset, et perículum intérdum étiam pecúnia redemísset, ex inopináto 
 aliquándo comprehénsus est, et, Præsidi oblátus, torméntis cruciátur 
 plúrimis; demum, ígnibus tráditus, egrégio martyrio coronátur.
\switchcolumn
\selectlanguage{english}
In Africa, St. Rutilius, martyr. 
 He had frequently secured safety from the perils of persecution by flight, 
 and sometimes even by means of money, but at last, being unexpectedly 
 apprehended, he was led to the governor and subjected to many tortures. 
 Afterwards he was cast into the fire, and thus merited the glorious crown of 
 martyrdom.
\switchcolumn*
\selectlanguage{latin}
Patávii sancti Máximi, 
 ejúsdem civitátis Epíscopi, qui, miráculis clarus, beáto fine quiévit.
\switchcolumn
\selectlanguage{english}
At Padua, St. Maximus, bishop of 
 that city, who ended his blessed life in peace, with a reputation for 
 miracles.
\switchcolumn*
\selectlanguage{latin}
\end{paracol}


% ---- martyrology/mart08/mart0803.htm
\needspace{10\baselineskip}
\begin{paracol}{2}
\selectlanguage{latin}
\begin{center}{\color{gregoriocolor} Tértio Nonas Augústi. 
 Luna\dots\ }\end{center}
\switchcolumn
\selectlanguage{english}
\begin{center}{\color{gregoriocolor} The   3\textsuperscript{rd} Day of 
 August. The\dots\ Day of the Moon.}\end{center}
\end{paracol}

\noindent\begin{tabularx}{\linewidth}{*{19}{>{\centering\arraybackslash}X}}
 \textcolor{gregoriocolor}{a} & \textcolor{gregoriocolor}{b} & \textcolor{gregoriocolor}{c} & \textcolor{gregoriocolor}{d} & \textcolor{gregoriocolor}{e} & \textcolor{gregoriocolor}{f} & \textcolor{gregoriocolor}{g} & \textcolor{gregoriocolor}{h} & \textcolor{gregoriocolor}{i} & \textcolor{gregoriocolor}{k} & \textcolor{gregoriocolor}{l} & \textcolor{gregoriocolor}{m} & \textcolor{gregoriocolor}{n} & \textcolor{gregoriocolor}{p} & \textcolor{gregoriocolor}{q} & \textcolor{gregoriocolor}{r} & \textcolor{gregoriocolor}{s} & \textcolor{gregoriocolor}{t} & \textcolor{gregoriocolor}{u} \\
 9 & 10 & 11 & 12 & 13 & 14 & 15 & 16 & 17 & 18 & 19 & 20 & 21 & 22 & 23 & 24 & 25 & 26 & 27 \\
\end{tabularx}
\vspace{0.5\baselineskip}
\noindent\begin{tabularx}{\linewidth}{*{12}{>{\centering\arraybackslash}X}}
 \textcolor{gregoriocolor}{A} & \textcolor{gregoriocolor}{B} & \textcolor{gregoriocolor}{C} & \textcolor{gregoriocolor}{D} & \textcolor{gregoriocolor}{E} & F & \textcolor{gregoriocolor}{F} & \textcolor{gregoriocolor}{G} & \textcolor{gregoriocolor}{H} & \textcolor{gregoriocolor}{M} & \textcolor{gregoriocolor}{N} & \textcolor{gregoriocolor}{P} \\
 28 & 29 & 1 & 2 & 3 & 4 & 3 & 4 & 5 & 6 & 7 & 8 \\
\end{tabularx}

\begin{paracol}{2}
\selectlanguage{latin}
\lettrine[lines=2]{H}{ierosólymis} Invéntio 
 beatíssimi Stéphani Protomártyris, et sanctórum Gamaliélis, Nicodémi et 
 Abibónis, sicut Luciáno Presbytero divínitus revelátum est, Honórii 
 Príncipis témpore.
\switchcolumn
\selectlanguage{english}
\lettrine[lines=2]{A}{t} Jerusalem, the finding of the 
 body of blessed Stephen, protomartyr, and of the Saints Gamaliel, Nicodemus, 
 and Abibo, through a divine revelation made to the priest Lucian, in the 
 time of Emperor Honorius.
\switchcolumn*
\selectlanguage{latin}
Philíppis, in 
 Macedónia, sanctæ Lydiæ purpuráriæ, quæ, prædicánte ibídem sancto Paulo 
 Apóstolo, ut beátus Lucas in Actibus Apostolórum refert, ómnium prima 
 crédidit Evangélio.
\switchcolumn
\selectlanguage{english}
At Philippi in Macedonia, St. Lydia, a dealer in purple, 
 who was the first to believe in the Gospel when the apostle St. Paul 
 preached in that city, as is related by St. Luke in the Acts of the Apostles
\switchcolumn*
\selectlanguage{latin}
Neápoli, in Campánia, sancti Aspréni Epíscopi, qui a sancto Petro Apóstolo, curátus ab infirmitáte 
 ac deínde baptizátus, ejúsdem civitátis Epíscopus ordinátus fuit.
\switchcolumn
\selectlanguage{english}
At Naples in Campania, St. Aspren, 
 bishop, who was cured of a sickness by the apostle St. Peter, and after 
 being baptized, was made bishop of that city.
\switchcolumn*
\selectlanguage{latin}
Constantinópoli natális 
 sancti Hermélli Mártyris.
\switchcolumn
\selectlanguage{english}
At Constantinople, the birthday of 
 St. Hermellus, martyr.
\switchcolumn*
\selectlanguage{latin}
Apud Indos, Persis 
 finítimos, pássio sanctórum Monachórum, et aliórum fidélium, quos Abénner 
 Rex, pérsequens Ecclésiam Dei, divérsis afflíctos supplíciis, cædi jussit.
\switchcolumn
\selectlanguage{english}
Among the Indians, bordering on 
 Persia, the martyrdom of holy monks and other Christians who were put to 
 death after suffering diverse torments, during the persecution of the Church 
 of God by King Abenner.
\switchcolumn*
\selectlanguage{latin}
Augustodúni deposítio 
 sancti Euphrónii, Epíscopi et Confessóris.
\switchcolumn
\selectlanguage{english}
At Autun, the death of St. 
 Euphronius, bishop and confessor.
\switchcolumn*
\selectlanguage{latin}
Anágniæ sancti Petri 
 Epíscopi, qui, monástica primum observántia, deínde pastoráli vigilántia 
 clarus, quiévit in Dómino.
\switchcolumn
\selectlanguage{english}
At Anagni, St. Peter, who rested in 
 the Lord after gaining great renown for monastical observance and for 
 pastoral vigilance.
\switchcolumn*
\selectlanguage{latin}
BerϾ, in Syria, 
 sanctárum mulíerum Maránæ et Cyræ.
\switchcolumn
\selectlanguage{english}
At Beroea in Syria, the holy women 
 Marana and Cyra.
\switchcolumn*
\selectlanguage{latin}
\end{paracol}


% ---- martyrology/mart08/mart0804.htm
\needspace{10\baselineskip}
\begin{paracol}{2}
\selectlanguage{latin}
\begin{center}{\color{gregoriocolor} Prídie Nonas Augústi. 
 Luna\dots\ }\end{center}
\switchcolumn
\selectlanguage{english}
\begin{center}{\color{gregoriocolor} The   4\textsuperscript{th} Day of 
 August. The\dots\ Day of the Moon.}\end{center}
\end{paracol}

\noindent\begin{tabularx}{\linewidth}{*{19}{>{\centering\arraybackslash}X}}
 \textcolor{gregoriocolor}{a} & \textcolor{gregoriocolor}{b} & \textcolor{gregoriocolor}{c} & \textcolor{gregoriocolor}{d} & \textcolor{gregoriocolor}{e} & \textcolor{gregoriocolor}{f} & \textcolor{gregoriocolor}{g} & \textcolor{gregoriocolor}{h} & \textcolor{gregoriocolor}{i} & \textcolor{gregoriocolor}{k} & \textcolor{gregoriocolor}{l} & \textcolor{gregoriocolor}{m} & \textcolor{gregoriocolor}{n} & \textcolor{gregoriocolor}{p} & \textcolor{gregoriocolor}{q} & \textcolor{gregoriocolor}{r} & \textcolor{gregoriocolor}{s} & \textcolor{gregoriocolor}{t} & \textcolor{gregoriocolor}{u} \\
 10 & 11 & 12 & 13 & 14 & 15 & 16 & 17 & 18 & 19 & 20 & 21 & 22 & 23 & 24 & 25 & 26 & 27 & 28 \\
\end{tabularx}
\vspace{0.5\baselineskip}
\noindent\begin{tabularx}{\linewidth}{*{12}{>{\centering\arraybackslash}X}}
 \textcolor{gregoriocolor}{A} & \textcolor{gregoriocolor}{B} & \textcolor{gregoriocolor}{C} & \textcolor{gregoriocolor}{D} & \textcolor{gregoriocolor}{E} & F & \textcolor{gregoriocolor}{F} & \textcolor{gregoriocolor}{G} & \textcolor{gregoriocolor}{H} & \textcolor{gregoriocolor}{M} & \textcolor{gregoriocolor}{N} & \textcolor{gregoriocolor}{P} \\
 29 & 1 & 2 & 3 & 4 & 5 & 4 & 5 & 6 & 7 & 8 & 9 \\
\end{tabularx}

\begin{paracol}{2}
\selectlanguage{latin}
\lettrine[lines=2]{S}{ancti} Domínici 
 Confessóris, qui Ordinis Fratrum Prædicatórum Fundátor fuit, atque octávo 
 Idus mensis hujus in pace quiévit.
\switchcolumn
\selectlanguage{english}
\lettrine[lines=2]{S}{t.} Dominic, confessor, founder of 
 the Order of Friars Preachers, who on the sixth day of this month rested in 
 peace.
\switchcolumn*
\selectlanguage{latin}
In vico Ars, diœcésis 
 Bellicénsis, in Gállia, natális sancti Joánnis Baptístæ-Maríæ Vianney, 
 Presbyteri et Confessóris, in parochiáli múnere obeúndo insígnis, quem Pius 
 Papa Undécimus in Sanctórum númerum rétulit; ipsíus festum quinto Idus mensis hujus agéndum indíxit, eúmque ómnium parochórum cæléstem Patrónum 
 constítuit.
\switchcolumn
\selectlanguage{english}
In the village of Ars, in the 
 diocese of Belley, France, the birthday of St. John Baptist-Mary Vianney, priest and 
 confessor, renowned for his devotion as a parish priest. Pope Pius XI 
 placed him in the number of the saints, ordered that his feast should be 
 observed on the 9th day of this month, and appointed him as the heavenly 
 patron of all parish priests.
\switchcolumn*
\selectlanguage{latin}
Thessalonícæ item 
 natális beáti Aristárchi, qui discípulus et comes indivíduus fuit sancti 
 Pauli Apóstoli, de quo ipse Paulus ad Colossénses scribit: « Salútat vos 
 Aristárchus, concaptívus meus ». Is, ab eódem Apóstolo Epíscopus 
 Thessalonicénsium ordinátus, tandem, post longos agónes, sub Neróne, 
 coronátus a Christo, quiévit.
\switchcolumn
\selectlanguage{english}
At Thessalonica, the birthday of 
 blessed Aristarchus, disciple and inseparable companion of the apostle St. 
 Paul, who writes to the Colossians: ``Aristarchus my fellow-prisoner saluteth 
 you.'' He was consecrated bishop of the Thessalonians by the same 
 apostle, and after long sufferings under Nero, crowned by Christ, rested in 
 peace.
\switchcolumn*
\selectlanguage{latin}
Romæ sanctæ Perpétuæ, 
 quæ, a beáto Petro Apóstolo baptizáta, Nazárium fílium et Africánum virum ad 
 Christi fidem perdúxit, et multa sanctórum Mártyrum córpora sepelívit; ac 
 tandem, bonórum óperum méritis cumuláta, migrávit ad Dóminum.
\switchcolumn
\selectlanguage{english}
At Rome, St. Perpetua, who was 
 baptized by the blessed apostle Peter. She converted to the faith her 
 son Nazarius and her husband Africanus, buried the remains of many holy 
 martyrs, and finally went to our Lord endowed with an abundance of merit.
\switchcolumn*
\selectlanguage{latin}
Item Romæ, via Latína, pássio beáti Tertullíni, Presbyteri et Mártyris; qui, sub Valeriáno 
 Imperatóre, post ímpiam fústium mactatiónem, ígnium circa látera exustiónem, 
 oris quassatiónem, atque in equúleo extensiónem nervorúmque cæsiónem, data 
 senténtia, cápitis amputatióne martyrium consummávit.
\switchcolumn
\selectlanguage{english}
At Rome, on the Latin Way, the 
 martyrdom of blessed Tertullinus, priest and martyr, in the time of Emperor 
 Valerian. After being cruelly beaten with rods, after having his sides 
 burned, his mouth shattered; after being stretched on the rack and his limbs 
 crushed, he completed his martyrdom by being beheaded.
\switchcolumn*
\selectlanguage{latin}
Constantinópoli sancti 
 Eleuthérii Mártyris, ex órdine Senatório viri, qui pro Christo, in 
 persecutióne Maximiáni, gládio cæsus est.
\switchcolumn
\selectlanguage{english}
At Constantinople, the holy martyr 
 Eleutherius, of the senatorial rank, who was put to the sword for Christ in 
 the persecution of Maximian.
\switchcolumn*
\selectlanguage{latin}
In Pérside sanctárum 
 Mártyrum Iæ et Sociárum; quæ, cum novem míllibus Christiánis captívis, sub 
 Sápore Rege, divérsis pœnis afflíctæ, martyrium subiérunt.
\switchcolumn
\selectlanguage{english}
In Persia, in the time of King Sapor, 
 the holy martyr Ia and her companions, who, with nine thousand Christian 
 captives, underwent martyrdom after having been subjected to various 
 torments.
\switchcolumn*
\selectlanguage{latin}
Verónæ sancti Agábii, 
 Epíscopi et Confessóris.
\switchcolumn
\selectlanguage{english}
At Verona, St. Agabius, bishop and 
 confessor.
\switchcolumn*
\selectlanguage{latin}
Turónis, in Gállia, 
 sancti Euphrónii Epíscopi.
\switchcolumn
\selectlanguage{english}
At Tours in France, St. Euphronius, bishop.
\switchcolumn*
\selectlanguage{latin}
Colóniæ Agrippínæ 
 commemorátio sancti Protásii Mártyris; qui Medioláni, una cum Gervásio 
 fratre, passus est tertiodécimo Kaléndas Júlii.
\switchcolumn
\selectlanguage{english}
At Cologne, the commemoration of St. 
 Protase, martyr. In company with his brother Gervase, he suffered at 
 Milan on the 19th of June.
\switchcolumn*
\selectlanguage{latin}
\end{paracol}


% ---- martyrology/mart08/mart0805.htm
\needspace{10\baselineskip}
\begin{paracol}{2}
\selectlanguage{latin}
\begin{center}{\color{gregoriocolor} Nonis Augústi. 
 Luna\dots\ }\end{center}
\switchcolumn
\selectlanguage{english}
\begin{center}{\color{gregoriocolor} The   5\textsuperscript{th} Day of 
 August. The\dots\ Day of the Moon.}\end{center}
\end{paracol}

\noindent\begin{tabularx}{\linewidth}{*{19}{>{\centering\arraybackslash}X}}
 \textcolor{gregoriocolor}{a} & \textcolor{gregoriocolor}{b} & \textcolor{gregoriocolor}{c} & \textcolor{gregoriocolor}{d} & \textcolor{gregoriocolor}{e} & \textcolor{gregoriocolor}{f} & \textcolor{gregoriocolor}{g} & \textcolor{gregoriocolor}{h} & \textcolor{gregoriocolor}{i} & \textcolor{gregoriocolor}{k} & \textcolor{gregoriocolor}{l} & \textcolor{gregoriocolor}{m} & \textcolor{gregoriocolor}{n} & \textcolor{gregoriocolor}{p} & \textcolor{gregoriocolor}{q} & \textcolor{gregoriocolor}{r} & \textcolor{gregoriocolor}{s} & \textcolor{gregoriocolor}{t} & \textcolor{gregoriocolor}{u} \\
 11 & 12 & 13 & 14 & 15 & 16 & 17 & 18 & 19 & 20 & 21 & 22 & 23 & 24 & 25 & 26 & 27 & 28 & 29 \\
\end{tabularx}
\vspace{0.5\baselineskip}
\noindent\begin{tabularx}{\linewidth}{*{12}{>{\centering\arraybackslash}X}}
 \textcolor{gregoriocolor}{A} & \textcolor{gregoriocolor}{B} & \textcolor{gregoriocolor}{C} & \textcolor{gregoriocolor}{D} & \textcolor{gregoriocolor}{E} & F & \textcolor{gregoriocolor}{F} & \textcolor{gregoriocolor}{G} & \textcolor{gregoriocolor}{H} & \textcolor{gregoriocolor}{M} & \textcolor{gregoriocolor}{N} & \textcolor{gregoriocolor}{P} \\
 1 & 2 & 3 & 4 & 5 & 6 & 5 & 6 & 7 & 8 & 9 & 10 \\
\end{tabularx}

\begin{paracol}{2}
\selectlanguage{latin}
\lettrine[lines=2]{R}{omæ,} in Exquíliis, 
 Dedicátio Basílicæ sanctæ Maríæ ad Nives.
\switchcolumn
\selectlanguage{english}
\lettrine[lines=2]{A}{t} Rome, on the Esquiline, the 
 Dedication of the Basilica of St. Mary of the Snows.
\switchcolumn*
\selectlanguage{latin}
Cataláuni, in Gállia, 
 sancti Mémmii, civis Románi, qui, a sancto Petro Apóstolo consecrátus illíus 
 civitátis Epíscopus, pópulum sibi commíssum ad Evangélii veritátem perdúxit.
\switchcolumn
\selectlanguage{english}
At Chalons in France, St. Memmius, a 
 Roman citizen, who was consecrated bishop of that city by St. Peter the 
 Apostle, and brought to the truth of the Gospel the people committed to his 
 care.
\switchcolumn*
\selectlanguage{latin}
Romæ pássio sanctórum 
 Mártyrum vigínti trium, qui, in persecutióne Diocletiáni, via Salária véteri, 
 cápite obtruncáti sunt, et, ad clivum Cucúmeris, ibídem sepúlti.
\switchcolumn
\selectlanguage{english}
At Rome, during the persecution of 
 Diocletian, the martyrdom of twenty-three holy martyrs, who were beheaded on 
 the Salarian Way, and buried at the foot of Cucumer Hill.
\switchcolumn*
\selectlanguage{latin}
Asculi, in Picéno, sancti Emygdii, Epíscopi et Mártyris; qui, a sancto Marcéllo Papa Epíscopus 
 ordinátus et illuc ad prædicándum Evangélium missus, ibídem, in confessióne 
 Christi, sub Diocletiáno Imperatóre, martyrii corónam accépit.
\switchcolumn
\selectlanguage{english}
At Ascoli in Piceno, St. Emygdius, 
 bishop and martyr, who was consecrated bishop by Pope St. Marcellus, and 
 sent thither to preach the Gospel. He received the crown of martyrdom 
 for the confession of Christ under Emperor Diocletian.
\switchcolumn*
\selectlanguage{latin}
Antiochíæ sancti 
 Eusígnii mílitis, qui, annum agens centésimum décimum, cum Constantíni Magni 
 fidem, sub quo militáverat, Juliáno Apóstatæ exprobráret, eúmque ut pátriæ 
 pietátis desertórem redargúeret, ab eódem jussus est cápite cædi.
\switchcolumn
\selectlanguage{english}
At Antioch, St. Eusignius, a 
 soldier, who, at the age of one hundred and ten years, because he reproached 
 Julian the Apostate for forsaking the faith of Constantine the Great, under 
 whom he had served, and for having degenerated from his ancestor's piety, 
 was beheaded at his command.
\switchcolumn*
\selectlanguage{latin}
Item sanctórum Mártyrum 
 Ægyptiórum Cantídii, Cantidiáni et Sobélis.
\switchcolumn
\selectlanguage{english}
Also the holy martyrs Cantidius, 
 Cantidian, and Sobel, Egyptians.
\switchcolumn*
\selectlanguage{latin}
Augústæ Vindelicórum 
 natális sanctæ Afræ Mártyris, quæ, cum esset pagána, per doctrínam sancti 
 Narcíssi Epíscopi ad Christum est convérsa, et, cum ómnibus domus suæ 
 membris, ab eódem Epíscopo baptizáta; póstmodum vero, ob Christi 
 confessiónem igni trádita, martyrium suum, septem diébus ántequam beáta 
 Hilária mater ac tres ancíllæ eódem cruciátus génere coronaréntur, felíciter 
 explévit.
\switchcolumn
\selectlanguage{english}
At Augsburg, the birthday of St. 
 Afra, martyr, who being a pagan, was converted to Christ by the teaching of 
 St. Narcissus the bishop, and being baptized with all her household, was 
 given over to the flames for the sake of Christ. Seven days later her 
 mother Hilaria and three handmaids were also crowned by enduring the same 
 kind of torment.
\switchcolumn*
\selectlanguage{latin}
Augustodúni beáti 
 Cassiáni Epíscopi.
\switchcolumn
\selectlanguage{english}
At Autun, blessed Cassian, bishop.
\switchcolumn*
\selectlanguage{latin}
Apud Theánum, in 
 Campánia, sancti Páridis Epíscopi.
\switchcolumn
\selectlanguage{english}
At Teano in Campania, St. Paris, 
 bishop.
\switchcolumn*
\selectlanguage{latin}
In Anglia sancti 
 Oswáldi Regis, cujus gesta sanctus Beda Venerábilis commémorat.
\switchcolumn
\selectlanguage{english}
In England, St. Oswald, king, whose 
 life is related by St. Venerable Bede.
\switchcolumn*
\selectlanguage{latin}
Eódem die sanctæ Nonnæ, 
 quæ fuit mater beatórum Gregórii Nazianzéni, Cæsárii et Gorgoniæ.
\switchcolumn
\selectlanguage{english}
On the same day, St. Nonna, mother 
 of Saints Gregory Nazianzen, Caesarius, and Gorgonia.
\switchcolumn*
\selectlanguage{latin}
\end{paracol}


% ---- martyrology/mart08/mart0806.htm
\needspace{10\baselineskip}
\begin{paracol}{2}
\selectlanguage{latin}
\begin{center}{\color{gregoriocolor} Octávo Idus Augústi. 
 Luna\dots\ }\end{center}
\switchcolumn
\selectlanguage{english}
\begin{center}{\color{gregoriocolor} The   6\textsuperscript{th} Day of 
 August. The\dots\ Day of the Moon.}\end{center}
\end{paracol}

\noindent\begin{tabularx}{\linewidth}{*{19}{>{\centering\arraybackslash}X}}
 \textcolor{gregoriocolor}{a} & \textcolor{gregoriocolor}{b} & \textcolor{gregoriocolor}{c} & \textcolor{gregoriocolor}{d} & \textcolor{gregoriocolor}{e} & \textcolor{gregoriocolor}{f} & \textcolor{gregoriocolor}{g} & \textcolor{gregoriocolor}{h} & \textcolor{gregoriocolor}{i} & \textcolor{gregoriocolor}{k} & \textcolor{gregoriocolor}{l} & \textcolor{gregoriocolor}{m} & \textcolor{gregoriocolor}{n} & \textcolor{gregoriocolor}{p} & \textcolor{gregoriocolor}{q} & \textcolor{gregoriocolor}{r} & \textcolor{gregoriocolor}{s} & \textcolor{gregoriocolor}{t} & \textcolor{gregoriocolor}{u} \\
 12 & 13 & 14 & 15 & 16 & 17 & 18 & 19 & 20 & 21 & 22 & 23 & 24 & 25 & 26 & 27 & 28 & 29 & 1 \\
\end{tabularx}
\vspace{0.5\baselineskip}
\noindent\begin{tabularx}{\linewidth}{*{12}{>{\centering\arraybackslash}X}}
 \textcolor{gregoriocolor}{A} & \textcolor{gregoriocolor}{B} & \textcolor{gregoriocolor}{C} & \textcolor{gregoriocolor}{D} & \textcolor{gregoriocolor}{E} & F & \textcolor{gregoriocolor}{F} & \textcolor{gregoriocolor}{G} & \textcolor{gregoriocolor}{H} & \textcolor{gregoriocolor}{M} & \textcolor{gregoriocolor}{N} & \textcolor{gregoriocolor}{P} \\
 2 & 3 & 4 & 5 & 6 & 7 & 6 & 7 & 8 & 9 & 10 & 11 \\
\end{tabularx}

\begin{paracol}{2}
\selectlanguage{latin}
\lettrine[lines=2]{I}{n} monte Thabor 
 Transfigurátio Dómini nostri Jesu Christi.
\switchcolumn
\selectlanguage{english}
\lettrine[lines=2]{O}{n} Mount Tabor, the Transfiguration 
 of our Lord Jesus Christ.
\switchcolumn*
\selectlanguage{latin}
Romæ, via Appia, in 
 cœmetério Callísti, natális beáti Xysti Secúndi, Papæ et Mártyris, qui in 
 persecutióne Valeriáni, gládio animadvérsus, martyrii corónam accépit.
\switchcolumn
\selectlanguage{english}
At Rome, on the Appian Way, in the 
 cemetery of Callistus, the birthday of blessed Sixtus II, pope and martyr, 
 who received the crown of martyrdom in the persecution of Valerian by being 
 put to the sword.
\switchcolumn*
\selectlanguage{latin}
Item Romæ sanctórum 
 Mártyrum Felicíssimi et Agapíti, ipsíus beáti Xysti Diaconórum; Januárii, 
 Magni, Vincéntii et Stéphani, Subdiaconórum. Hi omnes, una cum eódem 
 Pontífice páriter decolláti sunt, atque in Prætextáti cœmetério sepúlti. 
 Passus est étiam cum eis beátus Quartus, ut scribit sanctus Cypriánus.
\switchcolumn
\selectlanguage{english}
Also, the holy martyrs Felicissimus 
 and Agapitus, deacons of blessed Sixtus; Januarius, Magnus, Vincent, and 
 Stephen, subdeacons, all of whom were beheaded with him and buried in the 
 cemetery of Praetextatus. With them suffered also blessed Quartus, as 
 is related by St. Cyprian.
\switchcolumn*
\selectlanguage{latin}
Bonóniæ natális sancti 
 Domínici Confessóris, qui Ordinis Fratrum Prædicatórum Fundátor éxstitit. 
 Hic vir, sanctitáte et doctrína claríssimus, virginitátem perpétuo illibátam 
 custodívit, et, ob singulárem meritórum grátiam, tres mórtuos suscitávit; cumque prædicatióne sua compressísset hæreses, ac plúrimos ad religiósam et 
 piam vitam instituísset, in pace quiévit. Ejus autem festívitas prídie 
 Nonas mensis hujus celebrátur, ex constitutióne Pauli Papæ Quarti.
\switchcolumn
\selectlanguage{english}
At Bologna, the birthday of St. 
 Dominic, confessor, founder of the Order of Friars Preachers, most renowned 
 for sanctity and learning. He preserved his chastity unsullied to the 
 end of his life, and by his great merits raised three persons from the dead. 
 After having repressed heresies by his preaching, and instructed many in the 
 religious and godly life, he rested in peace. His feast is celebrated 
 on the 4th of August by decree of Pope Paul IV.
\switchcolumn*
\selectlanguage{latin}
In monastério sancti 
 Petri de Cardégna, Ordinis sancti Benedícti, apud Burgos, in Hispánia, 
 pássio ducentórum Monachórum cum Stéphano Abbáte, qui a Saracénis pro Jesu 
 Christi fide interfécti sunt, atque ibídem in claustro a Christiánis sepúlti.
\switchcolumn
\selectlanguage{english}
At Burgos in Spain, in the monastery 
 of St. Peter of Cardegna, of the Order of St. Benedict, two hundred monks, 
 with their abbot Stephen, who were put to death for the faith of Christ by 
 the Saracens, and buried in the monastery by the Christians.
\switchcolumn*
\selectlanguage{latin}
Complúti, in Hispánia, 
 sanctórum Mártyrum Justi et Pastóris fratrum, qui, cum adhuc púeri lítteris 
 imbueréntur, sponte ad martyrium, projéctis in schola tábulis, cucurrérunt; 
 et mox, a Daciáno Præside tenéri jussi et fústibus cædi, ambo, extra 
 civitátem perdúcti sunt, et ibi a carnífice juguláti.
\switchcolumn
\selectlanguage{english}
At Alcala in Spain, the holy martyrs 
 Justus and Pastor, brothers. While they were yet schoolboys, they 
 threw aside their books in school, and spontaneously ran to martyrdom. 
 By order of the governor Dacian, they were arrested, beaten with rods, and 
 as they exhorted each other to constancy, were led out of the city, and had 
 their throats cut by the executioner.
\switchcolumn*
\selectlanguage{latin}
Romæ sancti Hormísdæ, 
 Papæ et Confessóris.
\switchcolumn
\selectlanguage{english}
At Rome, St. Hormisdas, pope and 
 confessor.
\switchcolumn*
\selectlanguage{latin}
Amidæ, in Mesopotámia, 
 sancti Jacóbi Eremítæ, miráculis clari.
\switchcolumn
\selectlanguage{english}
At Amida in Mesopotamia, St. James, a hermit 
 renowned for miracles.
\switchcolumn*
\selectlanguage{latin}
\end{paracol}


% ---- martyrology/mart08/mart0807.htm
\needspace{10\baselineskip}
\begin{paracol}{2}
\selectlanguage{latin}
\begin{center}{\color{gregoriocolor} Séptimo Idus Augústi. 
 Luna\dots\ }\end{center}
\switchcolumn
\selectlanguage{english}
\begin{center}{\color{gregoriocolor} The   7\textsuperscript{th} Day of 
 August. The\dots\ Day of the Moon.}\end{center}
\end{paracol}

\noindent\begin{tabularx}{\linewidth}{*{19}{>{\centering\arraybackslash}X}}
 \textcolor{gregoriocolor}{a} & \textcolor{gregoriocolor}{b} & \textcolor{gregoriocolor}{c} & \textcolor{gregoriocolor}{d} & \textcolor{gregoriocolor}{e} & \textcolor{gregoriocolor}{f} & \textcolor{gregoriocolor}{g} & \textcolor{gregoriocolor}{h} & \textcolor{gregoriocolor}{i} & \textcolor{gregoriocolor}{k} & \textcolor{gregoriocolor}{l} & \textcolor{gregoriocolor}{m} & \textcolor{gregoriocolor}{n} & \textcolor{gregoriocolor}{p} & \textcolor{gregoriocolor}{q} & \textcolor{gregoriocolor}{r} & \textcolor{gregoriocolor}{s} & \textcolor{gregoriocolor}{t} & \textcolor{gregoriocolor}{u} \\
 13 & 14 & 15 & 16 & 17 & 18 & 19 & 20 & 21 & 22 & 23 & 24 & 25 & 26 & 27 & 28 & 29 & 1 & 2 \\
\end{tabularx}
\vspace{0.5\baselineskip}
\noindent\begin{tabularx}{\linewidth}{*{12}{>{\centering\arraybackslash}X}}
 \textcolor{gregoriocolor}{A} & \textcolor{gregoriocolor}{B} & \textcolor{gregoriocolor}{C} & \textcolor{gregoriocolor}{D} & \textcolor{gregoriocolor}{E} & F & \textcolor{gregoriocolor}{F} & \textcolor{gregoriocolor}{G} & \textcolor{gregoriocolor}{H} & \textcolor{gregoriocolor}{M} & \textcolor{gregoriocolor}{N} & \textcolor{gregoriocolor}{P} \\
 3 & 4 & 5 & 6 & 7 & 8 & 7 & 8 & 9 & 10 & 11 & 12 \\
\end{tabularx}

\begin{paracol}{2}
\selectlanguage{latin}
\lettrine[lines=2]{N}{eápoli,} in Campánia, sancti Cajetáni Thienǽi Confessóris, Clericórum Regulárium Fundatóris, qui, 
 singulári in Deum fidúcia, prístinam Apostólicam vivéndi formam suis 
 coléndam trádidit, et, miráculis clarus, a Cleménte Papa Décimo inter 
 Sanctos relátus est.
\switchcolumn
\selectlanguage{english}
\lettrine[lines=2]{A}{t} Naples in Campania, St. Cajetan 
 the Theatine, confessor, founder of the Clerics Regular, who, through 
 singular confidence in God, made his disciples practise the primitive mode 
 of life of the apostles. Being renowned for miracles, he was ranked 
 among the saints by Clement X.
\switchcolumn*
\selectlanguage{latin}
Arétii, in Túscia, 
 natális sancti Donáti, Epíscopi et Mártyris; qui, inter cétera virtútis 
 ópera (ut scribit beátus Gregórius Papa), cálicem sanctum, a Pagánis fractum, 
 orándo instaurávit. Is, in persecutióne Juliáni Apóstatæ, a 
 Quadratiáno Augustáli comprehénsus, et, cum sacrificáre idólis renuísset, 
 gládio percússus, martyrium consummávit. Passus est étiam cum eo 
 beátus Hilarínus Mónachus; cujus memória décimo séptimo Kaléndas Augústi 
 recólitur, quo die sacrum ipsíus corpus ad Ostia Tiberína translátum fuit.
\switchcolumn
\selectlanguage{english}
At Arezzo in Tuscany, the birthday 
 of St. Donatus, bishop and martyr, who among other miraculous deeds by his 
 prayers (as is related by blessed Pope Gregory) made whole again a sacred 
 chalice which had been broken by pagans. Being apprehended by the 
 imperial officer Quadratian, during the persecution of Julian the Apostate, 
 and refusing to sacrifice to idols, he was struck with the sword, and thus 
 fulfilled his martyrdom. With him suffered also the blessed monk 
 Hilarinus, whose feast is celebrated on the 16th of July, at which time his 
 body was taken to Ostia.
\switchcolumn*
\selectlanguage{latin}
Romæ sanctórum Mártyrum 
 Petri et Juliáni, cum áliis decem et octo.
\switchcolumn
\selectlanguage{english}
At Rome, the holy martyrs Peter and 
 Julian, with eighteen others.
\switchcolumn*
\selectlanguage{latin}
Medioláni sancti Fausti 
 mílitis, qui, sub Aurélio Cómmodo, post multa certámina, martyrii palmam 
 adéptus est.
\switchcolumn
\selectlanguage{english}
At Milan, St. Faustus, a soldier, 
 who obtained the palm of martyrdom after many trials in the time of Aurélius 
 Commodus.
\switchcolumn*
\selectlanguage{latin}
Novocómi pássio 
 sanctórum Mártyrum Carpóphori, Exánthi, Cássii, Severíni, Secúndi et Licínii; 
 qui, in confessióne Christi, cápite truncáti sunt.
\switchcolumn
\selectlanguage{english}
At Como, the passion of the holy 
 martyrs Carpophorus, Exanthus, Cassius, Severinus, Secundus, and Licinius, 
 who were beheaded for the confession of Christ.
\switchcolumn*
\selectlanguage{latin}
Nísibi, in Mesopotámia, sancti Dométii, Mónachi Persæ, qui cum duóbus discípulis, sub Juliáno Apóstata, lapidátus est.
\switchcolumn
\selectlanguage{english}
At Nisibis in Mesopotamia, St. 
 Dometius, a Persian monk, who was stoned to death with two of his disciples 
 at the time of Julian the Apostate.
\switchcolumn*
\selectlanguage{latin}
Rotómagi sancti 
 Victrícii Epíscopi, qui adhuc miles, sub eódem Juliáno, abjíciens pro 
 Christo cíngulum, a Tribúno multis torméntis affícitur, et cápitis damnátur; 
 sed, carnífice, qui ad eum cædéndum missus fúerat, cæcitáte percússo, ipse, 
 vínculis solútis, liber evásit. Póstea, Epíscopus factus, indómitas 
 Morinórum et Nerviórum gentes divíni prædicatióne verbi ad Christi fidem 
 perdúxit, et demum Conféssor in pace quiévit.
\switchcolumn
\selectlanguage{english}
At Rouen, the holy bishop St. 
 Victricius. While he was yet a soldier under Julian, he threw away his 
 military belt for Christ, and after being subjected by the tribune to many 
 torments, was condemned to death. But the executioner sent to slay him 
 being struck blind, and the confessor's chains being loosened, he made his 
 escape. Afterwards being made bishop, by preaching the word of God, he 
 brought to the faith of Christ the barbarous people of Belgic Gaul, and 
 finally died in peace, a confessor.
\switchcolumn*
\selectlanguage{latin}
Cataláuni, in Gállia, 
 sancti Donatiáni Epíscopi.
\switchcolumn
\selectlanguage{english}
At Chalons in France, St. Donatian, 
 bishop.
\switchcolumn*
\selectlanguage{latin}
Messánæ, in Sicília, 
 sancti Albérti Confessóris, ex Ordine Carmelitárum, miráculis clari.
\switchcolumn
\selectlanguage{english}
At Messina in Sicily, St. Albert, 
 confessor of the Carmelite Order, renowned for miracles.
\switchcolumn*
\selectlanguage{latin}
\end{paracol}


% ---- martyrology/mart08/mart0808.htm
\needspace{10\baselineskip}
\begin{paracol}{2}
\selectlanguage{latin}
\begin{center}{\color{gregoriocolor} Sexto Idus Augústi. 
 Luna\dots\ }\end{center}
\switchcolumn
\selectlanguage{english}
\begin{center}{\color{gregoriocolor} The 8\textsuperscript{th} Day of 
 August. The\dots\ Day of the Moon.}\end{center}
\end{paracol}

\noindent\begin{tabularx}{\linewidth}{*{19}{>{\centering\arraybackslash}X}}
 \textcolor{gregoriocolor}{a} & \textcolor{gregoriocolor}{b} & \textcolor{gregoriocolor}{c} & \textcolor{gregoriocolor}{d} & \textcolor{gregoriocolor}{e} & \textcolor{gregoriocolor}{f} & \textcolor{gregoriocolor}{g} & \textcolor{gregoriocolor}{h} & \textcolor{gregoriocolor}{i} & \textcolor{gregoriocolor}{k} & \textcolor{gregoriocolor}{l} & \textcolor{gregoriocolor}{m} & \textcolor{gregoriocolor}{n} & \textcolor{gregoriocolor}{p} & \textcolor{gregoriocolor}{q} & \textcolor{gregoriocolor}{r} & \textcolor{gregoriocolor}{s} & \textcolor{gregoriocolor}{t} & \textcolor{gregoriocolor}{u} \\
 14 & 15 & 16 & 17 & 18 & 19 & 20 & 21 & 22 & 23 & 24 & 25 & 26 & 27 & 28 & 29 & 1 & 2 & 3 \\
\end{tabularx}
\vspace{0.5\baselineskip}
\noindent\begin{tabularx}{\linewidth}{*{12}{>{\centering\arraybackslash}X}}
 \textcolor{gregoriocolor}{A} & \textcolor{gregoriocolor}{B} & \textcolor{gregoriocolor}{C} & \textcolor{gregoriocolor}{D} & \textcolor{gregoriocolor}{E} & F & \textcolor{gregoriocolor}{F} & \textcolor{gregoriocolor}{G} & \textcolor{gregoriocolor}{H} & \textcolor{gregoriocolor}{M} & \textcolor{gregoriocolor}{N} & \textcolor{gregoriocolor}{P} \\
 4 & 5 & 6 & 7 & 8 & 9 & 8 & 9 & 10 & 11 & 12 & 13 \\
\end{tabularx}

\begin{paracol}{2}
\selectlanguage{latin}
\lettrine[lines=2]{S}{anctórum} Mártyrum 
 Cyríaci Diáconi, Largi et Smarágdi, qui, cum áliis vigínti Sóciis, passi 
 sunt décimo séptimo Kaléndas Aprílis. Eórum córpora, via Salária a 
 Joánne Presbytero sepúlta, sanctus Marcéllus Papa in prædium Lucínæ, via 
 Ostiénsi, hoc die tránstulit; quæ póstea, in Urbem deláta, in Diaconía 
 sanctæ Maríæ in via Lata fuérunt recóndita.
\switchcolumn
\selectlanguage{english}
\lettrine[lines=2]{T}{he} holy martyrs Cyriacus, deacon, 
 Largus, and Smaragdus, with twenty others who suffered on the 16th of March, 
 during the persecution of Diocletian and Maximian. Their bodies were 
 buried on the Salarian Way by the priest John, but were on this day 
 translated by Pope St. Marcellus to the estate of Lucina, on the Ostian Way. 
 Afterwards they were brought to the city and placed in the church of St. 
 Mary in Via Lata.
\switchcolumn*
\selectlanguage{latin}
Anazárbi, in Cilícia, sancti Maríni senis, qui, sub Diocletiáno Imperatóre et Lysia Præside, cæsus 
 flagris, in ligno suspénsus ac laniátus, feris tandem objéctus intériit.
\switchcolumn
\selectlanguage{english}
At Anzarba in Cilicia, St. Marinus, 
 an old man who was scourged, racked, and lacerated, and who died by being 
 exposed to wild beasts, in the time of Emperor Diocletian and the governor 
 Lysias.
\switchcolumn*
\selectlanguage{latin}
Item sanctórum Mártyrum 
 Eleuthérii et Leónidæ, qui per ignem martyrium consummárunt.
\switchcolumn
\selectlanguage{english}
Also, the holy martyrs Eleutherius 
 and Leonides, who underwent martyrdom by fire.
\switchcolumn*
\selectlanguage{latin}
In Pérside sancti 
 Hormísdæ Mártyris, sub Sápore Rege.
\switchcolumn
\selectlanguage{english}
In Persia, St. Hormisdas, a martyr 
 under King Sapor.
\switchcolumn*
\selectlanguage{latin}
Cyzici, in Hellespónto, 
 sancti Æmiliáni Epíscopi, qui, pro sacrárum Imáginum cultu a Leóne Imperatóre multa passus, demum in exsílio vitam finívit.
\switchcolumn
\selectlanguage{english}
At Cyzicum, on the Hellespont, St. 
 Aemilian, bishop, who ended his life in exile after having suffered much 
 from Emperor Leo for the veneration of holy images.
\switchcolumn*
\selectlanguage{latin}
In Creta sancti Myrónis 
 Epíscopi, miráculis clari.
\switchcolumn
\selectlanguage{english}
In Crete, St. Myron, a bishop 
 renowned for miracles.
\switchcolumn*
\selectlanguage{latin}
Viénnæ, in Gállia, 
 sancti Sevéri, Presbyteri et Confessóris; qui ex India, Evangélii prædicándi 
 causa, laboriósam peregrinatiónem suscépit, et, cum ad præfátam urbem 
 devenísset, ingéntem Paganórum multitúdinem verbo et miráculis ad Christi 
 fidem convértit.
\switchcolumn
\selectlanguage{english}
At Vienne in France, St. Severus, 
 priest and confessor, who undertook a painful journey from India in order to 
 preach the Gospel in that city, and converted a great number of pagans to 
 the faith of Christ by his works and miracles.
\switchcolumn*
\selectlanguage{latin}
\end{paracol}


% ---- martyrology/mart08/mart0809.htm
\needspace{10\baselineskip}
\begin{paracol}{2}
\selectlanguage{latin}
\begin{center}{\color{gregoriocolor} Quinto Idus Augústi. 
 Luna\dots\ }\end{center}
\switchcolumn
\selectlanguage{english}
\begin{center}{\color{gregoriocolor} The 9\textsuperscript{th} Day of 
 August. The\dots\ Day of the Moon.}\end{center}
\end{paracol}

\noindent\begin{tabularx}{\linewidth}{*{19}{>{\centering\arraybackslash}X}}
 \textcolor{gregoriocolor}{a} & \textcolor{gregoriocolor}{b} & \textcolor{gregoriocolor}{c} & \textcolor{gregoriocolor}{d} & \textcolor{gregoriocolor}{e} & \textcolor{gregoriocolor}{f} & \textcolor{gregoriocolor}{g} & \textcolor{gregoriocolor}{h} & \textcolor{gregoriocolor}{i} & \textcolor{gregoriocolor}{k} & \textcolor{gregoriocolor}{l} & \textcolor{gregoriocolor}{m} & \textcolor{gregoriocolor}{n} & \textcolor{gregoriocolor}{p} & \textcolor{gregoriocolor}{q} & \textcolor{gregoriocolor}{r} & \textcolor{gregoriocolor}{s} & \textcolor{gregoriocolor}{t} & \textcolor{gregoriocolor}{u} \\
 15 & 16 & 17 & 18 & 19 & 20 & 21 & 22 & 23 & 24 & 25 & 26 & 27 & 28 & 29 & 1 & 2 & 3 & 4 \\
\end{tabularx}
\vspace{0.5\baselineskip}
\noindent\begin{tabularx}{\linewidth}{*{12}{>{\centering\arraybackslash}X}}
 \textcolor{gregoriocolor}{A} & \textcolor{gregoriocolor}{B} & \textcolor{gregoriocolor}{C} & \textcolor{gregoriocolor}{D} & \textcolor{gregoriocolor}{E} & F & \textcolor{gregoriocolor}{F} & \textcolor{gregoriocolor}{G} & \textcolor{gregoriocolor}{H} & \textcolor{gregoriocolor}{M} & \textcolor{gregoriocolor}{N} & \textcolor{gregoriocolor}{P} \\
 5 & 6 & 7 & 8 & 9 & 10 & 9 & 10 & 11 & 12 & 13 & 14 \\
\end{tabularx}

\begin{paracol}{2}
\selectlanguage{latin}
\lettrine[lines=1]{V}{igília} sancti Lauréntii Mártyris.
\switchcolumn
\selectlanguage{english}
\lettrine[lines=1]{T}{he} vigil of St. Lawrence, martyr.
\switchcolumn*
\selectlanguage{latin}
Sancti Joánnis 
 Baptístæ-Maríæ Vianney, Presbyteri et Confessóris, cæléstis ómnium 
 parochórum Patróni; cujus dies natális prídie Nonas mensis hujus recensétur.
\switchcolumn
\selectlanguage{english}
St. John Baptist-Mary Vianney, 
 priest and confessor, and heavenly patron of all parish priests, whose 
 birthday is remembered on the 4th day of this month.
\switchcolumn*
\selectlanguage{latin}
Romæ sancti Románi, 
 mílitis et Mártyris; qui, confessióne beáti Lauréntii compúnctus, pétiit ab 
 eo baptizári, et, mox exhíbitus ac fústibus cæsus, ad últimum decollátus 
 est.
\switchcolumn
\selectlanguage{english}
At Rome, St. Romanus, a soldier, who 
 was moved by the torments of blessed Lawrence to ask for baptism from him. 
 He was immediately prosecuted, scourged, and finally beheaded.
\switchcolumn*
\selectlanguage{latin}
In Túscia natális 
 sanctórum Mártyrum Secundiáni, Marcelliáni et Veriáni; qui, témpore Décii, a 
 Promóto Consulári primum cæsi sunt, deínde in equúleo suspénsi, et abrási úngulis, atque latéribus appósito assáti, ac tandem triumphálem martyrii 
 palmam, cápite cæsi, meruérunt.
\switchcolumn
\selectlanguage{english}
In Tuscany, the birthday of the holy 
 martyrs Secundian, Marcellian, and Verian. In the time of Decius, they 
 were scourged by the exconsul Promotus, then racked and torn with iron 
 hooks. Being burned with fie applied to their sides, they merited the 
 triumphant palm of martyrdom by being beheaded.
\switchcolumn*
\selectlanguage{latin}
Verónæ sanctórum 
 Mártyrum Firmi et Rústici, qui, témpore Maximiáni Imperatóris, sub Anolíno 
 Júdice, cum sacrificáre idólis renuérunt et constánter in Christi fide persísterent, ambo jussi sunt, post ália superáta supplícia, fústibus cædi 
 et cápite amputári.
\switchcolumn
\selectlanguage{english}
At Verona, the holy martyrs Firmus 
 and Rusticus. When they refused to sacrifice to idols and remained 
 constant in confessing Christ, after they had overcome many other torments, 
 they were condemned to be scourged and beheaded by Anolinus, a judge, during 
 the reign of Emperor Maximian.
\switchcolumn*
\selectlanguage{latin}
In Africa commemorátio 
 plurimórum sanctórum Mártyrum, qui, in persecutióne Valeriáni, hortánte eos 
 ad constántiam sancto Numídico, in ignem conjécti, martyrii palmam adépti 
 sunt. Ipse autem Numídicus, licet cum áliis in rogum injéctus et 
 lapídibus óbrutus fuísset, a fília tamen, effóssus et semivívus repértus, 
 curátus est; ac póstea, ob ejus virtútem, in Ecclésiæ Carthaginénsis 
 Presbyterum a beáto Cypriáno méruit cooptári.
\switchcolumn
\selectlanguage{english}
In Africa, the commemoration of many 
 holy martyrs during the persecution of Valerian. Being exhorted by St. 
 Numidicus, they obtained the palm of martyrdom by being cast into the fire, 
 but Numidicus, although thrown into the flames with the others and 
 overwhelmed with stones, was nevertheless taken out by his daughter. 
 Found half dead, he was restored and deserved afterwards by his virtue to be 
 made priest of the Church of Carthage by blessed Cyprian.
\switchcolumn*
\selectlanguage{latin}
Constantinópoli 
 sanctórum Mártyrum Juliáni, Marciáni et aliórum octo; qui, ob Salvatóris 
 imáginem, quam in porta ǽnea constitúerant, omnes, ímpii Leónis Imperatóris 
 jussu, post multa torménta, gládio necáti sunt.
\switchcolumn
\selectlanguage{english}
At Constantinople, the holy martyrs 
 Julian, Marcian, and eight others. For having set up the image of our 
 Saviour on the brass gate, they were exposed to many torments, and then 
 beheaded by order of the impious emperor Leo.
\switchcolumn*
\selectlanguage{latin}
Cataláuni, in Gállia, 
 sancti Domitiáni, Epíscopi et Confessóris.
\switchcolumn
\selectlanguage{english}
At Chalons in France, St. Domitian, 
 bishop and confessor.
\switchcolumn*
\selectlanguage{latin}
\end{paracol}


% ---- martyrology/mart08/mart0810.htm
\needspace{10\baselineskip}
\begin{paracol}{2}
\selectlanguage{latin}
\begin{center}{\color{gregoriocolor} Quarto Idus Augústi. 
 Luna\dots\ }\end{center}
\switchcolumn
\selectlanguage{english}
\begin{center}{\color{gregoriocolor} The 10\textsuperscript{th} Day of 
 August. The\dots\ Day of the Moon.}\end{center}
\end{paracol}

\noindent\begin{tabularx}{\linewidth}{*{19}{>{\centering\arraybackslash}X}}
 \textcolor{gregoriocolor}{a} & \textcolor{gregoriocolor}{b} & \textcolor{gregoriocolor}{c} & \textcolor{gregoriocolor}{d} & \textcolor{gregoriocolor}{e} & \textcolor{gregoriocolor}{f} & \textcolor{gregoriocolor}{g} & \textcolor{gregoriocolor}{h} & \textcolor{gregoriocolor}{i} & \textcolor{gregoriocolor}{k} & \textcolor{gregoriocolor}{l} & \textcolor{gregoriocolor}{m} & \textcolor{gregoriocolor}{n} & \textcolor{gregoriocolor}{p} & \textcolor{gregoriocolor}{q} & \textcolor{gregoriocolor}{r} & \textcolor{gregoriocolor}{s} & \textcolor{gregoriocolor}{t} & \textcolor{gregoriocolor}{u} \\
 16 & 17 & 18 & 19 & 20 & 21 & 22 & 23 & 24 & 25 & 26 & 27 & 28 & 29 & 1 & 2 & 3 & 4 & 5 \\
\end{tabularx}
\vspace{0.5\baselineskip}
\noindent\begin{tabularx}{\linewidth}{*{12}{>{\centering\arraybackslash}X}}
 \textcolor{gregoriocolor}{A} & \textcolor{gregoriocolor}{B} & \textcolor{gregoriocolor}{C} & \textcolor{gregoriocolor}{D} & \textcolor{gregoriocolor}{E} & F & \textcolor{gregoriocolor}{F} & \textcolor{gregoriocolor}{G} & \textcolor{gregoriocolor}{H} & \textcolor{gregoriocolor}{M} & \textcolor{gregoriocolor}{N} & \textcolor{gregoriocolor}{P} \\
 6 & 7 & 8 & 9 & 10 & 11 & 10 & 11 & 12 & 13 & 14 & 15 \\
\end{tabularx}

\begin{paracol}{2}
\selectlanguage{latin}
\lettrine[lines=2]{R}{omæ,} via Tiburtína, 
 natális beáti Lauréntii Archidiáconi, qui, in persecutióne Valeriáni, post 
 plúrima torménta cárceris, vérberum diversórum, fústium, ac plumbatárum et 
 laminárum ardéntium, ad últimum, in cratícula férrea assátus, martyrium 
 complévit; ejúsque corpus a beáto Hippólyto et Justíno Presbytero sepúltum 
 fuit in cœmetério Cyríacæ, in agro Veráno.
\switchcolumn
\selectlanguage{english}
\lettrine[lines=2]{A}{t} Rome, on the Tiburtine Way, the 
 birthday of the blessed archdeacon Lawrence, martyred during the persecution 
 of Valerian. After much suffering from imprisonment, from scourging 
 with whips set with iron or lead, from hot metal plates, he at last 
 completed his martyrdom by being slowly consumed on an iron instrument made 
 in the form of a gridiron. His body was buried by blessed Hippolytus 
 and the priest Justin in the cemetery of Cyriaca, in the Agro Verano.
\switchcolumn*
\selectlanguage{latin}
In Hispánia Apparítio 
 beátæ Maríæ Vírginis de Mercéde nuncupátæ, quæ Ordinis redemptiónis 
 captivórum sub ejus nómine Institútrix fuit. Ipsíus autem festívitas 
 octávo Kaléndas Octóbris recólitur.
\switchcolumn
\selectlanguage{english}
In Spain, the apparition of the 
 Blessed Virgin Mary under the title of our Lady of Ransom, foundress of the 
 Order for the Redemption of Captives. Her feast is celebrated on the 
 24th of September.
\switchcolumn*
\selectlanguage{latin}
Romæ pássio sanctórum 
 centum sexagínta quinque mílitum Mártyrum, sub Aureliáno Imperatóre.
\switchcolumn
\selectlanguage{english}
At Rome, the passion of one hundred 
 and sixty-five holy martyrs, who were soldiers under Emperor Aurelian.
\switchcolumn*
\selectlanguage{latin}
Alexandríæ commemorátio 
 sanctórum Mártyrum, qui, in persecutióne Valeriáni, sub Præside Æmiliáno, 
 divérsis exquisitísque torméntis diútius cruciáti, vário mortis génere 
 corónam martyrii sunt adépti.
\switchcolumn
\selectlanguage{english}
At Alexandria, the commemoration of 
 the holy martyrs who in the persecution of Valerian, under the governor 
 Emilian, were long tormented with diverse and sharp tortures, and obtained 
 the crown of martyrdom by various kinds of deaths.
\switchcolumn*
\selectlanguage{latin}
Bérgomi sanctæ Astériæ, 
 Vírginis et Mártyris, in persecutióne Diocletiáni et Maximiáni Imperatórum.
\switchcolumn
\selectlanguage{english}
At Bergamo, St. Asteria, virgin and 
 martyr, in the persecution of Emperors Diocletian and Maximian.
\switchcolumn*
\selectlanguage{latin}
Carthágine sanctárum 
 Vírginum et Mártyrum Bassæ, Paulæ et Agathonícæ.
\switchcolumn
\selectlanguage{english}
At Carthage, the holy virgins and 
 martyrs Bassa, Paula, and Agathonica.
\switchcolumn*
\selectlanguage{latin}
Romæ sancti Deúsdedit 
 Confessóris, qui quod in hebdómada mánibus suis operándo lucrabátur, die 
 sábbati paupéribus erogábat.
\switchcolumn
\selectlanguage{english}
At Rome, the holy confessor 
 Deusdedit, a labouring man who gave to the poor every Saturday what he had 
 earned during the week.
\switchcolumn*
\selectlanguage{latin}
\end{paracol}


% ---- martyrology/mart08/mart0811.htm
\needspace{10\baselineskip}
\begin{paracol}{2}
\selectlanguage{latin}
\begin{center}{\color{gregoriocolor} Tértio Idus Augústi. 
 Luna\dots\ }\end{center}
\switchcolumn
\selectlanguage{english}
\begin{center}{\color{gregoriocolor} The 11\textsuperscript{th} Day of 
 August. The\dots\ Day of the Moon.}\end{center}
\end{paracol}

\noindent\begin{tabularx}{\linewidth}{*{19}{>{\centering\arraybackslash}X}}
 \textcolor{gregoriocolor}{a} & \textcolor{gregoriocolor}{b} & \textcolor{gregoriocolor}{c} & \textcolor{gregoriocolor}{d} & \textcolor{gregoriocolor}{e} & \textcolor{gregoriocolor}{f} & \textcolor{gregoriocolor}{g} & \textcolor{gregoriocolor}{h} & \textcolor{gregoriocolor}{i} & \textcolor{gregoriocolor}{k} & \textcolor{gregoriocolor}{l} & \textcolor{gregoriocolor}{m} & \textcolor{gregoriocolor}{n} & \textcolor{gregoriocolor}{p} & \textcolor{gregoriocolor}{q} & \textcolor{gregoriocolor}{r} & \textcolor{gregoriocolor}{s} & \textcolor{gregoriocolor}{t} & \textcolor{gregoriocolor}{u} \\
 17 & 18 & 19 & 20 & 21 & 22 & 23 & 24 & 25 & 26 & 27 & 28 & 29 & 1 & 2 & 3 & 4 & 5 & 6 \\
\end{tabularx}
\vspace{0.5\baselineskip}
\noindent\begin{tabularx}{\linewidth}{*{12}{>{\centering\arraybackslash}X}}
 \textcolor{gregoriocolor}{A} & \textcolor{gregoriocolor}{B} & \textcolor{gregoriocolor}{C} & \textcolor{gregoriocolor}{D} & \textcolor{gregoriocolor}{E} & F & \textcolor{gregoriocolor}{F} & \textcolor{gregoriocolor}{G} & \textcolor{gregoriocolor}{H} & \textcolor{gregoriocolor}{M} & \textcolor{gregoriocolor}{N} & \textcolor{gregoriocolor}{P} \\
 7 & 8 & 9 & 10 & 11 & 12 & 11 & 12 & 13 & 14 & 15 & 16 \\
\end{tabularx}

\begin{paracol}{2}
\selectlanguage{latin}
\lettrine[lines=2]{R}{omæ,} inter duas Lauros, 
 natális sancti Tibúrtii Mártyris, qui, sub Júdice Fabiáno, in persecutióne 
 Diocletiáni, cum Christum, nudis plántis super carbónes ardéntes ingréssus, 
 majóri confiterétur constántia, duci in tértium ab Urbe milliárium atque 
 ibídem gládio animadvérti jubétur.
\switchcolumn
\selectlanguage{english}
\lettrine[lines=2]{A}{t} Rome, between the two laurels 
 situation about three miles from the city, the birthday of St. Tiburtius, 
 martyr, under the judge Fabian, in the persecution of Diocletian. 
 After he had walked barefooted on burning coals and confessed Christ with 
 increased constancy, he was put to the sword.
\switchcolumn*
\selectlanguage{latin}
Item Romæ sanctæ 
 Susánnæ Vírginis, quæ, cum ex nóbili prosápia esset orta et beáti Caji 
 Pontíficis neptis, martyrii palmam, témpore Diocletiáni, cápitis 
 obtruncatióne proméruit.
\switchcolumn
\selectlanguage{english}
Also at Rome, the holy virgin 
 Susanna, a woman of noble race, and niece of the blessed Pontiff Caius. 
 She merited the palm of martyrdom by being beheaded in the time of 
 Diocletian.
\switchcolumn*
\selectlanguage{latin}
Assísii in Umbria, 
 natális sanctæ Claræ Vírginis, primæ plantæ Páuperum Dominárum Ordinis 
 Minórum; quam, vita et miráculis célebrem, Alexánder Papa Quartus in númerum 
 sanctárum Vírginum rétulit. Ipsíus tamen festum sequénti die 
 celebrátur.
\switchcolumn
\selectlanguage{english}
At Assisi in Umbria, the birthday of 
 St. Clare, virgin, the first of the Poor Ladies of the Order of Friars 
 Minor. Being celebrated for holiness of life and miracles, she was 
 placed among the holy virgins by Pope Alexander IV. Her feast, 
 however, is observed on the day following.
\switchcolumn*
\selectlanguage{latin}
Cománæ, in Ponto, 
 sancti Alexándri Epíscopi, cognoménto Carbonárii, qui, ex philósopho disertíssimo eminéntem Christiánæ humilitátis sciéntiam adéptus, et a sancto 
 Gregório Thaumatúrgo in thronum illíus Ecclésiæ sublimátus, non solum 
 prædicatióne, sed étiam consummáto per ignem martyrio fuit illústris.
\switchcolumn
\selectlanguage{english}
At Comana in Pontus, St. Alexander, 
 bishop, surnamed Carbonarius, who added to a masterful knowledge of 
 philosophy an eminent degree of Christian humility. He was promoted to 
 the See of that church by St. Gregory Thaumaturgus, and became illustrious, 
 not only by preaching, but also by suffering martyrdom by fire.
\switchcolumn*
\selectlanguage{latin}
Eódem die pássio 
 sanctórum Rufíni, Marsórum Epíscopi, et Sociórum ejus, sub Maximiáno 
 Imperatóre.
\switchcolumn
\selectlanguage{english}
The same day, the martyrdom of St. 
 Rufinus, Bishop of the Marsi, and his companions, under Emperor Maximinus.
\switchcolumn*
\selectlanguage{latin}
Apud Ebroicénses, in 
 Gállia, sancti Tauríni Epíscopi, qui, a beáto Cleménte Papa ordinátus illíus 
 civitátis Epíscopus, Evangélii prædicatióne Christiánam fidem propagávit, 
 ac, multis pro ea suscéptis labóribus, miraculórum glória conspícuus 
 obdormívit in Dómino.
\switchcolumn
\selectlanguage{english}
At Evreux in France, St. Thaurinus, 
 bishop. Being made bishop of that city by blessed Pope Clement, he 
 propagated the Christian faith by the preaching of the Gospel, and the many 
 labours he sustained for it. Celebrated for glorious miracles, he fell 
 asleep in the Lord.
\switchcolumn*
\selectlanguage{latin}
Cameráci, in Gállia, 
 sancti Gaugeríci, Epíscopi et Confessóris.
\switchcolumn
\selectlanguage{english}
At Cambrai in France, St. Gaugericus, 
 bishop and confessor.
\switchcolumn*
\selectlanguage{latin}
In província Valériæ 
 sancti Equítii Abbátis, cujus sánctitas testimónio beáti Gregórii Papæ 
 comprobátur.
\switchcolumn
\selectlanguage{english}
In the province of Valeria, St. 
 Equitius, abbot, whose sanctity is attested by blessed Pope Gregory.
\switchcolumn*
\selectlanguage{latin}
Tudérti, in Umbria, 
 sanctæ Dignæ Vírginis.
\switchcolumn
\selectlanguage{english}
At Todi in Umbria, St. Digna, 
 virgin.
\switchcolumn*
\selectlanguage{latin}
\end{paracol}


% ---- martyrology/mart08/mart0812.htm
\needspace{10\baselineskip}
\begin{paracol}{2}
\selectlanguage{latin}
\begin{center}{\color{gregoriocolor} Prídie Idus Augústi. 
 Luna\dots\ }\end{center}
\switchcolumn
\selectlanguage{english}
\begin{center}{\color{gregoriocolor} The 12\textsuperscript{th} Day of 
 August. The\dots\ Day of the Moon.}\end{center}
\end{paracol}

\noindent\begin{tabularx}{\linewidth}{*{19}{>{\centering\arraybackslash}X}}
 \textcolor{gregoriocolor}{a} & \textcolor{gregoriocolor}{b} & \textcolor{gregoriocolor}{c} & \textcolor{gregoriocolor}{d} & \textcolor{gregoriocolor}{e} & \textcolor{gregoriocolor}{f} & \textcolor{gregoriocolor}{g} & \textcolor{gregoriocolor}{h} & \textcolor{gregoriocolor}{i} & \textcolor{gregoriocolor}{k} & \textcolor{gregoriocolor}{l} & \textcolor{gregoriocolor}{m} & \textcolor{gregoriocolor}{n} & \textcolor{gregoriocolor}{p} & \textcolor{gregoriocolor}{q} & \textcolor{gregoriocolor}{r} & \textcolor{gregoriocolor}{s} & \textcolor{gregoriocolor}{t} & \textcolor{gregoriocolor}{u} \\
 18 & 19 & 20 & 21 & 22 & 23 & 24 & 25 & 26 & 27 & 28 & 29 & 1 & 2 & 3 & 4 & 5 & 6 & 7 \\
\end{tabularx}
\vspace{0.5\baselineskip}
\noindent\begin{tabularx}{\linewidth}{*{12}{>{\centering\arraybackslash}X}}
 \textcolor{gregoriocolor}{A} & \textcolor{gregoriocolor}{B} & \textcolor{gregoriocolor}{C} & \textcolor{gregoriocolor}{D} & \textcolor{gregoriocolor}{E} & F & \textcolor{gregoriocolor}{F} & \textcolor{gregoriocolor}{G} & \textcolor{gregoriocolor}{H} & \textcolor{gregoriocolor}{M} & \textcolor{gregoriocolor}{N} & \textcolor{gregoriocolor}{P} \\
 8 & 9 & 10 & 11 & 12 & 13 & 12 & 13 & 14 & 15 & 16 & 17 \\
\end{tabularx}

\begin{paracol}{2}
\selectlanguage{latin}
\lettrine[lines=2]{S}{anctæ} Claræ Vírginis, 
 primæ plantæ Páuperum Dominárum Ordinis Minórum; quæ ad ætérnas Agni núptias 
 evocáta est prídie hujus diéi.
\switchcolumn
\selectlanguage{english}
\lettrine[lines=2]{S}{t.} Clare, virgin, the first fruits 
 of the Poor Ladies of the Order of Friars Minor, who was called to the 
 everlasting nuptials of the Lamb on the day previous.
\switchcolumn*
\selectlanguage{latin}
Eódem die sanctórum 
 Mártyrum Porcárii, Abbátis monastérii Lirinénsis, et Sociórum ejus 
 quingentórum Monachórum; qui pro fide cathólica, a bárbaris cæsi, martyrio 
 coronáti sunt.
\switchcolumn
\selectlanguage{english}
The same day, the holy martyrs 
 Porcarius, abbot of the monastery of Lerins, and five hundred monks, who 
 were slain for the Catholic faith by barbarians, and were thus crowned with 
 martyrdom.
\switchcolumn*
\selectlanguage{latin}
Catánæ, in Sicília, 
 natális sancti Euplii Diáconi, sub Diocletiáno et Maximiáno Augústis; qui, 
 cum diutíssime pro confessióne Dómini tortus esset, tandem martyrii palmam, 
 gládio cædénte, percépit.
\switchcolumn
\selectlanguage{english}
At Catania in Sicily, the birthday 
 of St. Euplius, deacon, under Emperors Diocletian and Maximian. He was 
 long tortured for the confession of the Lord, and finally obtained the palm 
 of martyrdom by being put to the sword.
\switchcolumn*
\selectlanguage{latin}
Augústæ Vindelicórum 
 sanctæ Hiláriæ, quæ, cum esset beátæ Afræ Mártyris mater et ad sepúlcrum illíus excubáret, ibídem, pro fide Christi, a persecutóribus igni trádita 
 est cum Digna, Euprépia et Eunómia, ancíllis suis. Passi sunt étiam 
 eódem die, in præfáta urbe, Quiríacus, Lárgio, Crescentiánus, Nímmia et 
 Juliána, cum áliis vigínti.
\switchcolumn
\selectlanguage{english}
At Augsburg, St. Hilaria, mother of 
 the blessed martyr Afra. Because she watched at the tomb of her 
 daughter she was cast into the fire for the faith of Christ, together with 
 her maidservants Digna, Euprepia, and Eunomia. On the same day there 
 suffered also in that city Quiriacus, Largius, Crescentian, Nimmia, and 
 Juliana, with twenty others.
\switchcolumn*
\selectlanguage{latin}
In Syria sanctórum 
 Mártyrum Macárii et Juliáni.
\switchcolumn
\selectlanguage{english}
In Syria, the holy martyrs Marcarius 
 and Julian.
\switchcolumn*
\selectlanguage{latin}
Nicomedíæ sanctórum 
 Mártyrum Anicéti Cómitis, et Photíni fratris, cum áliis plúribus, sub 
 Diocletiáno Imperatóre.
\switchcolumn
\selectlanguage{english}
At Nicomedia, the holy martyrs Count 
 Anicetus and his brother Photinus, along with many others, under Emperor 
 Diocletian.
\switchcolumn*
\selectlanguage{latin}
Falériæ, in Túscia, 
 pássio sanctórum Graciliáni, et Felicíssimæ Vírginis, quorum pro fídei 
 confessióne ora lapídibus primo contúsa sunt; dehinc ambo, gládio percússi, 
 optátam martyrii palmam suscepérunt.
\switchcolumn
\selectlanguage{english}
At Faleria in Tuscany, the Saints 
 Gracilian, and Felicíssima, virgin, who, for the confession of the faith, 
 first had their mouths bruised with stones, and being afterwards struck with 
 the sword, received the palm of martyrdom.
\switchcolumn*
\selectlanguage{latin}
Medioláni deposítio 
 sancti Eusébii, Epíscopi et Confessóris.
\switchcolumn
\selectlanguage{english}
At Milan, the death of St. Eusebius, bishop and 
 confessor.
\switchcolumn*
\selectlanguage{latin}
Bríxiæ sancti Herculáni 
 Epíscopi.
\switchcolumn
\selectlanguage{english}
At Brescia, St. Herculanus, bishop.
\switchcolumn*
\selectlanguage{latin}
\end{paracol}


% ---- martyrology/mart08/mart0813.htm
\needspace{10\baselineskip}
\begin{paracol}{2}
\selectlanguage{latin}
\begin{center}{\color{gregoriocolor} Idibus Augústi. 
 Luna\dots\ }\end{center}
\switchcolumn
\selectlanguage{english}
\begin{center}{\color{gregoriocolor} The 13\textsuperscript{th} Day of 
 August. The\dots\ Day of the Moon.}\end{center}
\end{paracol}

\noindent\begin{tabularx}{\linewidth}{*{19}{>{\centering\arraybackslash}X}}
 \textcolor{gregoriocolor}{a} & \textcolor{gregoriocolor}{b} & \textcolor{gregoriocolor}{c} & \textcolor{gregoriocolor}{d} & \textcolor{gregoriocolor}{e} & \textcolor{gregoriocolor}{f} & \textcolor{gregoriocolor}{g} & \textcolor{gregoriocolor}{h} & \textcolor{gregoriocolor}{i} & \textcolor{gregoriocolor}{k} & \textcolor{gregoriocolor}{l} & \textcolor{gregoriocolor}{m} & \textcolor{gregoriocolor}{n} & \textcolor{gregoriocolor}{p} & \textcolor{gregoriocolor}{q} & \textcolor{gregoriocolor}{r} & \textcolor{gregoriocolor}{s} & \textcolor{gregoriocolor}{t} & \textcolor{gregoriocolor}{u} \\
 19 & 20 & 21 & 22 & 23 & 24 & 25 & 26 & 27 & 28 & 29 & 1 & 2 & 3 & 4 & 5 & 6 & 7 & 8 \\
\end{tabularx}
\vspace{0.5\baselineskip}
\noindent\begin{tabularx}{\linewidth}{*{12}{>{\centering\arraybackslash}X}}
 \textcolor{gregoriocolor}{A} & \textcolor{gregoriocolor}{B} & \textcolor{gregoriocolor}{C} & \textcolor{gregoriocolor}{D} & \textcolor{gregoriocolor}{E} & F & \textcolor{gregoriocolor}{F} & \textcolor{gregoriocolor}{G} & \textcolor{gregoriocolor}{H} & \textcolor{gregoriocolor}{M} & \textcolor{gregoriocolor}{N} & \textcolor{gregoriocolor}{P} \\
 9 & 10 & 11 & 12 & 13 & 14 & 13 & 14 & 15 & 16 & 17 & 18 \\
\end{tabularx}

\begin{paracol}{2}
\selectlanguage{latin}
\lettrine[lines=2]{R}{omæ} beáti Hippólyti 
 Mártyris, qui pro confessióne glória, sub Valeriáno Imperatóre, post ália 
 torménta, ligátis pédibus ad colla indomitórum equórum, per carduétum et 
 tríbulos crudéliter tractus est, ac, toto córpore laceráto, emísit spíritum. 
 Passi sunt étiam eódem die beáta Concórdia, ejus nutrix, quæ ante ipsum, 
 plumbátis cæsa, migrávit ad Dóminum; et álii decem et novem de domo sua, qui 
 extra portam Tiburtínam decolláti sunt, et, una cum ipso, in agro Veráno 
 sepúlti.
\switchcolumn
\selectlanguage{english}
\lettrine[lines=2]{A}{t} Rome, the blessed Hippolytus, 
 martyr, who gloriously confessed the faith, under Emperor Valerian. 
 After enduring other torments, he was tied by the feet to the necks of wild 
 horses, and being cruelly dragged through briars and brambles, and having 
 all his body lacerated, he yielded up his spirit. On the same day 
 suffered also blessed Concordia, his nurse, who being scourged in his 
 presence with leaded whips, went to our Lord, and nineteen others of his 
 household, who were beheaded beyond the Tiburtine Gate, and buried with him 
 in the Agro Verano.
\switchcolumn*
\selectlanguage{latin}
Apud Forum Sillæ 
 natális sancti Cassiáni Mártyris, qui cum adoráre idóla noluísset, ídeo, 
 vocátis a persecutóre púeris, quibus exósus docéndo factus fúerat, data est 
 eis facúltas eum periméndi; quorum quanto infírmior erat manus, tanto 
 graviórem martyrii pœnam, diláta morte, faciébat.
\switchcolumn
\selectlanguage{english}
At Imola, the birthday of St. 
 Cassian, martyr. As he refused to worship idols, the persecutor called 
 the boys whom the saint had taught and who hated him, giving them leave to 
 kill him. The torment suffered by the martyr was the more grievous, as 
 the hands which inflicted it, by reason of weakness, rendered death long 
 drawn-out.
\switchcolumn*
\selectlanguage{latin}
Tudérti, in Umbria, 
 sancti Cassiáni, Epíscopi et Mártyris, sub Diocletiáno Imperatóre.
\switchcolumn
\selectlanguage{english}
At Todi in Umbria, St. Cassian, 
 bishop and martyr, under Emperor Diocletian.
\switchcolumn*
\selectlanguage{latin}
Burgis, in Hispánia, 
 sanctárum Centóllæ et Hélenæ Mártyrum.
\switchcolumn
\selectlanguage{english}
At Burgos in Spain, Saints Centolla 
 and Helena, martyrs.
\switchcolumn*
\selectlanguage{latin}
Constantinópoli sancti 
 Máximi Abbátis, doctrína et cathólicæ veritátis zelo insígnis; qui, cum 
 advérsus Monothelítas strénue decertáret, ab hærético Imperatóre Constánte, 
 præcísis mánibus ac lingua, in Chersonésum relegátus est, ibíque, glória 
 confessiónis clarus, spíritum Deo réddidit. Tunc étiam duo Anastásii, 
 qui ejus erant discípuli, aliíque plures divérsa torménta et dura exsília 
 sunt expérti.
\switchcolumn
\selectlanguage{english}
At Constantinople, St. Maximus, a 
 monk distinguished for learning and for zeal for Catholic truth. 
 Valiantly disputing the Monothelites, he had his hands and tongue torn from 
 him by the heretical emperor Constans, and was banished to Chersonesus, 
 where he breathed his last. At this time, two of his disciples, both 
 named Anastasius, and many others endured divers torments and the hardships 
 of exile.
\switchcolumn*
\selectlanguage{latin}
Fritesláriæ, in 
 Germánia, sancti Wigbérti, Presbyteri et Confessóris.
\switchcolumn
\selectlanguage{english}
At Fritzlar in Germany, St. Wigbert, 
 priest and confessor.
\switchcolumn*
\selectlanguage{latin}
Romæ natális sancti 
 Joánnis Berchmans, scholástici e Societáte Jesu et Confessóris, vitæ 
 innocéntia et religiósæ disciplínæ custódia præclári; cui Leo Décimus 
 tértius, Póntifex Máximus, cælitum Sanctórum honóres decrévit.
\switchcolumn
\selectlanguage{english}
At Rome, the birthday of St. John 
 Berchmans, a scholastic of the Society of Jesus, illustrious for his 
 innocence and for his fidelity to the rules of the religious life. He 
 was canonized by Pope Leo XIII.
\switchcolumn*
\selectlanguage{latin}
Pictávis, in Gállia, 
 sanctæ Radegúndis Regínæ, cujus vita miráculis et virtútibus cláruit.
\switchcolumn
\selectlanguage{english}
At Poitiers in France, St. Radegund, 
 queen, whose life was renowned for miracles and virtues.
\switchcolumn*
\selectlanguage{latin}
\end{paracol}


% ---- martyrology/mart08/mart0814.htm
\needspace{10\baselineskip}
\begin{paracol}{2}
\selectlanguage{latin}
\begin{center}{\color{gregoriocolor} Décimo nono Kaléndas Septémbris. 
 Luna\dots\ }\end{center}
\switchcolumn
\selectlanguage{english}
\begin{center}{\color{gregoriocolor} The 
 14\textsuperscript{th} Day of 
 August. The\dots\ Day of the Moon.}\end{center}
\end{paracol}

\noindent\begin{tabularx}{\linewidth}{*{19}{>{\centering\arraybackslash}X}}
 \textcolor{gregoriocolor}{a} & \textcolor{gregoriocolor}{b} & \textcolor{gregoriocolor}{c} & \textcolor{gregoriocolor}{d} & \textcolor{gregoriocolor}{e} & \textcolor{gregoriocolor}{f} & \textcolor{gregoriocolor}{g} & \textcolor{gregoriocolor}{h} & \textcolor{gregoriocolor}{i} & \textcolor{gregoriocolor}{k} & \textcolor{gregoriocolor}{l} & \textcolor{gregoriocolor}{m} & \textcolor{gregoriocolor}{n} & \textcolor{gregoriocolor}{p} & \textcolor{gregoriocolor}{q} & \textcolor{gregoriocolor}{r} & \textcolor{gregoriocolor}{s} & \textcolor{gregoriocolor}{t} & \textcolor{gregoriocolor}{u} \\
 20 & 21 & 22 & 23 & 24 & 25 & 26 & 27 & 28 & 29 & 1 & 2 & 3 & 4 & 5 & 6 & 7 & 8 & 9 \\
\end{tabularx}
\vspace{0.5\baselineskip}
\noindent\begin{tabularx}{\linewidth}{*{12}{>{\centering\arraybackslash}X}}
 \textcolor{gregoriocolor}{A} & \textcolor{gregoriocolor}{B} & \textcolor{gregoriocolor}{C} & \textcolor{gregoriocolor}{D} & \textcolor{gregoriocolor}{E} & F & \textcolor{gregoriocolor}{F} & \textcolor{gregoriocolor}{G} & \textcolor{gregoriocolor}{H} & \textcolor{gregoriocolor}{M} & \textcolor{gregoriocolor}{N} & \textcolor{gregoriocolor}{P} \\
 10 & 11 & 12 & 13 & 14 & 15 & 14 & 15 & 16 & 17 & 18 & 19 \\
\end{tabularx}

\begin{paracol}{2}
\selectlanguage{latin}
\lettrine[lines=2]{V}{igília} Assumptiónis 
 beátæ Maríæ Vírginis.
\switchcolumn
\selectlanguage{english}
\lettrine[lines=2]{T}{he} Vigil of the Assumption of the 
 Blessed Virgin Mary.
\switchcolumn*
\selectlanguage{latin}
Romæ natális beáti 
 Eusébii, Presbyteri et Confessóris, qui, ab Ariáno Imperatóre Constántio, ob 
 cathólicæ fídei defensiónem, in quodam domus suæ cubículo inclúsus, ibi, cum 
 menses septem in oratióne constánter perseverásset, dormitiónem accépit. 
 Ipsíus autem corpus collegérunt Gregórius et Orósius Presbyteri, et in 
 cœmetério Callísti, via Appia, sepeliérunt.
\switchcolumn
\selectlanguage{english}
At Rome, the birthday of the blessed 
 priest Eusebius, who for the defence of the Catholic faith was shut up in a 
 room of his own house by the Arian emperor Constantius, where constantly 
 persevering in prayer for seven months, he rested in peace. His body 
 was removed by the priests Gregory and Orosius, and buried in the cemetery 
 of Callistus, on the Appian Way.
\switchcolumn*
\selectlanguage{latin}
Apaméæ, in Syria, 
 sancti Marcélli, Epíscopi et Mártyris; qui, cum Jovis delúbrum diruísset, a 
 furéntibus Gentílibus occísus est.
\switchcolumn
\selectlanguage{english}
At Apamea in Syria, St. Marcellus, 
 bishop and martyr, who was killed by the enraged heathen for having pulled 
 down a temple of Jupiter.
\switchcolumn*
\selectlanguage{latin}
Tudérti, in Umbria, 
 sancti Callísti, Epíscopi et Mártyris.
\switchcolumn
\selectlanguage{english}
At Todi in Umbria, St. Callistus, 
 bishop and martyr.
\switchcolumn*
\selectlanguage{latin}
In Illyrico sancti 
 Ursícii Mártyris, qui, sub Maximiáno Imperatóre et Aristíde Præside, post 
 multa et divérsa torménta, pro Christi nómine, gládio cæsus est.
\switchcolumn
\selectlanguage{english}
In Illyria, St. Ursicius, martyr, 
 who was beheaded for Christ after suffering various torments under Emperor 
 Maximian and the governor Aristides.
\switchcolumn*
\selectlanguage{latin}
In Africa sancti 
 Demétrii Mártyris.
\switchcolumn
\selectlanguage{english}
In Africa, St. Demetrius, martyr.
\switchcolumn*
\selectlanguage{latin}
In Ægína ínsula sanctæ 
 Athanásiæ Víduæ, monástica observántia et miraculórum dono illústris.
\switchcolumn
\selectlanguage{english}
In the island of Aegina, St. 
 Athanasia, widow, celebrated for monastical observance and the gift of 
 miracles.
\switchcolumn*
\selectlanguage{latin}
\end{paracol}


% ---- martyrology/mart08/mart0815.htm
\needspace{10\baselineskip}
\begin{paracol}{2}
\selectlanguage{latin}
\begin{center}{\color{gregoriocolor} Décimo octávo Kaléndas Septémbris. 
 Luna\dots\ }\end{center}
\switchcolumn
\selectlanguage{english}
\begin{center}{\color{gregoriocolor} The 
 15\textsuperscript{th} Day of 
 August. The\dots\ Day of the Moon.}\end{center}
\end{paracol}

\noindent\begin{tabularx}{\linewidth}{*{19}{>{\centering\arraybackslash}X}}
 \textcolor{gregoriocolor}{a} & \textcolor{gregoriocolor}{b} & \textcolor{gregoriocolor}{c} & \textcolor{gregoriocolor}{d} & \textcolor{gregoriocolor}{e} & \textcolor{gregoriocolor}{f} & \textcolor{gregoriocolor}{g} & \textcolor{gregoriocolor}{h} & \textcolor{gregoriocolor}{i} & \textcolor{gregoriocolor}{k} & \textcolor{gregoriocolor}{l} & \textcolor{gregoriocolor}{m} & \textcolor{gregoriocolor}{n} & \textcolor{gregoriocolor}{p} & \textcolor{gregoriocolor}{q} & \textcolor{gregoriocolor}{r} & \textcolor{gregoriocolor}{s} & \textcolor{gregoriocolor}{t} & \textcolor{gregoriocolor}{u} \\
 21 & 22 & 23 & 24 & 25 & 26 & 27 & 28 & 29 & 1 & 2 & 3 & 4 & 5 & 6 & 7 & 8 & 9 & 10 \\
\end{tabularx}
\vspace{0.5\baselineskip}
\noindent\begin{tabularx}{\linewidth}{*{12}{>{\centering\arraybackslash}X}}
 \textcolor{gregoriocolor}{A} & \textcolor{gregoriocolor}{B} & \textcolor{gregoriocolor}{C} & \textcolor{gregoriocolor}{D} & \textcolor{gregoriocolor}{E} & F & \textcolor{gregoriocolor}{F} & \textcolor{gregoriocolor}{G} & \textcolor{gregoriocolor}{H} & \textcolor{gregoriocolor}{M} & \textcolor{gregoriocolor}{N} & \textcolor{gregoriocolor}{P} \\
 11 & 12 & 13 & 14 & 15 & 16 & 15 & 16 & 17 & 18 & 19 & 20 \\
\end{tabularx}

\begin{paracol}{2}
\selectlanguage{latin}
\lettrine[lines=2]{A}{ssúmptio} sanctíssimæ 
 Dei Genitrícis Vírginis Maríæ.
\switchcolumn
\selectlanguage{english}
\lettrine[lines=2]{T}{he} Assumption of the most holy 
 Virgin Mary, Mother of God.
\switchcolumn*
\selectlanguage{latin}
Cracóviæ, in Polónia, 
 natális sancti Hyacínthi, ex Ordine Prædicatórum, Confessóris, quem Clemens 
 Octávus, Póntifex Máximus, in Sanctórum númerum rétulit. Ipsíus autem 
 festum sextodécimo Kaléndas Septémbris celebrátur.
\switchcolumn
\selectlanguage{english}
At Cracow in Poland, St. Hyacinth, 
 confessor of the Order of Preachers, whom Pope Clement VIII placed in the 
 number of the saints. His feast is observed on the 17th of August.
\switchcolumn*
\selectlanguage{latin}
Apud Albam Regálem, in 
 Pannónia, item natális sancti Stéphani, Regis Hungarórum et Confessóris; 
 qui, divínis virtútibus exornátus, primus Húngaros ad Christi fidem 
 convértit, et a Deípara Vírgine, ipso die Assumptiónis suæ, in cælum 
 recéptus fuit. Ejus vero festívitas quarto Nonas Septémbris, quo die 
 munitíssima Budæ arx, sancti Regis ope, recólitur, ex dispositióne 
 Innocéntii Papæ Undécimi.
\switchcolumn
\selectlanguage{english}
At Alba Regalis in Hungary, St. 
 Stephen, King of Hungary, who was graced with divine virtues, was the first 
 to convert the Hungarians to the faith of Christ, and was received into 
 heaven by the Virgin Mother of God on the very day of her Assumption. 
 By decree of Pope Innocent XI, his feast is kept on the 2nd of September, on 
 which day the strong city of Buda, by the aid of the holy king, was 
 recovered by the Christian army.
\switchcolumn*
\selectlanguage{latin}
Romæ, via Appia, sancti 
 Tharsícii Acolythi, quem Pagáni, cum inveníssent Córporis Christi Sacraménta 
 portántem, cœpérunt disquírere quid géreret. At ipse, indígnum 
 júdicans porcis pródere margarítas, támdiu ab illis mactátus luto ejus 
 córpore, sacrílegi discussóres nihil Sacramentórum Christi in occísi mánibus 
 aut in véstibus invenérunt. Christiáni autem collegérunt Mártyris 
 corpus, et in cœmetério Callísti honorifice sepeliérunt.
\switchcolumn
\selectlanguage{english}
At Rome, on the Appian Way, St. 
 Tarsicius, acolyte. The pagans accosted him as he was carrying the 
 Sacrament of Christ's Body, and began to inquire what it was. But he 
 judged it an unworthy thing to cast pearls before swine. They 
 therefore beat him with sticks and stones until he expired. The 
 sacrilegious searchers examined his body, but found no vestige of the 
 Sacrament of Christ, either in his hands or in his clothes. The 
 Christians took up the body of the martyr, and buried it reverently in the 
 cemetery of Callistus.
\switchcolumn*
\selectlanguage{latin}
Tagáste, in Africa, 
 sancti Alípii Epíscopi, qui beáti Augustíni olim discípulus, póstea in 
 conversióne sócius, in múnere pastoráli colléga, et in certamínibus advérsus 
 hæréticos commílito strénuus, ac demum in cælésti glória consors fuit.
\switchcolumn
\selectlanguage{english}
At Tagaste in Africa, St. Alipius, 
 bishop, who was the disciple of blessed Augustine, and the companion of his 
 conversion, his colleague in the pastoral charge, his valiant fellow-soldier 
 in disputing heretics, and finally his partner in the glory of heaven.
\switchcolumn*
\selectlanguage{latin}
Suessióne, in Gálliis, 
 sancti Arnúlfi, Epíscopi et Confessóris.
\switchcolumn
\selectlanguage{english}
At Soissons in France, St. Arnold, 
 bishop and confessor.
\switchcolumn*
\selectlanguage{latin}
Romæ sancti Stanislái 
 Kostkæ Polóni, novítii e Societáte Jesu et Confessóris; qui, consummátus in 
 brevi, per angélicam vitæ innocéntiam explévit témpora multa, et a Benedícto 
 Décimo tértio, Pontífice Máximo, in Sanctórum album relátus est.
\switchcolumn
\selectlanguage{english}
At Rome, St. Stanislas Kostka, a 
 native of Poland, confessor of the Society of Jesus, who being made perfect 
 in a short time, fulfilled a long time by the angelic innocence of his life. 
 He was inscribed on the list of the saints by the Sovereign Pontiff, 
 Benedict XIII.
\switchcolumn*
\selectlanguage{latin}
\end{paracol}


% ---- martyrology/mart08/mart0816.htm
\needspace{10\baselineskip}
\begin{paracol}{2}
\selectlanguage{latin}
\begin{center}{\color{gregoriocolor} Décimo séptimo Kaléndas Septémbris. 
 Luna\dots\ }\end{center}
\switchcolumn
\selectlanguage{english}
\begin{center}{\color{gregoriocolor} The 
 16\textsuperscript{th} Day of 
 August. The\dots\ Day of the Moon.}\end{center}
\end{paracol}

\noindent\begin{tabularx}{\linewidth}{*{19}{>{\centering\arraybackslash}X}}
 \textcolor{gregoriocolor}{a} & \textcolor{gregoriocolor}{b} & \textcolor{gregoriocolor}{c} & \textcolor{gregoriocolor}{d} & \textcolor{gregoriocolor}{e} & \textcolor{gregoriocolor}{f} & \textcolor{gregoriocolor}{g} & \textcolor{gregoriocolor}{h} & \textcolor{gregoriocolor}{i} & \textcolor{gregoriocolor}{k} & \textcolor{gregoriocolor}{l} & \textcolor{gregoriocolor}{m} & \textcolor{gregoriocolor}{n} & \textcolor{gregoriocolor}{p} & \textcolor{gregoriocolor}{q} & \textcolor{gregoriocolor}{r} & \textcolor{gregoriocolor}{s} & \textcolor{gregoriocolor}{t} & \textcolor{gregoriocolor}{u} \\
 22 & 23 & 24 & 25 & 26 & 27 & 28 & 29 & 1 & 2 & 3 & 4 & 5 & 6 & 7 & 8 & 9 & 10 & 11 \\
\end{tabularx}
\vspace{0.5\baselineskip}
\noindent\begin{tabularx}{\linewidth}{*{12}{>{\centering\arraybackslash}X}}
 \textcolor{gregoriocolor}{A} & \textcolor{gregoriocolor}{B} & \textcolor{gregoriocolor}{C} & \textcolor{gregoriocolor}{D} & \textcolor{gregoriocolor}{E} & F & \textcolor{gregoriocolor}{F} & \textcolor{gregoriocolor}{G} & \textcolor{gregoriocolor}{H} & \textcolor{gregoriocolor}{M} & \textcolor{gregoriocolor}{N} & \textcolor{gregoriocolor}{P} \\
 12 & 13 & 14 & 15 & 16 & 17 & 16 & 17 & 18 & 19 & 20 & 21 \\
\end{tabularx}

\begin{paracol}{2}
\selectlanguage{latin}
\lettrine[lines=2]{S}{ancti} Jóachim, patris 
 immaculátæ Vírginis Genetrícis Dei Maríæ, Confessóris, cujus dies natális 
 refértur tertiodécimo Kaléndas Aprílis.
\switchcolumn
\selectlanguage{english}
\lettrine[lines=2]{S}{t.} Joachim, father of the most 
 Blessed Virgin Mary, Mother of God, Confessor. His birthday is noted 
 on the 20th of March.
\switchcolumn*
\selectlanguage{latin}
Romæ sancti Titi 
 Diáconi, qui, Urbe a Gothis occupáta, pecúnias paupéribus distríbuens, a 
 Tribúno bárbaro jussus est occídi.
\switchcolumn
\selectlanguage{english}
At Rome, St. Titus, deacon, who, 
 when the city was taken by the Goths, was put to death by a barbarous 
 tribune while distributing money to the poor.
\switchcolumn*
\selectlanguage{latin}
Nicææ, in Bithynia, 
 sancti Diomédis médici, qui in persecutióne Diocletiáni Imperatóris, pro 
 Christi fide cæsus gládio, martyrium complévit.
\switchcolumn
\selectlanguage{english}
At Nicaea in Bithynia, St. Diomede, 
 a physician who underwent martyrdom by being beheaded during the persecution 
 of Diocletian.
\switchcolumn*
\selectlanguage{latin}
In Palæstína sanctórum 
 trigínta trium Mártyrum.
\switchcolumn
\selectlanguage{english}
In Palestine thirty-three holy 
 martyrs.
\switchcolumn*
\selectlanguage{latin}
Ferentíni, in Hérnicis, 
 sancti Ambrósii Centuriónis, qui, in persecutióne Diocletiáni, váriis modis 
 cruciátus, et novíssime, cum per ignem illæsus transísset, demérsus in aquam, 
 edúctus est in refrigérium.
\switchcolumn
\selectlanguage{english}
At Ferentino in Campania, St. 
 Ambrose, centurion. In the persecution of Diocletian he was subjected 
 to different kinds of tortures, and finally passing through fire without 
 injury, was cast into the waters, and thus reached the place of eternal 
 rest.
\switchcolumn*
\selectlanguage{latin}
Medioláni deposítio 
 sancti Simpliciáni Epíscopi, sanctórum Ambrósii et Augustíni testimónio 
 célebris.
\switchcolumn
\selectlanguage{english}
At Milan, the death of St. 
 Simplician, bishop, renowned by the testimony of given of him by St. Ambrose 
 and St. Augustine.
\switchcolumn*
\selectlanguage{latin}
Antisiodóri sancti 
 Eleuthérii Epíscopi.
\switchcolumn
\selectlanguage{english}
At Auxerre, St. Eleutherius, bishop.
\switchcolumn*
\selectlanguage{latin}
Nicomedíæ sancti Arsácii 
 Confessóris, qui, sub Licínio persecutóre, milítia relícta, solitáriam vitam 
 ducens, tantis virtútibus cláruit, ut et dæmones expulísse et orándo 
 interemísse ingéntem dracónem legátur; dénique, futúram civitátis cladem 
 prænúntians, in oratióne spíritum Deo réddidit.
\switchcolumn
\selectlanguage{english}
At Nicomedia, St. Arsacius, 
 confessor. Under the persecution of Licinius he left the military 
 service, and leading a solitary life, became so famous for working miracles 
 that we read of his expelling the demons and killing a huge dragon by his 
 prayers. Finally he foretold the destruction of the city, and gave up 
 his soul to God in prayer.
\switchcolumn*
\selectlanguage{latin}
Apud Montem Pessulánum, 
 in Gállia Narbonénsi, deposítio beáti Rochi Confessóris, qui multas Itáliæ 
 urbes signo Crucis a morbo epidémiæ liberávit. Ipsíus corpus Venétias 
 póstea translátum, et in Ecclésia, ejus nómine consecráta, 
 honorificentíssime cónditum fuit.
\switchcolumn
\selectlanguage{english}
In France, near Montpellier, in the 
 province of Narbonne, the death of blessed Roch, confessor, who by the sing 
 of the cross, delivered many cities of Italy from an epidemic. His 
 body was afterwards transferred to Venice, and deposited with the greatest 
 honours in the church dedicated under his invocation.
\switchcolumn*
\selectlanguage{latin}
Romæ sanctæ Serénæ, 
 uxóris quondam Diocletiáni Augústi.
\switchcolumn
\selectlanguage{english}
At Rome, St. Serena, who had been 
 the wife of Emperor Diocletian.
\switchcolumn*
\selectlanguage{latin}
\end{paracol}


% ---- martyrology/mart08/mart0817.htm
\needspace{10\baselineskip}
\begin{paracol}{2}
\selectlanguage{latin}
\begin{center}{\color{gregoriocolor} Sextodécimo Kaléndas Septémbris. 
 Luna\dots\ }\end{center}
\switchcolumn
\selectlanguage{english}
\begin{center}{\color{gregoriocolor} The 
 17\textsuperscript{th} Day of 
 August. The\dots\ Day of the Moon.}\end{center}
\end{paracol}

\noindent\begin{tabularx}{\linewidth}{*{19}{>{\centering\arraybackslash}X}}
 \textcolor{gregoriocolor}{a} & \textcolor{gregoriocolor}{b} & \textcolor{gregoriocolor}{c} & \textcolor{gregoriocolor}{d} & \textcolor{gregoriocolor}{e} & \textcolor{gregoriocolor}{f} & \textcolor{gregoriocolor}{g} & \textcolor{gregoriocolor}{h} & \textcolor{gregoriocolor}{i} & \textcolor{gregoriocolor}{k} & \textcolor{gregoriocolor}{l} & \textcolor{gregoriocolor}{m} & \textcolor{gregoriocolor}{n} & \textcolor{gregoriocolor}{p} & \textcolor{gregoriocolor}{q} & \textcolor{gregoriocolor}{r} & \textcolor{gregoriocolor}{s} & \textcolor{gregoriocolor}{t} & \textcolor{gregoriocolor}{u} \\
 23 & 24 & 25 & 26 & 27 & 28 & 29 & 1 & 2 & 3 & 4 & 5 & 6 & 7 & 8 & 9 & 10 & 11 & 12 \\
\end{tabularx}
\vspace{0.5\baselineskip}
\noindent\begin{tabularx}{\linewidth}{*{12}{>{\centering\arraybackslash}X}}
 \textcolor{gregoriocolor}{A} & \textcolor{gregoriocolor}{B} & \textcolor{gregoriocolor}{C} & \textcolor{gregoriocolor}{D} & \textcolor{gregoriocolor}{E} & F & \textcolor{gregoriocolor}{F} & \textcolor{gregoriocolor}{G} & \textcolor{gregoriocolor}{H} & \textcolor{gregoriocolor}{M} & \textcolor{gregoriocolor}{N} & \textcolor{gregoriocolor}{P} \\
 13 & 14 & 15 & 16 & 17 & 18 & 17 & 18 & 19 & 20 & 21 & 22 \\
\end{tabularx}

\begin{paracol}{2}
\selectlanguage{latin}
\lettrine[lines=1]{O}{ctáva} sancti Lauréntii 
 Mártyris.
\switchcolumn
\selectlanguage{english}
\lettrine[lines=1]{T}{he} Octave of St. Lawrence, martyr.
\switchcolumn*
\selectlanguage{latin}
Sancti Hyacínthi, ex 
 Ordine Prædicatórum, Confessóris, qui décimo octávo Kaléndas Septémbris 
 obdormívit in Dómino.
\switchcolumn
\selectlanguage{english}
St. Hyacinth, confessor of the Order 
 of Preachers, who fell asleep in the Lord on the 15th of August.
\switchcolumn*
\selectlanguage{latin}
Carthágine sanctórum 
 Mártyrum Liberáti abbátis, Bonifátii Diáconi, Servi et Rústici Subdiaconórum, 
 Rogáti et Séptimi Monachórum, et Máximi púeri; qui, in persecutióne 
 Wandálica, sub Hunneríco Rege, pro confessióne cathólicæ fídei et pro únici 
 Baptísmatis defensióne, váriis et inaudítis supplíciis exagitáti, demum 
 super ligna, quibus concremándi erant, clavis confíxi, et, cum ignis sæpius 
 accénsus fuísset ac divínitus semper exstínctus, Regis jussu remórum 
 véctibus percússi, et, comminútis cérebris, enecáti, speciósum cursum 
 certáminis sui, coronánte Dómino, perfecérunt.
\switchcolumn
\selectlanguage{english}
At Carthage in Africa, the holy 
 martyrs Liberatus, abbot, Boniface, a deacon, Servus and Rusticus, 
 subdeacons, Rogatus and Septimus, monks, and Maximus, a young child. 
 In the persecution of the Vandals, under King Hunneric, they were subjected 
 to various and unheard-of torments for the confession of the Catholic faith 
 and the defence of one baptism. Finally, being nailed to the wood with 
 which they were to be burned, as the fire was always miraculously 
 extinguished whenever kindled, they were struck with iron bars by order of 
 the tyrant until their brains were dashed out. Thus they ended the 
 glorious series of their combats, and were crowned by our Lord.
\switchcolumn*
\selectlanguage{latin}
In Achája sancti 
 Myrónis, Presbyteri et Mártyris, qui, sub Décio Imperatóre et Antípatre 
 Præside, Cyzici, post multa torménta, cápite truncátus est.
\switchcolumn
\selectlanguage{english}
In Achaia, St. Myron, priest and 
 martyr, who was beheaded at Cyzicum after undergoing many torments, in the 
 time of Emperor Decius and the governor Antipater.
\switchcolumn*
\selectlanguage{latin}
Cæsaréæ, in Cappadócia, natális sancti Mamántis Mártyris, qui, sanctórum Theódoti et Rufínæ Mártyrum 
 fílius, longum a puerítia ad senectútem usque martyrium duxit, et tandem, 
 imperánte Aureliáno, sub Alexándro Præside, illud felíciter consummávit; 
 quem sancti Patres Basilíus et Gregórius Nazianzénus summis láudibus 
 celebrárunt.
\switchcolumn
\selectlanguage{english}
At Caesarea in Cappadocia, the 
 birthday of St. Mamas, martyr, the son of Saints Theodotus and Rufina, 
 martyrs, who, from childhood to old age, endured a long martyrdom, and at 
 length ended it happily in the reign of Aurelian, under the governor 
 Alexander. He has been highly praised by the holy Fathers Basil and 
 Gregory Nazianzen.
\switchcolumn*
\selectlanguage{latin}
Nicomedíæ sanctórum 
 Mártyrum Stratónis, Philíppi et Eutychiáni; qui, damnáti ad béstias et nil 
 læsi, per ignem martyrium consummárunt.
\switchcolumn
\selectlanguage{english}
At Nicomedia, the holy martyrs 
 Straton, Philip, and Eutychian, who were condemned to the beasts, but being 
 uninjured by them, ended their martyrdom by fire.
\switchcolumn*
\selectlanguage{latin}
Ptolemáide, in 
 Palæstína, pássio sanctórum Mártyrum Pauli, ejúsque soróris Juliánæ Vírginis; 
 qui ambo, sub Aureliáno Imperatóre, cum in Christi confessióne permanérent 
 immóbiles, jussi sunt váriis et diríssimis torméntis afflígi ac tandem 
 cápite obtruncári.
\switchcolumn
\selectlanguage{english}
At Ptolemais in Palestine, the holy 
 martyrs Paul and his sister Juliana, virgin, who suffered under Aurelian. 
 They were both punished with various cruel torments and were finally 
 beheaded for their constancy in confessing the name of Christ.
\switchcolumn*
\selectlanguage{latin}
Romæ sancti Eusébii 
 Papæ.
\switchcolumn
\selectlanguage{english}
At Rome, Pope St. Eusebius.
\switchcolumn*
\selectlanguage{latin}
Interámnæ sancti 
 Anastásii, Epíscopi et Confessóris.
\switchcolumn
\selectlanguage{english}
At Teramo, St. Anastasius, bishop 
 and confessor.
\switchcolumn*
\selectlanguage{latin}
In Monte Falco, in 
 Umbria, sanctæ Claræ, Moniális ex Ordine Eremitárum sancti Augustíni, 
 Vírginis; in cujus viscéribus renováta Domínicæ passiónis mystéria fidéles, 
 máxima cum devotióne, venerántur. Eam Leo Décimus tértius, Summus 
 Póntifex, sanctárum Vírginum albo adscrípsit.
\switchcolumn
\selectlanguage{english}
At Montefalco in Umbria, St. Clare, 
 a nun of the Order of Hermits of St. Augustine, virgin. In her flesh 
 were renewed the mysteries of the Lord's passion, which the faithful honour 
 with great devotion. Pope Leo XIII solemnly inscribed her in the list 
 of the holy virgins.
\switchcolumn*
\selectlanguage{latin}
\end{paracol}


% ---- martyrology/mart08/mart0818.htm
\needspace{10\baselineskip}
\begin{paracol}{2}
\selectlanguage{latin}
\begin{center}{\color{gregoriocolor} Quintodécimo Kaléndas Septémbris. 
 Luna\dots\ }\end{center}
\switchcolumn
\selectlanguage{english}
\begin{center}{\color{gregoriocolor} The 
 18\textsuperscript{th} Day of 
 August. The\dots\ Day of the Moon.}\end{center}
\end{paracol}

\noindent\begin{tabularx}{\linewidth}{*{19}{>{\centering\arraybackslash}X}}
 \textcolor{gregoriocolor}{a} & \textcolor{gregoriocolor}{b} & \textcolor{gregoriocolor}{c} & \textcolor{gregoriocolor}{d} & \textcolor{gregoriocolor}{e} & \textcolor{gregoriocolor}{f} & \textcolor{gregoriocolor}{g} & \textcolor{gregoriocolor}{h} & \textcolor{gregoriocolor}{i} & \textcolor{gregoriocolor}{k} & \textcolor{gregoriocolor}{l} & \textcolor{gregoriocolor}{m} & \textcolor{gregoriocolor}{n} & \textcolor{gregoriocolor}{p} & \textcolor{gregoriocolor}{q} & \textcolor{gregoriocolor}{r} & \textcolor{gregoriocolor}{s} & \textcolor{gregoriocolor}{t} & \textcolor{gregoriocolor}{u} \\
 24 & 25 & 26 & 27 & 28 & 29 & 1 & 2 & 3 & 4 & 5 & 6 & 7 & 8 & 9 & 10 & 11 & 12 & 13 \\
\end{tabularx}
\vspace{0.5\baselineskip}
\noindent\begin{tabularx}{\linewidth}{*{12}{>{\centering\arraybackslash}X}}
 \textcolor{gregoriocolor}{A} & \textcolor{gregoriocolor}{B} & \textcolor{gregoriocolor}{C} & \textcolor{gregoriocolor}{D} & \textcolor{gregoriocolor}{E} & F & \textcolor{gregoriocolor}{F} & \textcolor{gregoriocolor}{G} & \textcolor{gregoriocolor}{H} & \textcolor{gregoriocolor}{M} & \textcolor{gregoriocolor}{N} & \textcolor{gregoriocolor}{P} \\
 14 & 15 & 16 & 17 & 18 & 19 & 18 & 19 & 20 & 21 & 22 & 23 \\
\end{tabularx}

\begin{paracol}{2}
\selectlanguage{latin}
\lettrine[lines=2]{P}{rænéste} natális sancti 
 Agapíti Mártyris, qui, cum esset annórum quíndecim et amóre Christi fervéret, 
 jussu Aureliáni Imperatóris tentus est, ac primo nervis crudis diutíssime 
 cæsus, deínde, sub Antíocho Præfécto, gravióra supplícia passus; exínde, cum 
 ex Imperatóris præcépto leónibus objicerétur et mínime læsus esset, gládio 
 ministrórum coronándus percútitur.
\switchcolumn
\selectlanguage{english}
\lettrine[lines=2]{A}{t} Palestrina, the birthday of the 
 holy martyr Agapitus. Although only fifteen years of age, because he 
 was fervent in the love of Christ, he was arrested by order of Emperor 
 Aurelian, and scourged for a long time. Afterwards, under the prefect 
 Antiochus, he endured more severe torments, and being delivered to the lions 
 by the emperor's order without receiving any injury, he was finally struck 
 with the sword, and thus merited his crown.
\switchcolumn*
\selectlanguage{latin}
Romæ beatórum Joánnis 
 et Crispi Presbyterórum, qui, in persecutióne Diocletiáni, multa Sanctórum 
 córpora officiosíssime sepeliérunt, quorum méritis et ipsi póstmodum sociáti, 
 gáudia vitæ ætérnæ sibi comparárunt.
\switchcolumn
\selectlanguage{english}
At Rome, during the persecution of 
 Diocletian, the blessed John and Crispus, priests, who charitably buried the 
 bodies of many saints; afterwards becoming partakers of their merits, they 
 deserved the joys of eternal life.
\switchcolumn*
\selectlanguage{latin}
Item Romæ sanctórum 
 Mártyrum Hermæ, Serapiónis et Polyæni, qui, per angústa, saxósa et áspera 
 loca raptáti, ánimas Deo reddidérunt.
\switchcolumn
\selectlanguage{english}
In the same city, the holy martyrs 
 Hermas, Serapion, and Polyaenus. Being dragged through narrow, stony, 
 and rough places, they yielded up their souls to God.
\switchcolumn*
\selectlanguage{latin}
In Illyrico sanctórum 
 Mártyrum Flori et Lauri, artis lapicidínæ, qui, sub Licióne Præside, 
 martyrio consúmptis eórum magístris Próculo et Máximo, ambo, post multa 
 torménta, in profúndum púteum sunt demérsi.
\switchcolumn
\selectlanguage{english}
In Illyria, the holy martyrs Florus 
 and Laurus, stonecutters, who, after the martyrdom of Proculus and Maximus, 
 their employers, were subjected to many torments under the governor Licion, 
 and plunged into a deep well.
\switchcolumn*
\selectlanguage{latin}
Myræ, in Lycia, 
 sanctórum Mártyrum Leónis et Juliánæ.
\switchcolumn
\selectlanguage{english}
At Myra in Lycia, the holy martyrs 
 Leo and Juliana.
\switchcolumn*
\selectlanguage{latin}
Metis, in Gállia, 
 sancti Firmíni, Epíscopi et Confessóris.
\switchcolumn
\selectlanguage{english}
At Metz in France, St. Firmin, 
 bishop and confessor.
\switchcolumn*
\selectlanguage{latin}
Romæ, via Lavicána, 
 sanctæ Hélenæ, matris Constantíni Magni, piíssimi Imperatóris, qui primus 
 egrégium Ecclésiæ tuéndæ atque amplificándæ exémplum céteris Princípibus 
 præbuit.
\switchcolumn
\selectlanguage{english}
At Rome, on the Via Lavicana, St. 
 Helena, mother of the religious emperor Constantine the Great, who was the 
 first to set the example to other princes of protecting and extending the 
 Church.
\switchcolumn*
\selectlanguage{latin}
\end{paracol}


% ---- martyrology/mart08/mart0819.htm
\needspace{10\baselineskip}
\begin{paracol}{2}
\selectlanguage{latin}
\begin{center}{\color{gregoriocolor} Quartodécimo Kaléndas Septémbris. 
 Luna\dots\ }\end{center}
\switchcolumn
\selectlanguage{english}
\begin{center}{\color{gregoriocolor} The 
 19\textsuperscript{th} Day of 
 August. The\dots\ Day of the Moon.}\end{center}
\end{paracol}

\noindent\begin{tabularx}{\linewidth}{*{19}{>{\centering\arraybackslash}X}}
 \textcolor{gregoriocolor}{a} & \textcolor{gregoriocolor}{b} & \textcolor{gregoriocolor}{c} & \textcolor{gregoriocolor}{d} & \textcolor{gregoriocolor}{e} & \textcolor{gregoriocolor}{f} & \textcolor{gregoriocolor}{g} & \textcolor{gregoriocolor}{h} & \textcolor{gregoriocolor}{i} & \textcolor{gregoriocolor}{k} & \textcolor{gregoriocolor}{l} & \textcolor{gregoriocolor}{m} & \textcolor{gregoriocolor}{n} & \textcolor{gregoriocolor}{p} & \textcolor{gregoriocolor}{q} & \textcolor{gregoriocolor}{r} & \textcolor{gregoriocolor}{s} & \textcolor{gregoriocolor}{t} & \textcolor{gregoriocolor}{u} \\
 25 & 26 & 27 & 28 & 29 & 1 & 2 & 3 & 4 & 5 & 6 & 7 & 8 & 9 & 10 & 11 & 12 & 13 & 14 \\
\end{tabularx}
\vspace{0.5\baselineskip}
\noindent\begin{tabularx}{\linewidth}{*{12}{>{\centering\arraybackslash}X}}
 \textcolor{gregoriocolor}{A} & \textcolor{gregoriocolor}{B} & \textcolor{gregoriocolor}{C} & \textcolor{gregoriocolor}{D} & \textcolor{gregoriocolor}{E} & F & \textcolor{gregoriocolor}{F} & \textcolor{gregoriocolor}{G} & \textcolor{gregoriocolor}{H} & \textcolor{gregoriocolor}{M} & \textcolor{gregoriocolor}{N} & \textcolor{gregoriocolor}{P} \\
 15 & 16 & 17 & 18 & 19 & 20 & 19 & 20 & 21 & 22 & 23 & 24 \\
\end{tabularx}

\begin{paracol}{2}
\selectlanguage{latin}
\lettrine[lines=2]{C}{adómi,} in Gállia, 
 sancti Joánnis Eudes Confessóris, Missionárii Apostólici, Fundatóris 
 Congregatiónis Presbyterórum Jesu et Maríæ necnon Ordinis Moniálium Dóminæ 
 Nostræ a Caritáte, et promotóris litúrgici cultus erga Sacratíssima Christi 
 ejúsque Genitrícis Corda; quem Pius Papa Undécimus fastis Sanctórum 
 adscrípsit.
\switchcolumn
\selectlanguage{english}
\lettrine[lines=2]{A}{t} Caen in France, St. John Eudes, 
 apostolic missionary, founder of the Congregation of Priests of Jesus and 
 Mary and of the Order of Nuns of our Lady of Charity, and the promoter of 
 the liturgical cult towards the most sacred Hearts of Christ and his Mother. 
 He was canonized by Pope Pius XI.
\switchcolumn*
\selectlanguage{latin}
Romæ sancti Júlii, 
 Senatóris et Mártyris; qui, Vitéllio Júdici tráditus et ab eo in cárcerem 
 trusus, támdiu, jubénte Cómmodo Imperatóre, fústibus cæsus est, donec 
 emítteret spíritum. Ipsíus autem corpus in cœmetério Calepódii, via 
 Aurélia, sepúltum fuit.
\switchcolumn
\selectlanguage{english}
At Rome, St. Julius, senator and 
 martyr, who was delivered up to the judge Vitellius, and thrown into prison 
 by him. By order of Emperor Commodus, he was beaten with rods until he 
 expired. His body was buried in the cemetery of Caleposius on the 
 Aurelian Way.
\switchcolumn*
\selectlanguage{latin}
Anágniæ sancti Magni, 
 Epíscopi et Mártyris, qui in persecutióne Décii necátus est.
\switchcolumn
\selectlanguage{english}
At Anagni, St. Magnus, bishop and 
 martyr, who was put to death in the persecution of Decius.
\switchcolumn*
\selectlanguage{latin}
In Cilícia natális 
 sancti Andréæ Tribúni, et Sociórum mílitum; qui, victória de Persis 
 divínitus obténta, ad Christi fidem sunt convérsi, et, hoc nómine accusáti, 
 sub Maximiáno Imperatóre, in angústiis Tauri montis, a Seléuci Præsidis 
 exércitu trucidáti sunt.
\switchcolumn
\selectlanguage{english}
In Cilicia, the birthday of St. 
 Andrew, tribune, and his military companions, who were converted to 
 Christianity through a miraculous victory they had gained over the Persians. 
 Being accused on this account, they were massacred in the Mount Taurus pass, 
 by the army of the governor Seleucus, under Emperor Maximian.
\switchcolumn*
\selectlanguage{latin}
In Palæstína sancti 
 Timóthei Mártyris, qui, in persecutióne Diocletiáni, sub Urbáno Præside, 
 post multa superáta supplícia, lento igne combústus est. Passi sunt 
 étiam ibídem Thecla et Agápius, e quibus Thecla, feris expósita, eárum 
 laniáta déntibus transívit ad Sponsum; Agápius vero, plúrima torménta 
 perpéssus, ad majóra certámina fuit dilátus.
\switchcolumn
\selectlanguage{english}
In Palestine, St. Timothy, a martyr 
 in the persecution of Diocletian, under the governor Urbanus. After 
 overcoming many torments, he was consumed with a slow fire. In the 
 same country there suffered also Thecla and Agapius. The former, being 
 exposed to the beasts, was torn to pieces by their teeth, and went to her 
 Spouse; but Agapius, after enduring many torments, was reserved for greater 
 trials.
\switchcolumn*
\selectlanguage{latin}
Romæ sancti Xysti 
 Tértii, Papæ et Confessóris.
\switchcolumn
\selectlanguage{english}
At Rome, St. Sixtus III, pope and 
 confessor.
\switchcolumn*
\selectlanguage{latin}
Apud castrum Bríncolam, 
 in Província, deposítio sancti Ludovíci, ex Ordine Minórum, Epíscopi 
 Tolosiáni, vitæ sanctitáte et miráculis clari; cujus corpus, inde Massíliam 
 translátum, in Ecclésia Fratrum Minórum honorífice cónditum fuit, ac póstea 
 Valéntiam, in Hispánia, devéctum est, atque in cathedráli Ecclésia 
 collocátum.
\switchcolumn
\selectlanguage{english}
In Provence, at the village of 
 Brignoles, the death of St. Louis, bishop of Toulouse, of the Order of 
 Friars Minor, renowned for holiness of life and miracles. His body was 
 taken to Marseilles, and buried with due honours in the Church of the Friars 
 Minor, but afterwards it was taken to Valencia in Spain, and enshrined in 
 the cathedral.
\switchcolumn*
\selectlanguage{latin}
In pago Sigistérico, in 
 Gállia, beáti Donáti, Presbyteri et Confessóris; qui, ab ipsis usque 
 infántiæ rudiméntis mira Dei grátia præditus, anachoréticam vitam multis 
 annis exégit, et miraculórum glória clarus migrávit ad Christum.
\switchcolumn
\selectlanguage{english}
In the neighbourhood of Sisteron in 
 France, blessed Donatus, priest and confessor. Being from his very 
 infancy endowed with the grace of God in an extraordinary manner, he lived 
 the life of an anchoret for many years, and after having been renowned for 
 glorious miracles, went to Christ.
\switchcolumn*
\selectlanguage{latin}
In território 
 Bituricénsi sancti Mariáni Confessóris, cujus virtútes et mirácula beátus 
 Gregórius, Turonénsis Epíscopus, magnis láudibus celebrávit.
\switchcolumn
\selectlanguage{english}
In the territory of Bourges, St. 
 Marianus, confessor, whose virtues and miracles were described with great 
 praise by St. Gregory, bishop of Tours.
\switchcolumn*
\selectlanguage{latin}
Mántuæ sancti Rufíni 
 Confessóris.
\switchcolumn
\selectlanguage{english}
At Mantua, St. Rufinus, confessor.
\switchcolumn*
\selectlanguage{latin}
Norimbérgæ sancti 
 Sebáldi Eremítæ, virtútibus et miráculis præclári, qui Sanctórum catálogo a 
 Martíno Papa Quinto adjéctus est.
\switchcolumn
\selectlanguage{english}
At Nuremburg, St. Sebald, hermit, 
 noted for his virtues and miracles. Pope Martin V added his name to 
 the list of the saints.
\switchcolumn*
\selectlanguage{latin}
\end{paracol}


% ---- martyrology/mart08/mart0820.htm
\needspace{10\baselineskip}
\begin{paracol}{2}
\selectlanguage{latin}
\begin{center}{\color{gregoriocolor} Tertiodécimo Kaléndas Septémbris. 
 Luna\dots\ }\end{center}
\switchcolumn
\selectlanguage{english}
\begin{center}{\color{gregoriocolor} The 
 20\textsuperscript{th} Day of 
 August. The\dots\ Day of the Moon.}\end{center}
\end{paracol}

\noindent\begin{tabularx}{\linewidth}{*{19}{>{\centering\arraybackslash}X}}
 \textcolor{gregoriocolor}{a} & \textcolor{gregoriocolor}{b} & \textcolor{gregoriocolor}{c} & \textcolor{gregoriocolor}{d} & \textcolor{gregoriocolor}{e} & \textcolor{gregoriocolor}{f} & \textcolor{gregoriocolor}{g} & \textcolor{gregoriocolor}{h} & \textcolor{gregoriocolor}{i} & \textcolor{gregoriocolor}{k} & \textcolor{gregoriocolor}{l} & \textcolor{gregoriocolor}{m} & \textcolor{gregoriocolor}{n} & \textcolor{gregoriocolor}{p} & \textcolor{gregoriocolor}{q} & \textcolor{gregoriocolor}{r} & \textcolor{gregoriocolor}{s} & \textcolor{gregoriocolor}{t} & \textcolor{gregoriocolor}{u} \\
 26 & 27 & 28 & 29 & 1 & 2 & 3 & 4 & 5 & 6 & 7 & 8 & 9 & 10 & 11 & 12 & 13 & 14 & 15 \\
\end{tabularx}
\vspace{0.5\baselineskip}
\noindent\begin{tabularx}{\linewidth}{*{12}{>{\centering\arraybackslash}X}}
 \textcolor{gregoriocolor}{A} & \textcolor{gregoriocolor}{B} & \textcolor{gregoriocolor}{C} & \textcolor{gregoriocolor}{D} & \textcolor{gregoriocolor}{E} & F & \textcolor{gregoriocolor}{F} & \textcolor{gregoriocolor}{G} & \textcolor{gregoriocolor}{H} & \textcolor{gregoriocolor}{M} & \textcolor{gregoriocolor}{N} & \textcolor{gregoriocolor}{P} \\
 16 & 17 & 18 & 19 & 20 & 21 & 20 & 21 & 22 & 23 & 24 & 25 \\
\end{tabularx}

\begin{paracol}{2}
\selectlanguage{latin}
\lettrine[lines=2]{I}{n} território 
 Lingoniénsi deposítio sancti Bernárdi, primi Clarævallénsis Abbátis, vita, 
 doctrína et miráculis gloriósi, quem Pius Octávus, Póntifex Máximus, 
 universális Ecclésiæ Doctórem declarávit et confirmávit.
\switchcolumn
\selectlanguage{english}
\lettrine[lines=2]{I}{n} the territory of Langres, the 
 death of St. Bernard, first abbot of Clairvaux, illustrious for virtues, 
 learning, and miracles. He was declared and confirmed doctor of the 
 Universal Church by the Sovereign Pontiff, Pius VIII.
\switchcolumn*
\selectlanguage{latin}
Romæ deposítio sancti 
 Pii Décimi, Papæ et Confessóris, fídei integritátis et ecclesiásticæ 
 libertátis propugnatóris invícti, religionísque zelo insígnis, cujus festum 
 tértio Nonas Septémbris recólitur.
\switchcolumn
\selectlanguage{english}
At Rome, the death of St. Pius X, 
 pope and confessor, who championed the integrity of the faith and the 
 liberty of the Church, and was renowned for his religious zeal. His 
 feastday is celebrated on the 3rd of September.
\switchcolumn*
\selectlanguage{latin}
Apud montem Senárium, 
 in Etrúria, natális sancti Manétti Confessóris, e septem Fundatóribus 
 Ordinis Servórum beátæ Maríæ Vírginis; qui, eídem hymnos dicens, exspirávit. 
 Ipsíus autem ac Sociórum festum prídie Idus Februárii celebrátur.
\switchcolumn
\selectlanguage{english}
On Mount Senario in Tuscany, the 
 birthday of St. Manetto, confessor, one of the seven founders of the Order 
 of the Servites of the Blessed Virgin Mary, who died as he was repeating a 
 hymn to her. His feast, with that of his companions, is kept on the 
 12th of February.
\switchcolumn*
\selectlanguage{latin}
In Judæa sancti 
 Samuélis Prophétæ, cujus sacra ossa (ut beátus Hierónymus scribit) Arcádius 
 Augústus Constantinópolim tránstulit, et prope Séptimum collocávit.
\switchcolumn
\selectlanguage{english}
In Judea, the holy prophet Samuel, 
 whose holy relics (as is related by St. Jerome) were taken to Constantinople 
 by Emperor Arcadius, and deposited near Septimum.
\switchcolumn*
\selectlanguage{latin}
In Cypro sancti Lúcii 
 Senatóris, qui, perspécta constántia Theodóri, Cyrenénsis Epíscopi, in 
 martyrio pósiti, ad Christi fidem est convérsus, et ad eam étiam Digniánum 
 Præsidem pertráxit; cum eóque Cyprum proféctus, ibi, cum álios Christiános 
 pro confessióne Dómini coronári vidéret, ultro se ipsum óbtulit, et cápitis 
 obtruncatióne eándem martyrii corónam proméruit.
\switchcolumn
\selectlanguage{english}
In Cyprus, St. Lucius, senator, who 
 was converted to the faith on seeing the constancy of Theodore, bishop of 
 Cyrene, during his martyrdom. He also converted the governor Dignian, 
 with whom he set out for Cyprus, where, seeing other Christians crowned for 
 the confession of the Lord, he offered himself voluntarily, and merited the 
 same crown of martyrdom by beheading.
\switchcolumn*
\selectlanguage{latin}
In Thrácia sanctórum 
 trigínta septem Mártyrum, qui, sub Præside Apelliáno, pro Christi fide, 
 mánibus pedibúsque præcísis, in camínum ardéntem injécti sunt.
\switchcolumn
\selectlanguage{english}
In Thrace, in the time of the 
 governor Apellian, thirty-seven holy martyrs, who had their hands and feet 
 cut off for the faith of Christ, and were cast into a burning furnace.
\switchcolumn*
\selectlanguage{latin}
Ibídem sanctórum 
 Mártyrum Sevéri, et Memnónis Centuriónis; qui, eódem mortis génere 
 consummáti, victóres abiérunt in cælum.
\switchcolumn
\selectlanguage{english}
Also, the holy martyrs Severus, and 
 the centurion Memnon, who, suffering the same kind of death, went 
 victoriously to heaven.
\switchcolumn*
\selectlanguage{latin}
Córdubæ, in Hispánia, 
 sanctórum Mártyrum Leovigíldi et Christóphori Monachórum, qui, in Arabum 
 persecutióne, pro Christiánæ fídei defensióne in cárcerem conjécti, ac mox, 
 cervícibus abscíssis, igni tráditi, martyrii palmam adépti sunt.
\switchcolumn
\selectlanguage{english}
At Cordova, during the persecution 
 of the Arabs, the holy martyrs Leovigild and Christopher, monks, who were 
 thrust into prison for the defence of the Christian faith, and soon after, 
 being beheaded and cast into the fire, thus obtained the palm of martyrdom.
\switchcolumn*
\selectlanguage{latin}
In Hério ínsula sancti 
 Philibérti Abbátis.
\switchcolumn
\selectlanguage{english}
In the island of Hermoutier, St. 
 Philibert, abbot.
\switchcolumn*
\selectlanguage{latin}
Romæ beáti Porphyrii, 
 qui fuit homo Dei, et sanctum Mártyrem Agapítum erudívit in fide et doctrína 
 Christi.
\switchcolumn
\selectlanguage{english}
At Rome, blessed Porphyry, a man of 
 God, who instructed the holy martyr Agapitus in the faith and doctrine of 
 Christ.
\switchcolumn*
\selectlanguage{latin}
In castro Cainóne, in 
 Gállia, sancti Máximi Confessóris, qui éxstitit discípulus beáti Mártini 
 Epíscopi.
\switchcolumn
\selectlanguage{english}
At Chinon, St. Maximus, confessor, 
 disciple of the blessed bishop Martin.
\switchcolumn*
\selectlanguage{latin}
\end{paracol}


% ---- martyrology/mart08/mart0821.htm
\needspace{10\baselineskip}
\begin{paracol}{2}
\selectlanguage{latin}
\begin{center}{\color{gregoriocolor} Duodécimo Kaléndas Septémbris. 
 Luna\dots\ }\end{center}
\switchcolumn
\selectlanguage{english}
\begin{center}{\color{gregoriocolor} The 
 21\textsuperscript{st} Day of 
 August. The\dots\ Day of the Moon.}\end{center}
\end{paracol}

\noindent\begin{tabularx}{\linewidth}{*{19}{>{\centering\arraybackslash}X}}
 \textcolor{gregoriocolor}{a} & \textcolor{gregoriocolor}{b} & \textcolor{gregoriocolor}{c} & \textcolor{gregoriocolor}{d} & \textcolor{gregoriocolor}{e} & \textcolor{gregoriocolor}{f} & \textcolor{gregoriocolor}{g} & \textcolor{gregoriocolor}{h} & \textcolor{gregoriocolor}{i} & \textcolor{gregoriocolor}{k} & \textcolor{gregoriocolor}{l} & \textcolor{gregoriocolor}{m} & \textcolor{gregoriocolor}{n} & \textcolor{gregoriocolor}{p} & \textcolor{gregoriocolor}{q} & \textcolor{gregoriocolor}{r} & \textcolor{gregoriocolor}{s} & \textcolor{gregoriocolor}{t} & \textcolor{gregoriocolor}{u} \\
 27 & 28 & 29 & 1 & 2 & 3 & 4 & 5 & 6 & 7 & 8 & 9 & 10 & 11 & 12 & 13 & 14 & 15 & 16 \\
\end{tabularx}
\vspace{0.5\baselineskip}
\noindent\begin{tabularx}{\linewidth}{*{12}{>{\centering\arraybackslash}X}}
 \textcolor{gregoriocolor}{A} & \textcolor{gregoriocolor}{B} & \textcolor{gregoriocolor}{C} & \textcolor{gregoriocolor}{D} & \textcolor{gregoriocolor}{E} & F & \textcolor{gregoriocolor}{F} & \textcolor{gregoriocolor}{G} & \textcolor{gregoriocolor}{H} & \textcolor{gregoriocolor}{M} & \textcolor{gregoriocolor}{N} & \textcolor{gregoriocolor}{P} \\
 17 & 18 & 19 & 20 & 21 & 22 & 21 & 22 & 23 & 24 & 25 & 26 \\
\end{tabularx}

\begin{paracol}{2}
\selectlanguage{latin}
\lettrine[lines=2]{S}{anctæ} Joánnæ-Francíscæ 
 Frémiot de Chantal, Víduæ, quæ Ordinis Sanctimoniálium Visitatiónis sanctæ 
 Maríæ fuit Institútrix, cujus dies natális recólitur Idibus Decémbris.
\switchcolumn
\selectlanguage{english}
\lettrine[lines=2]{T}{he} festival of St. Jane Frances 
 Fremiot de Chantal, foundress of the Order of Nuns of the Visitation of St. 
 Mary, whose birthday is commemorated on the 13th of December.
\switchcolumn*
\selectlanguage{latin}
Romæ, in agro Veráno, 
 sanctæ Cyríacæ, Víduæ et Mártyris; quæ, in persecutióne Valeriáni, cum se 
 súaque ómnia in Sanctórum ministéria impendísset, demum, martyrium pro 
 Christo súbiens, vitam quoque ipsam libénter impéndit.
\switchcolumn
\selectlanguage{english}
At Rome, in the Agro Verano, St. 
 Cyriaca, widow and martyr. In the persecution of Valerian, after 
 devoting herself and all her goods in the service of the saints, she gave up 
 her life by suffering martyrdom for Christ.
\switchcolumn*
\selectlanguage{latin}
In território 
 Gavalitáno sancti Priváti, Epíscopi et Mártyris; qui passus est in 
 persecutióne Valeriáni et Galliéni.
\switchcolumn
\selectlanguage{english}
In Gevaudan, St. Privatus, bishop 
 and martyr, who suffered in the persecution of Valerian and Gallienus.
\switchcolumn*
\selectlanguage{latin}
Salónæ, in Dalmátia, sancti Anastásii Corniculárii, qui, cum vidéret beátum Agapítum constánter 
 torménta perferéntem, convérsus est ad fidem, et, pro confessióne nóminis 
 Christi, jubénte Aureliáno Imperatóre, interémptus, Martyr migrávit ad 
 Dóminum.
\switchcolumn
\selectlanguage{english}
At Salona in Dalmatia, St. 
 Anastasius, a law officer, who was converted to the faith by seeing the 
 fortitude with which blessed Agapitus bore his torments, and being put to 
 death by order of Emperor Aurelian for confessing the name of Christ, went 
 to our Lord, a martyr.
\switchcolumn*
\selectlanguage{latin}
In Sardínia natális 
 sanctórum Mártyrum Luxórii, Cisélli et Cameríni; qui, in persecutióne 
 Diocletiáni, sub Délphio Præside, gládio cæsi sunt.
\switchcolumn
\selectlanguage{english}
In Sardinia, the birthday of the 
 holy martyrs Luxorius, Cisellus, and Camerinus, who were put to the sword in 
 the persecution of Diocletian, under the governor Delphius.
\switchcolumn*
\selectlanguage{latin}
Eódem die sanctórum 
 Mártyrum Bonósi et Maximiáni.
\switchcolumn
\selectlanguage{english}
On the same day, the holy martyrs 
 Bononus and Maximian.
\switchcolumn*
\selectlanguage{latin}
Fundis, in Látio, 
 sancti Patérni Mártyris, qui, ab Alexandría Romam venit ad Apostolórum 
 memórias; et inde in agrum Fundánum secéssit, atque ibi, cum Mártyrum 
 córpora sepelíret, a Tribúno comprehénsus est, et in vínculis exspirávit.
\switchcolumn
\selectlanguage{english}
At Fundi in Campania, St. Paternus, 
 a martyr, who came from Alexandria to Rome to visit the tomb of the 
 apostles. Thence he retired to the neighbourhood of Fundi, where, 
 being seized by the tribune while he was burying the bodies of the martyrs, 
 he died in captivity.
\switchcolumn*
\selectlanguage{latin}
Edéssæ, in Syria, 
 sanctórum Mártyrum Bassæ, ac trium ejus filiórum, id est Theogónii, Agápii 
 et Fidélis; quos, in persecutióne Maximiáni, pia mater exhórtans, martyrio 
 coronátos præmísit ad palmam, et, truncáto cápite, gaudens secúta est cum 
 victória.
\switchcolumn
\selectlanguage{english}
At Edessa in Syria, during the 
 persecution of Maximian, the holy martyrs Bassa, and her sons Theogonius, 
 Agapius, and Fidelis, whom their pious mother exhorted to martyrdom and sent 
 before her bearing their crowns. Being herself beheaded, she joyfully 
 followed them and shared their victory.
\switchcolumn*
\selectlanguage{latin}
Verónæ sancti Euprépii, 
 Epíscopi et Confessóris.
\switchcolumn
\selectlanguage{english}
At Verona, St. Euprepius, bishop and 
 confessor.
\switchcolumn*
\selectlanguage{latin}
Item sancti Quadráti 
 Epíscopi.
\switchcolumn
\selectlanguage{english}
Also, St. Quadratus, bishop.
\switchcolumn*
\selectlanguage{latin}
Arvérnis, in Gállia, 
 sancti Sidónii Epíscopi, doctrína et sanctitáte conspícui.
\switchcolumn
\selectlanguage{english}
In Auvergne in France, St. Sidonius, 
 bishop, noted for learning and holiness.
\switchcolumn*
\selectlanguage{latin}
Senis, in Túscia, beáti 
 Bernárdi Ptolomæi Abbátis, Congregatiónis Olivetánæ Fundatóris.
\switchcolumn
\selectlanguage{english}
At Siena in Tuscany, blessed Bernard 
 Ptolemy, abbot and founder of the Congregation of Olivetans.
\switchcolumn*
\selectlanguage{latin}
\end{paracol}


% ---- martyrology/mart08/mart0822.htm
\needspace{10\baselineskip}
\begin{paracol}{2}
\selectlanguage{latin}
\begin{center}{\color{gregoriocolor} Undécimo Kaléndas Septémbris. 
 Luna\dots\ }\end{center}
\switchcolumn
\selectlanguage{english}
\begin{center}{\color{gregoriocolor} The 
 22\textsuperscript{nd} Day of 
 August. The\dots\ Day of the Moon.}\end{center}
\end{paracol}

\noindent\begin{tabularx}{\linewidth}{*{19}{>{\centering\arraybackslash}X}}
 \textcolor{gregoriocolor}{a} & \textcolor{gregoriocolor}{b} & \textcolor{gregoriocolor}{c} & \textcolor{gregoriocolor}{d} & \textcolor{gregoriocolor}{e} & \textcolor{gregoriocolor}{f} & \textcolor{gregoriocolor}{g} & \textcolor{gregoriocolor}{h} & \textcolor{gregoriocolor}{i} & \textcolor{gregoriocolor}{k} & \textcolor{gregoriocolor}{l} & \textcolor{gregoriocolor}{m} & \textcolor{gregoriocolor}{n} & \textcolor{gregoriocolor}{p} & \textcolor{gregoriocolor}{q} & \textcolor{gregoriocolor}{r} & \textcolor{gregoriocolor}{s} & \textcolor{gregoriocolor}{t} & \textcolor{gregoriocolor}{u} \\
 28 & 29 & 1 & 2 & 3 & 4 & 5 & 6 & 7 & 8 & 9 & 10 & 11 & 12 & 13 & 14 & 15 & 16 & 17 \\
\end{tabularx}
\vspace{0.5\baselineskip}
\noindent\begin{tabularx}{\linewidth}{*{12}{>{\centering\arraybackslash}X}}
 \textcolor{gregoriocolor}{A} & \textcolor{gregoriocolor}{B} & \textcolor{gregoriocolor}{C} & \textcolor{gregoriocolor}{D} & \textcolor{gregoriocolor}{E} & F & \textcolor{gregoriocolor}{F} & \textcolor{gregoriocolor}{G} & \textcolor{gregoriocolor}{H} & \textcolor{gregoriocolor}{M} & \textcolor{gregoriocolor}{N} & \textcolor{gregoriocolor}{P} \\
 18 & 19 & 20 & 21 & 22 & 23 & 22 & 23 & 24 & 25 & 26 & 27 \\
\end{tabularx}

\begin{paracol}{2}
\selectlanguage{latin}
\lettrine[lines=2]{O}{ctava} Assumptiónis 
 beátæ Maríæ Vírginis.
\switchcolumn
\selectlanguage{english}
\lettrine[lines=2]{T}{he} Octave of the Assumption of the 
 Blessed Virgin Mary.
\switchcolumn*
\selectlanguage{latin}
Festum Immaculáti 
 Cordis ejúsdem beátæ Vírginis Maríæ.
\switchcolumn
\selectlanguage{english}
Feast of the Immaculate Heart of the 
 same Blessed Virgin Mary.
\switchcolumn*
\selectlanguage{latin}
Romæ, via Ostiénsi, 
 natális sancti Timóthei Mártyris, qui, a Præfécto Urbis Tarquínio tentus, et 
 longa cárceris custódia macerátus, et, cum sacrificáre idólis noluísset, 
 tértio cæsus et gravíssimis supplíciis attrectátus, ad últimum decollátus 
 est.
\switchcolumn
\selectlanguage{english}
At Rome, on the Ostian Way, the 
 birthday of the holy martyr Timothy. After he had been arrested by 
 Tarquin, prefect of the city, and kept for a long time in prison, because he 
 refused to sacrifice to idols, he was scourged three times, subjected to the 
 most severe torments, and finally beheaded.
\switchcolumn*
\selectlanguage{latin}
In Portu Románo sancti 
 Hippóloyti Epíscopi, eruditióne claríssimi, qui, sub Alexándro Imperatóre, ob 
 præcláram fídei confessiónem, mánibus pedibúsque ligátis in altam fóveam, 
 aquis plenam, præcipitátus, martyrii palmam accépit; cujus corpus apud 
 eúndem locum a Christiánis sepúltum fuit.
\switchcolumn
\selectlanguage{english}
At Porto, St. Hippolytus, bishop, 
 most renowned for learning. Having gloriously confessed the faith, in 
 the time of Emperor Alexander, he was bound hand and foot, thrown into a 
 deep ditch filled with water, and thus received the palm of martyrdom. 
 His body was buried by the Christians at that place.
\switchcolumn*
\selectlanguage{latin}
Augustodúni sancti 
 Symphoriáni Mártyris, qui, témpore Aureliáni Imperatóris, cum sacrificáre 
 nollet idólis, primo verbéribus afflíctus est, deínde cárceri mancipátus; 
 atque ad últimum, cæso cápite, martyrium consummávit.
\switchcolumn
\selectlanguage{english}
At Autun, St. Symphorian, a martyr, 
 in the time of Emperor Aurelian. Refusing to offer sacrifice to the 
 idols, he was first scourged, then confined to prison, and finally ended his 
 martyrdom by being beheaded.
\switchcolumn*
\selectlanguage{latin}
Tudérti, in Umbria, 
 natális sancti Philíppi Benítii, Confessóris, Florentíni, qui Ordinis 
 Servórum beátæ Maríæ Vírginis éxstitit propagátor et exímiæ humilitátis vir; 
 atque a Cleménte Décimo, Pontífice Máximo, Sanctórum número adscríptus est. 
 Ipsíus autem festívitas sequénti die celebrátur.
\switchcolumn
\selectlanguage{english}
At Todi in Umbria, the birthday of 
 St. Philip Beniti, confessor, of Florence. He was a zealous promoter 
 of the Order of the Servants of the Blessed Virgin Mary, and was a man of 
 great humility. He was canonized by Pope Clement X; his feast, 
 however, is observed on the day following
\switchcolumn*
\selectlanguage{latin}
Romæ sancti Antoníni 
 Mártyris, qui cum se Christiánum líbera voce faterétur, senténtia capitáli a 
 Júdice Vitéllio damnátus est, et a Rufíno Presbytero via Aurélia sepúltus.
\switchcolumn
\selectlanguage{english}
At Rome, St. Antoninus, martyr, who, 
 openly declaring himself a Christian, was condemned to capital punishment by 
 the judge Vitellius, and buried on the Aurelian Way.
\switchcolumn*
\selectlanguage{latin}
Tarsi, in Cilícia, commemorátio sanctórum Athanásii, Epíscopi et Mártyris, Anthúsæ, nóbilis 
 féminæ, quam ipse baptizáverat, ac simul Charísii et Neóphyti Mártyrum, 
 ejúsdem Anthúsæ servórum, qui sub Valeriáno Imperatóre passi sunt.
\switchcolumn
\selectlanguage{english}
At Tarsus in Cilicia, the 
 commemoration of Saints Athanasius, bishop and martyr, Anthusa, a noble 
 woman he had baptized, and two of her servants, Charisius and Neophytus, 
 martyrs who suffered under the Emperor Valerian.
\switchcolumn*
\selectlanguage{latin}
In Portu Románo sanctórum Mártyrum Martiális, Saturníni, Epictéti, Maprílis et Felícis, cum 
 Sóciis eórum.
\switchcolumn
\selectlanguage{english}
At Porto, the holy martyrs Martial, 
 Saturninus, Epictetus, Maprilis, and Felix, with their companions.
\switchcolumn*
\selectlanguage{latin}
Nicomedíæ pássio 
 sanctórum Agathoníci, Zótici et Sociórum Mártyrum, sub Maximiáno Imperatóre 
 et Eutólmio Præside.
\switchcolumn
\selectlanguage{english}
At Nicomedia, the passion of Saints 
 Agathonicus, Zoticus, and their fellow-martyrs, under Emperor Maximian and 
 the governor Eutholomius.
\switchcolumn*
\selectlanguage{latin}
Rhemis, in Gállia, 
 sanctórum Mártyrum Mauri et Sociórum.
\switchcolumn
\selectlanguage{english}
At Rheims in France, the holy 
 martyrs Maur and his companions.
\switchcolumn*
\selectlanguage{latin}
In Hispánia sanctórum 
 Mártyrum Fabriciáni et Filibérti.
\switchcolumn
\selectlanguage{english}
In Spain, the holy martyrs Fabrician 
 and Philibert.
\switchcolumn*
\selectlanguage{latin}
Papíæ sancti Gunifórti 
 Mártyris.
\switchcolumn
\selectlanguage{english}
At Pavia, St. Gunifort, martyr.
\switchcolumn*
\selectlanguage{latin}
\end{paracol}


% ---- martyrology/mart08/mart0823.htm
\needspace{10\baselineskip}
\begin{paracol}{2}
\selectlanguage{latin}
\begin{center}{\color{gregoriocolor} Décimo Kaléndas Septémbris. 
 Luna\dots\ }\end{center}
\switchcolumn
\selectlanguage{english}
\begin{center}{\color{gregoriocolor} The 
 23\textsuperscript{rd} Day of 
 August. The\dots\ Day of the Moon.}\end{center}
\end{paracol}

\noindent\begin{tabularx}{\linewidth}{*{19}{>{\centering\arraybackslash}X}}
 \textcolor{gregoriocolor}{a} & \textcolor{gregoriocolor}{b} & \textcolor{gregoriocolor}{c} & \textcolor{gregoriocolor}{d} & \textcolor{gregoriocolor}{e} & \textcolor{gregoriocolor}{f} & \textcolor{gregoriocolor}{g} & \textcolor{gregoriocolor}{h} & \textcolor{gregoriocolor}{i} & \textcolor{gregoriocolor}{k} & \textcolor{gregoriocolor}{l} & \textcolor{gregoriocolor}{m} & \textcolor{gregoriocolor}{n} & \textcolor{gregoriocolor}{p} & \textcolor{gregoriocolor}{q} & \textcolor{gregoriocolor}{r} & \textcolor{gregoriocolor}{s} & \textcolor{gregoriocolor}{t} & \textcolor{gregoriocolor}{u} \\
 29 & 1 & 2 & 3 & 4 & 5 & 6 & 7 & 8 & 9 & 10 & 11 & 12 & 13 & 14 & 15 & 16 & 17 & 18 \\
\end{tabularx}
\vspace{0.5\baselineskip}
\noindent\begin{tabularx}{\linewidth}{*{12}{>{\centering\arraybackslash}X}}
 \textcolor{gregoriocolor}{A} & \textcolor{gregoriocolor}{B} & \textcolor{gregoriocolor}{C} & \textcolor{gregoriocolor}{D} & \textcolor{gregoriocolor}{E} & F & \textcolor{gregoriocolor}{F} & \textcolor{gregoriocolor}{G} & \textcolor{gregoriocolor}{H} & \textcolor{gregoriocolor}{M} & \textcolor{gregoriocolor}{N} & \textcolor{gregoriocolor}{P} \\
 19 & 20 & 21 & 22 & 23 & 24 & 23 & 24 & 25 & 26 & 27 & 28 \\
\end{tabularx}

\begin{paracol}{2}
\selectlanguage{latin}
\lettrine[lines=1]{V}{igília} sancti Bartholomǽi Apóstoli.
\switchcolumn
\selectlanguage{english}
\lettrine[lines=1]{T}{he} Vigil of St. Bartholomew, Apostle.
\switchcolumn*
\selectlanguage{latin}
Sancti Philíppi Benítii, 
 Confessóris, qui Ordinis Servórum beátæ Maríæ Vírginis éxstitit propagátor, 
 ac prídie hujus diéi migrávit ad Dóminum.
\switchcolumn
\selectlanguage{english}
St. Philip Beniti, confessor, 
 promoter of the Order of the Servants of the Blessed Virgin Mary, who 
 departed to the Lord on the previous day.
\switchcolumn*
\selectlanguage{latin}
Apud Ostia Tiberína 
 sanctórum Mártyrum Quiríaci Epíscopi, Máximi Presbyteri, Archelái Diáconi, 
 et Sociórum, qui sub Ulpiáno Præfécto, témpore Alexándri, passi sunt.
\switchcolumn
\selectlanguage{english}
At Ostia, the holy martyrs Quiriacus, 
 bishop, Maximus, priest, Archelaus, deacon, and their companions, who 
 suffered under prefect Ulpian, in the time of Alexander.
\switchcolumn*
\selectlanguage{latin}
Antiochíæ natális 
 sanctórum Mártyrum Restitúti, Donáti, Valeriáni et Fructuósæ, cum áliis 
 duódecim; qui præclaríssimo confessiónis honóre coronáti sunt.
\switchcolumn
\selectlanguage{english}
At Antioch, the birthday of the holy 
 martyrs Restitutus, Donatus, Valerian, and Fructuosa, with twelve others, 
 who were crowned after having distinguished themselves by a glorious 
 confession.
\switchcolumn*
\selectlanguage{latin}
Ægǽæ, in Cilícia, sanctórum Mártyrum fratrum Cláudii, Astérii et Neónis, qui, Christiánæ 
 religiónis a novérca accusáti, sub Diocletiáno Imperatóre et Lysia Præside, 
 post acérba torménta cruci sunt affíxi, in qua victóres cum Christo 
 triumphárunt. Passæ sunt post eos Donvína et Theonílla.
\switchcolumn
\selectlanguage{english}
At Aegaea in Cilicia, the holy 
 martyrs Claudius, Asterius, and Neon, brothers, who were accused of being 
 Christians by their stepmother, under Emperor Diocletian and the governor 
 Lysias. After enduring bitter torments, they were fastened to a cross, 
 and thus conquered and triumphed with Christ. After them suffered 
 Dovina and Theonilla.
\switchcolumn*
\selectlanguage{latin}
Rhemis, in Gállia, 
 natális sanctórum Timóthei et Apollináris, qui, ibídem consummáto martyrio, 
 cæléstia regna meruérunt.
\switchcolumn
\selectlanguage{english}
At Rheims in France, the birthday of 
 the Saints Timothy and Apollinaris, who merited to enter the heavenly 
 kingdom by completing their martyrdom in that city.
\switchcolumn*
\selectlanguage{latin}
Lugdúni, in Gállia, 
 sanctórum Mártyrum Minérvi, et Eleazári cum fíliis octo.
\switchcolumn
\selectlanguage{english}
At Lyons, the holy martyrs Minercus 
 and Eleazar, with his eight sons.
\switchcolumn*
\selectlanguage{latin}
Item sancti Luppi 
 Mártyris, qui, ex sérvili conditióne, Christi libertáte donátus, martyrii 
 quoque coróna dignátus est.
\switchcolumn
\selectlanguage{english}
Also St. Luppus, martyr, who, though 
 a slave, enjoyed the liberty of Christ, and was likewise deemed worthy of 
 the crown of martyrdom.
\switchcolumn*
\selectlanguage{latin}
Hierosólymis sancti 
 Zachæi Epíscopi, qui, quartus a beáto Jacóbo Apóstolo, Hierosolymitánum 
 Ecclésiam rexit.
\switchcolumn
\selectlanguage{english}
At Jerusalem, St. Zachaeus, bishop, 
 who governed the Church in that city the fourth after the blessed apostle 
 James.
\switchcolumn*
\selectlanguage{latin}
Alexandríæ sancti 
 Theónæ, Epíscopi et Confessóris.
\switchcolumn
\selectlanguage{english}
At Alexandria, St. Theonas, bishop 
 and confessor.
\switchcolumn*
\selectlanguage{latin}
Uticæ, in Africa, beáti 
 Victóris Epíscopi.
\switchcolumn
\selectlanguage{english}
At Utica in Africa, blessed Victor, 
 bishop.
\switchcolumn*
\selectlanguage{latin}
Augustodúni sancti 
 Flaviáni Epíscopi.
\switchcolumn
\selectlanguage{english}
At Autun, St. Flavian, bishop.
\switchcolumn*
\selectlanguage{latin}
\end{paracol}


% ---- martyrology/mart08/mart0824.htm
\needspace{10\baselineskip}
\begin{paracol}{2}
\selectlanguage{latin}
\begin{center}{\color{gregoriocolor} Nono Kaléndas Septémbris. 
 Luna\dots\ }\end{center}
\switchcolumn
\selectlanguage{english}
\begin{center}{\color{gregoriocolor} The 
 24\textsuperscript{th} Day of 
 August. The\dots\ Day of the Moon.}\end{center}
\end{paracol}

\noindent\begin{tabularx}{\linewidth}{*{19}{>{\centering\arraybackslash}X}}
 \textcolor{gregoriocolor}{a} & \textcolor{gregoriocolor}{b} & \textcolor{gregoriocolor}{c} & \textcolor{gregoriocolor}{d} & \textcolor{gregoriocolor}{e} & \textcolor{gregoriocolor}{f} & \textcolor{gregoriocolor}{g} & \textcolor{gregoriocolor}{h} & \textcolor{gregoriocolor}{i} & \textcolor{gregoriocolor}{k} & \textcolor{gregoriocolor}{l} & \textcolor{gregoriocolor}{m} & \textcolor{gregoriocolor}{n} & \textcolor{gregoriocolor}{p} & \textcolor{gregoriocolor}{q} & \textcolor{gregoriocolor}{r} & \textcolor{gregoriocolor}{s} & \textcolor{gregoriocolor}{t} & \textcolor{gregoriocolor}{u} \\
 1 & 2 & 3 & 4 & 5 & 6 & 7 & 8 & 9 & 10 & 11 & 12 & 13 & 14 & 15 & 16 & 17 & 18 & 19 \\
\end{tabularx}
\vspace{0.5\baselineskip}
\noindent\begin{tabularx}{\linewidth}{*{12}{>{\centering\arraybackslash}X}}
 \textcolor{gregoriocolor}{A} & \textcolor{gregoriocolor}{B} & \textcolor{gregoriocolor}{C} & \textcolor{gregoriocolor}{D} & \textcolor{gregoriocolor}{E} & F & \textcolor{gregoriocolor}{F} & \textcolor{gregoriocolor}{G} & \textcolor{gregoriocolor}{H} & \textcolor{gregoriocolor}{M} & \textcolor{gregoriocolor}{N} & \textcolor{gregoriocolor}{P} \\
 20 & 21 & 22 & 23 & 24 & 25 & 24 & 25 & 26 & 27 & 28 & 29 \\
\end{tabularx}

\begin{paracol}{2}
\selectlanguage{latin}
\lettrine[lines=2]{S}{ancti} Bartholomǽi 
 Apóstoli, qui Christi Evangélium in India prædicávit; inde in majórem 
 Arméniam proféctus, ibi, cum plúrimos ad fidem convertísset, vivus a 
 bárbaris decoriátus est, atque, Astyagis Regis jussu, cápitis decollatióne 
 martyrium complévit. Ipsíus sacrum corpus, primo ad Líparam ínsulam, 
 deínde Benevéntum, postrémo Romam ad Tiberínam translátum ínsulam, ibi pia 
 fidélium veneratióne honorátur.
\switchcolumn
\selectlanguage{english}
\lettrine[lines=2]{T}{he} apostle St. Bartholomew, who 
 preached the Gospel of Christ in India. Passing thence into the 
 Greater Armenia where, after converting many to the faith, he was flayed 
 alive by the barbarians, and having his head cut off by order of King 
 Astyages, he fulfilled his martyrdom. His holy body was first carried 
 to the island of Lipara, then to Benevento, and finally to Rome in the 
 Island of the Tiber, where it is venerated by the pious faithful.
\switchcolumn*
\selectlanguage{latin}
Limæ, in Perúvia, 
 natális sanctæ Rosæ a sancta María, Vírginis, e tértio Ordine sancti 
 Domínici. Ejus vero festívitas tértio Kaléndas Septémbris celebrátur.
\switchcolumn
\selectlanguage{english}
At Lima in Peru, the birthday of St. 
 Rose of St. Mary, virgin of the Third Order of St. Dominic. Her feast 
 is observed on the 30th of August.
\switchcolumn*
\selectlanguage{latin}
Népete sancti Ptolomǽi 
 Epíscopi, qui fuit discípulus beáti Petri Apóstoli; atque, ab eo missus in 
 Túsciam ad prædicándum Evangélium, in eádem civitáte gloriósus Christi 
 Martyr occúbuit.
\switchcolumn
\selectlanguage{english}
At Nepi, St. Ptolemy, bishop, 
 disciple of the blessed apostle Peter. Being sent by him to preach the 
 Gospel in Tuscany, he died a glorious martyr of Christ in the city of Nepi.
\switchcolumn*
\selectlanguage{latin}
Eódem die sancti 
 Eutychii, qui fuit discípulus beáti Joánnis Evangelístæ; atque, ob Evangélii 
 prædicatiónem in multis regiónibus cárceres, vérbera et ignes perpéssus, in 
 pace tandem quiévit.
\switchcolumn
\selectlanguage{english}
Also, St. Eutychius, disciple of the 
 blessed evangelist John. He preached the Gospel in many countries, and 
 was subjected to imprisonment, to stripes and fire, but finally he rested in 
 peace.
\switchcolumn*
\selectlanguage{latin}
Népete sancti Románi, 
 ejúsdem civitátis Epíscopi, qui, cum esset sancti Ptolomǽi discípulus, fuit 
 étiam in passióne sócius.
\switchcolumn
\selectlanguage{english}
Also at Nepi, St. Romanus, bishop of 
 that city, who was the disciple of St. Ptolemy, and his companion in 
 martyrdom.
\switchcolumn*
\selectlanguage{latin}
Carthágine sanctórum 
 trecentórum Mártyrum, témpore Valeriáni et Galliéni. Hi Mártyres 
 magnánimi, inter ália supplícia, cum Præses fornácem calcáriam accéndi 
 jussísset, et, in præséntia ejus, prunas cum thure exhibéri, atque illis 
 dixísset: « Elígite e duóbus unum, aut thura super his carbónibus offérte 
 Jovi, aut in calcem demergímini », fide armáti, Christum Dei Fílium 
 confiténtes, ictu rapidíssimo se injecérunt in ignem, et inter calcis 
 vapóres in púlverem sunt redácti; ex quo candidátus ille beatórum exércitus 
 appellári Massa cándida méruit.
\switchcolumn
\selectlanguage{english}
At Carthage, three hundred holy 
 martyrs, in the time of Valerian and Gallienus. Among other torments 
 inflicted on them, a pit filled with burning lime was prepared by order of 
 the governor, who, live coals with incense being brought to him, said to the 
 confessors: ``Choose one of these two things: either offer incense to Jupiter 
 upon these coals, or be thrown into the lime.'' Armed with faith, and 
 confessing Christ to be the Son of God, they quickly threw themselves into 
 the pit, and amid the vapours of the lime were reduced to dust. From 
 this circumstance, this white-robed company of the blessed earned for itself 
 the name of the White Mass.
\switchcolumn*
\selectlanguage{latin}
In Isáuria sancti 
 Tatiónis Mártyris, qui, in persecutióne Diocletiáni, sub Urbáno Præside, 
 gládio cæsus, martyrii corónam accépit.
\switchcolumn
\selectlanguage{english}
In Isauria, St. Tation, martyr, who 
 received the crown of martyrdom by being beheaded in the persecution of 
 Diocletian, under the governor Urbanus.
\switchcolumn*
\selectlanguage{latin}
Item sancti Geórgii 
 Limniótæ Mónachi, qui, cum ímpium Leónem Imperatórem, quod sacras Imágines, 
 frángeret Sanctorúmque relíquias combúreret, reprehendísset, hanc ob causam, 
 ejus jussu mánibus abscíssis et cápite incénso, Martyr migrávit ad Dóminum.
\switchcolumn
\selectlanguage{english}
Also, St. George Limniota, monk. 
 Because he reprehended the wicked emperor Leo for breaking holy images, and 
 burning the relics of the saints, he had his hands cut off and his head 
 burned by order of the tyrant, and went to our Lord to receive the 
 recompence of a martyr.
\switchcolumn*
\selectlanguage{latin}
Apud Ostia Tiberína 
 sanctæ Aureæ, Vírginis et Mártyris; quæ saxo ad collum ligáto, in mare 
 demérsa est. Ipsíus autem corpus, ejéctum ad littus, beátus Nonnus 
 sepelívit.
\switchcolumn
\selectlanguage{english}
At Ostia, on the Tiber, St. Aurea, 
 virgin and martyr, who was plunged into the sea with a stone tied to her 
 neck. Her body being driven to the shore was buried by blessed Nonnus.
\switchcolumn*
\selectlanguage{latin}
Rotómagi sancti Audoéni, 
 Epíscopi et Confessóris.
\switchcolumn
\selectlanguage{english}
At Rouen, St. Owen, bishop and 
 confessor.
\switchcolumn*
\selectlanguage{latin}
Nivérnis, in Gállia, 
 sancti Patrícii Abbátis.
\switchcolumn
\selectlanguage{english}
At Nevers in France, St. Patrick, 
 abbot.
\switchcolumn*
\selectlanguage{latin}
Neápoli in Campánia, sanctæ Joánnæ Antidæ Thouret, Vírginis, Institúti Sorórum a Caritáte 
 Fundatrícis, quam Pius Papa Undécimus in album sanctárum Vírginum rétulit.
\switchcolumn
\selectlanguage{english}
At Naples in Campania, St. Joan 
 Antide Thouret, virgin, who founded the Daughters of Saint Vincent de Paul, 
 and whom Pope Pius XI added to the catalogue of holy virgins.
\switchcolumn*
\selectlanguage{latin}
Massíliæ, in Gállia, 
 sanctæ Æmíliæ de Vialár, Vírginis, Fundatrícis Institúti Sorórum a sancto 
 Joseph ab Apparitióne, fortitúdine, patiéntia et caritáte insígnis, quam 
 Pius Duodécimus, Póntifex Máximus, in Sanctárum númerum rétulit.
\switchcolumn
\selectlanguage{english}
At Marseilles in France, St. Emily 
 de Vialar, virgin, foundress of the Congregation of the Sisters of Saint 
 Joseph of the Apparition. A shining example of fortitude, patience and 
 charity, the Sovereign Pontiff Pius XII added her to the number of the 
 saints.
\switchcolumn*
\selectlanguage{latin}
Valéntiæ, in Hispánia, 
 natális sanctæ Maríæ Michaélæ, Vírginis, Fundatrícis Congregatiónis 
 Ancillárum a Sanctíssimo Sacraménto et Caritátis, patiéndi stúdio ac desidério ánimas Deo lucrándi inflammátæ, quam Pius Papa Undécimus sanctis 
 Virgínibus accénsuit.
\switchcolumn
\selectlanguage{english}
At Valencia in Spain, the birthday 
 of St. Mary Micaela, virgin, who founded the Institute of Religious 
 Adorer-Slaves of the Blessed Sacrament and of Charity. Burning with 
 the desire to suffer and draw souls to God, she was numbered among the holy 
 virgins by Pope Pius XI.
\switchcolumn*
\selectlanguage{latin}
\end{paracol}


% ---- martyrology/mart08/mart0825.htm
\needspace{10\baselineskip}
\begin{paracol}{2}
\selectlanguage{latin}
\begin{center}{\color{gregoriocolor} Octávo Kaléndas Septémbris. 
 Luna\dots\ }\end{center}
\switchcolumn
\selectlanguage{english}
\begin{center}{\color{gregoriocolor} The 
 25\textsuperscript{th} Day of 
 August. The\dots\ Day of the Moon.}\end{center}
\end{paracol}

\noindent\begin{tabularx}{\linewidth}{*{19}{>{\centering\arraybackslash}X}}
 \textcolor{gregoriocolor}{a} & \textcolor{gregoriocolor}{b} & \textcolor{gregoriocolor}{c} & \textcolor{gregoriocolor}{d} & \textcolor{gregoriocolor}{e} & \textcolor{gregoriocolor}{f} & \textcolor{gregoriocolor}{g} & \textcolor{gregoriocolor}{h} & \textcolor{gregoriocolor}{i} & \textcolor{gregoriocolor}{k} & \textcolor{gregoriocolor}{l} & \textcolor{gregoriocolor}{m} & \textcolor{gregoriocolor}{n} & \textcolor{gregoriocolor}{p} & \textcolor{gregoriocolor}{q} & \textcolor{gregoriocolor}{r} & \textcolor{gregoriocolor}{s} & \textcolor{gregoriocolor}{t} & \textcolor{gregoriocolor}{u} \\
 2 & 3 & 4 & 5 & 6 & 7 & 8 & 9 & 10 & 11 & 12 & 13 & 14 & 15 & 16 & 17 & 18 & 19 & 20 \\
\end{tabularx}
\vspace{0.5\baselineskip}
\noindent\begin{tabularx}{\linewidth}{*{12}{>{\centering\arraybackslash}X}}
 \textcolor{gregoriocolor}{A} & \textcolor{gregoriocolor}{B} & \textcolor{gregoriocolor}{C} & \textcolor{gregoriocolor}{D} & \textcolor{gregoriocolor}{E} & F & \textcolor{gregoriocolor}{F} & \textcolor{gregoriocolor}{G} & \textcolor{gregoriocolor}{H} & \textcolor{gregoriocolor}{M} & \textcolor{gregoriocolor}{N} & \textcolor{gregoriocolor}{P} \\
 21 & 22 & 23 & 24 & 25 & 26 & 25 & 26 & 27 & 28 & 29 & 1 \\
\end{tabularx}

\begin{paracol}{2}
\selectlanguage{latin}
\lettrine[lines=2]{A}{pud} Cartháginem sancti 
 Ludovíci Noni, Regis Francórum et Confessóris, vitæ sanctitáte ac 
 miraculórum glória præclári; cujus ossa póstmodum Lutétiam Parisiórum sunt 
 reláta.
\switchcolumn
\selectlanguage{english}
\lettrine[lines=2]{A}{t} Carthage, St. Louis IX, king of 
 France and confessor, illustrious for holiness of life and glorious 
 miracles. His bones were later translated to Paris.
\switchcolumn*
\selectlanguage{latin}
Romæ natális sancti 
 Joséphi Calasánctii, Presbyteri et Confessóris, vitæ innocéntia et miráculis 
 illústris; qui, ad erudiéndam pietáte ac lítteris juventútem, Ordinem 
 Clericórum Regulárium Páuperum Matris Dei Scholárum Piárum fundávit. 
 Eum Pius Duodécimus, Póntifex Máximus, ómnium Scholárum populárium 
 christianárum ubíque exsisténtium cæléstem apud Deum Patrónum constítuit. 
 Ipsíus tamen festívitas sexto Kaléndas Septémbris recólitur.
\switchcolumn
\selectlanguage{english}
At Rome, the birthday of St. Joseph 
 Calasanctius, priest and confessor, noteworthy for his holy life and 
 miracles. He founded the Order of Poor Clerics Regular of the Mother 
 of God of the Christian Schools. The Sovereign Pontiff, Pius XII, 
 named him as heavenly patron of all Christian schoolchildren. His 
 feast is on the 27th of August.
\switchcolumn*
\selectlanguage{latin}
Item Romæ sanctórum 
 Mártyrum Eusébii, Pontiáni, Vincéntii et Peregríni; qui, sub Cómmodo 
 Imperatóre, primum in equúleo leváti, nervis quoque disténti, ac deínde 
 fústibus cæsi sunt, flammis circa eórum látera appósitis; et, cum in laude 
 Christi fidelíssime permanérent, plumbátis usque ad emissiónem spíritus sunt 
 mactáti.
\switchcolumn
\selectlanguage{english}
Also at Rome, in the time of Emperor 
 Commodus, the holy martyrs Eusebius, Pontian, Vincent, and Peregrinus, who 
 were first racked, distended by ropes, then beaten with rods and burned 
 about their sides. As they continued faithfully to praise Christ, they 
 were scourged with leaded whips until they expired.
\switchcolumn*
\selectlanguage{latin}
Romæ prætérea natális 
 beáti Nemésii Diáconi, et fíliæ Lucillæ Vírginis; qui, cum de fide Christi 
 flecti nequáquam possent, decolláti sunt, jubénte Valeriáno Imperatóre. 
 Ipsórum córpora, a beáto Stéphano Papa sepúlta, deínde a beáto Xysto Secúndo 
 via Appia, prídie Kaléndas Novémbris, honéstius tumuláta, Gregórius Quintus 
 in Diacóniam sanctæ Maríæ Novæ tránstulit, una cum sanctis Symphrónio, 
 Olympio Tribúno, hujúsque uxóre Exsupéria et Theodúlo fílio; qui omnes, 
 Symphrónii ópera convérsi et ab eódem sancto Stéphano baptizáti, martyrio 
 coronáti fúerant. Eadem Sanctórum córpora, Gregório Décimo tértio 
 Summo Pontífice, ibídem invénta, sub altári ejúsdem Ecclésiæ honorificéntius 
 collocáta sunt sexto Idus Decémbris.
\switchcolumn
\selectlanguage{english}
In the same city of Rome, the 
 birthday of blessed Nemesius, deacon, and his daughter, the virgin Lucilla. 
 As they could not be prevailed upon to abandon the faith of Christ, they 
 were beheaded by order of Emperor Valerian. Their bodies were buried 
 by blessed Pope Stephen, and afterwards more decently entombed on the 31st 
 of October, by blessed Sixtus on the Appian Way. Gregory V translated 
 them into the sacristy of Santa Maria Nova, together with the Saints 
 Symphronius, Olympius, a tribune, Exuperia, his wife, and Theodulus, his 
 son, who, being all converted by the exertions of Symphonius, and baptized 
 by the same St. Stephen, had been crowned with martyrdom. These holy 
 bodies were found there during the pontificate of Gregory XIII, and placed 
 more honourably beneath the altar of the same church, on the 8th of 
 December.
\switchcolumn*
\selectlanguage{latin}
Item Romæ sancti 
 Genésii Mártyris, qui, primum sub Gentilitáte mimus, cum in theátro, 
 spectánte Diocletiáno Imperatóre, Mystériis Christianórum illúderet, repénte, 
 inspirátus a Deo, convérsus est ad fidem et baptizátus. Mox, 
 Imperatóris jussu, fústibus crudelíssime cæsus, deínde suspénsus in equúleo, 
 et ungulárum diutíssima laceratióne vexátus, lampádibus étiam adústus est; 
 ac tandem, cum in fide Christi persísteret, dicens: « Non est Rex præter 
 Christum, pro quo si míllies occídar, ipsum mihi de ore, ipsum mihi de corde 
 auférre non potéritis », martyrii palmam obtruncatióne cápitis proméruit.
\switchcolumn
\selectlanguage{english}
Also at Rome, St. Genesius, martyr, 
 who had embraced the profession of actor while he was a pagan. One day 
 he was deriding the Christian mysteries in the theatre in the presence of 
 Emperor Diocletian; but by the inspiration of God he was suddenly converted 
 to the faith and baptized. By command of the emperor he was forthwith 
 most cruelly beaten with rods, then racked, and a long time lacerated with 
 iron hooks, and burned with torches. As he remained firm in the faith 
 of Christ, even saying: ``There is no king besides Christ. Should you 
 kill me a thousand times, you shall not be able to take him from my lips or 
 my heart.'' He was then beheaded, and thus merited the palm of 
 martyrdom.
\switchcolumn*
\selectlanguage{latin}
Areláte, in Gállia, 
 beáti item Genésii, qui, cum ímpia edícta, quibus Christiáni puníri 
 jubebántur, exceptóris offício fungens, nollet excípere, et, projéctis in 
 públicum tábulis, se Christiánum esse testarétur, comprehénsus et decollátus 
 est, atque ita martyrii glóriam, próprio cruóre baptizátus, accépit.
\switchcolumn
\selectlanguage{english}
At Arles in France, another blessed 
 Genesius, who, filling the office of notary, and refusing to record the 
 impious edicts by which Christians were commanded to be punished, threw away 
 his books publicly, and declared himself a Christian. He was seized 
 and beheaded, and thus attained the glory of martyrdom through baptism in 
 his own blood.
\switchcolumn*
\selectlanguage{latin}
In Syria sancti Juliáni 
 Mártyris.
\switchcolumn
\selectlanguage{english}
In Syria, St. Julian, martyr.
\switchcolumn*
\selectlanguage{latin}
Tarracóne, in Hispánia, 
 sancti Magíni Mártyris.
\switchcolumn
\selectlanguage{english}
At Tarragona in Spain, St. Maginus, 
 martyr.
\switchcolumn*
\selectlanguage{latin}
Itálicæ, in Hispánia, 
 sancti Gerúntii Epíscopi, qui, Apostolórum témpore, Evangélium in ea 
 província prædicávit, et, post multos labóres, in cárcere quiévit.
\switchcolumn
\selectlanguage{english}
At Italica in Spain, St. Gerontius, 
 bishop, who preached the Gospel in that country in apostolic times, and 
 after many labours died in prison.
\switchcolumn*
\selectlanguage{latin}
Constantinópoli sancti 
 Mennæ Epíscopi.
\switchcolumn
\selectlanguage{english}
At Constantinople, St. Mennas, 
 bishop.
\switchcolumn*
\selectlanguage{latin}
Trajécti sancti 
 Gregórii Epíscopi.
\switchcolumn
\selectlanguage{english}
At Utrecht, St. Gregory, bishop.
\switchcolumn*
\selectlanguage{latin}
Apud Montem Falíscum, 
 in Etrúria, sancti Thomæ, Confessóris, qui Herfordiénsis Ecclésiæ, in 
 Anglia, Epíscopus éxstitit.
\switchcolumn
\selectlanguage{english}
At Monte Falisco in Etruria, St. 
 Thomas, bishop of the church of Hereford in England, and confessor.
\switchcolumn*
\selectlanguage{latin}
Neápoli, in Campánia, 
 sanctæ Patríciæ Vírginis.
\switchcolumn
\selectlanguage{english}
At Naples in Campania, St. Patricia, 
 virgin.
\switchcolumn*
\selectlanguage{latin}
\end{paracol}


% ---- martyrology/mart08/mart0826.htm
\needspace{10\baselineskip}
\begin{paracol}{2}
\selectlanguage{latin}
\begin{center}{\color{gregoriocolor} Séptimo Kaléndas Septémbris. 
 Luna\dots\ }\end{center}
\switchcolumn
\selectlanguage{english}
\begin{center}{\color{gregoriocolor} The 
 26\textsuperscript{th} Day of 
 August. The\dots\ Day of the Moon.}\end{center}
\end{paracol}

\noindent\begin{tabularx}{\linewidth}{*{19}{>{\centering\arraybackslash}X}}
 \textcolor{gregoriocolor}{a} & \textcolor{gregoriocolor}{b} & \textcolor{gregoriocolor}{c} & \textcolor{gregoriocolor}{d} & \textcolor{gregoriocolor}{e} & \textcolor{gregoriocolor}{f} & \textcolor{gregoriocolor}{g} & \textcolor{gregoriocolor}{h} & \textcolor{gregoriocolor}{i} & \textcolor{gregoriocolor}{k} & \textcolor{gregoriocolor}{l} & \textcolor{gregoriocolor}{m} & \textcolor{gregoriocolor}{n} & \textcolor{gregoriocolor}{p} & \textcolor{gregoriocolor}{q} & \textcolor{gregoriocolor}{r} & \textcolor{gregoriocolor}{s} & \textcolor{gregoriocolor}{t} & \textcolor{gregoriocolor}{u} \\
 3 & 4 & 5 & 6 & 7 & 8 & 9 & 10 & 11 & 12 & 13 & 14 & 15 & 16 & 17 & 18 & 19 & 20 & 21 \\
\end{tabularx}
\vspace{0.5\baselineskip}
\noindent\begin{tabularx}{\linewidth}{*{12}{>{\centering\arraybackslash}X}}
 \textcolor{gregoriocolor}{A} & \textcolor{gregoriocolor}{B} & \textcolor{gregoriocolor}{C} & \textcolor{gregoriocolor}{D} & \textcolor{gregoriocolor}{E} & F & \textcolor{gregoriocolor}{F} & \textcolor{gregoriocolor}{G} & \textcolor{gregoriocolor}{H} & \textcolor{gregoriocolor}{M} & \textcolor{gregoriocolor}{N} & \textcolor{gregoriocolor}{P} \\
 22 & 23 & 24 & 25 & 26 & 27 & 26 & 27 & 28 & 29 & 1 & 2 \\
\end{tabularx}

\begin{paracol}{2}
\selectlanguage{latin}
\lettrine[lines=2]{S}{ancti} Zephyríni, Papæ 
 et Mártyris; cujus dies natális tertiodécimo Kaléndas Januárii recensétur.
\switchcolumn
\selectlanguage{english}
\lettrine[lines=2]{A}{t} Rome, St. Zephyrinus, pope and 
 martyr, whose birthday falls on the 20th of December.
\switchcolumn*
\selectlanguage{latin}
Cardónæ, in Hispánia, 
 tránsitus sancti Raymúndi Nonnáti, Cardinális et Confessóris, ex Ordine 
 beátæ Maríæ de Mercéde redemptiónis captivórum, vitæ sanctitáte et miráculis 
 clari. Ipsíus tamen festum recólitur prídie Kaléndas Septémbris.
\switchcolumn
\selectlanguage{english}
At Cardona in Spain, the birthday of 
 St. Raymund Nonnatus, cardinal and confessor, of the Order of our Lady of 
 Ransom for the Redemption of Captives, renowned for holiness of life and for 
 miracles, whose feast is observed on the 31st of August.
\switchcolumn*
\selectlanguage{latin}
Romæ sanctórum Mártyrum 
 Irenǽi et Abúndii, qui, in persecutióne Valeriáni, eo quod corpus beátæ 
 Concórdiæ, in cloácam projéctum, leváverant, in eándem cloácam demérsi 
 fuérunt; quorum córpora, a Justíno Presbytero inde extrácta, in crypta, 
 juxta beátum Lauréntium, sepúlta sunt.
\switchcolumn
\selectlanguage{english}
At Rome, during the persecution of 
 Valerian, the holy martyrs Irenaeus and Abundius, who were thrown into a 
 sewer from which they had taken the body of blessed of Concordia. 
 Their bodies were drawn out by the priest Justin, and buried in a crypt near 
 St. Lawrence.
\switchcolumn*
\selectlanguage{latin}
Apud Albintimélium, 
 Ligúriæ civitátem, sancti Secúndi Mártyris, viri spectábilis et Ducis ex 
 legióne Thebæórum.
\switchcolumn
\selectlanguage{english}
At Ventimiglia, a city of Liguria, 
 St. Secundus, martyr, a distinguished man and officer in the Theban Legion.
\switchcolumn*
\selectlanguage{latin}
Bérgomi sancti 
 Alexándri Mártyris, qui, et ipse unus ex eádem legióne, cum nomen Dómini 
 Jesu Christi constantíssime faterétur, cápitis abscissióne martyrium 
 complévit.
\switchcolumn
\selectlanguage{english}
At Bergamo in Lombardy, St. 
 Alexander, martyr, who was one of the same legion, and endured martyrdom, 
 being beheaded for the constant confession of the name of our Lord Jesus 
 Christ.
\switchcolumn*
\selectlanguage{latin}
Apud Marsos sanctórum 
 Simplícii, et ejus filiórum Constántii et Victoriáni, qui, sub Antoníno 
 Imperatóre, várie primum excruciáti, tum demum, secúris ictu percússi, 
 martyrii corónam, adépti sunt.
\switchcolumn
\selectlanguage{english}
Among the Marcians, the saints 
 Simplicius, and his sons Constantius and Victorian, who were first tortured 
 in different manners, and lastly, struck with the axe, obtained the crown of 
 martyrdom, in the time of Emperor Antoninus.
\switchcolumn*
\selectlanguage{latin}
Nicomedíæ pássio sancti 
 Hadriáni, e Probo Cæsare progéniti, qui, Licínio persecutiónem in 
 Christiános commótam éxprobrans, ab eódem jussus est occídi. Ipsíus 
 corpus Domítius, Byzántii Epíscopus, ejus pátruus, in ipsíus civitátis 
 subúrbio cui nomen Argyrópolis sepelívit.
\switchcolumn
\selectlanguage{english}
At Nicomedia, the martyrdom of St. 
 Adrian, son of Emperor Probus. For reproaching Licinius because of the 
 persecution of Christians, he was put to death by his order. His body 
 was buried at Argyropolis by his uncle Domitius, bishop of Byzantium.
\switchcolumn*
\selectlanguage{latin}
In Hispánia sancti 
 Victoris Mártyris, qui pro Christi fide, a Mauris occísus, martyrii coróna 
 donátus est.
\switchcolumn
\selectlanguage{english}
In Spain, St. Victor, martyr, who 
 merited the crown of martyrs by being slain by the Moors for the faith of 
 Christ.
\switchcolumn*
\selectlanguage{latin}
Cápuæ sancti Rufíni, 
 Epíscopi et Confessóris.
\switchcolumn
\selectlanguage{english}
At Capua, St. Rufinus, bishop and 
 confessor.
\switchcolumn*
\selectlanguage{latin}
Pistórii, in Túscia, 
 sancti Felícis, Presbyteri et Confessóris.
\switchcolumn
\selectlanguage{english}
At Pistoia, St. Felix, priest and 
 confessor.
\switchcolumn*
\selectlanguage{latin}
Pódii, in diœcési 
 Pictaviénsi, sanctæ Joánnæ-Elisabeth Bichier des Ages, Vírginis, 
 Congregatiónis Filiárum a Cruce una cum sancto Andréa Hubérto Fournet 
 Fundatrícis, jugi mortificatióne et vitæ innocéntia claræ, quam Pius Papa 
 Duodécimus sanctárum Vírginum fastis accénsuit.
\switchcolumn
\selectlanguage{english}
In the diocese of Poitiers, St. 
 Joan-Elizabeth Bichier des Ages, virgin, who with St. André Hubert Fournet 
 co-founded the Congregation of the Daughters of the Cross, and who was 
 renowned for her spirit of mortification and life of innocence. Pope 
 Pius XII added her name to the list of holy virgins.
\switchcolumn*
\selectlanguage{latin}
\end{paracol}


% ---- martyrology/mart08/mart0827.htm
\needspace{10\baselineskip}
\begin{paracol}{2}
\selectlanguage{latin}
\begin{center}{\color{gregoriocolor} Sexto Kaléndas Septémbris. 
 Luna\dots\ }\end{center}
\switchcolumn
\selectlanguage{english}
\begin{center}{\color{gregoriocolor} The 
 27\textsuperscript{th} Day of 
 August. The\dots\ Day of the Moon.}\end{center}
\end{paracol}

\noindent\begin{tabularx}{\linewidth}{*{19}{>{\centering\arraybackslash}X}}
 \textcolor{gregoriocolor}{a} & \textcolor{gregoriocolor}{b} & \textcolor{gregoriocolor}{c} & \textcolor{gregoriocolor}{d} & \textcolor{gregoriocolor}{e} & \textcolor{gregoriocolor}{f} & \textcolor{gregoriocolor}{g} & \textcolor{gregoriocolor}{h} & \textcolor{gregoriocolor}{i} & \textcolor{gregoriocolor}{k} & \textcolor{gregoriocolor}{l} & \textcolor{gregoriocolor}{m} & \textcolor{gregoriocolor}{n} & \textcolor{gregoriocolor}{p} & \textcolor{gregoriocolor}{q} & \textcolor{gregoriocolor}{r} & \textcolor{gregoriocolor}{s} & \textcolor{gregoriocolor}{t} & \textcolor{gregoriocolor}{u} \\
 4 & 5 & 6 & 7 & 8 & 9 & 10 & 11 & 12 & 13 & 14 & 15 & 16 & 17 & 18 & 19 & 20 & 21 & 22 \\
\end{tabularx}
\vspace{0.5\baselineskip}
\noindent\begin{tabularx}{\linewidth}{*{12}{>{\centering\arraybackslash}X}}
 \textcolor{gregoriocolor}{A} & \textcolor{gregoriocolor}{B} & \textcolor{gregoriocolor}{C} & \textcolor{gregoriocolor}{D} & \textcolor{gregoriocolor}{E} & F & \textcolor{gregoriocolor}{F} & \textcolor{gregoriocolor}{G} & \textcolor{gregoriocolor}{H} & \textcolor{gregoriocolor}{M} & \textcolor{gregoriocolor}{N} & \textcolor{gregoriocolor}{P} \\
 23 & 24 & 25 & 26 & 27 & 28 & 27 & 28 & 29 & 1 & 2 & 3 \\
\end{tabularx}

\begin{paracol}{2}
\selectlanguage{latin}
\lettrine[lines=2]{S}{ancti} Joséphi 
 Calasánctii, Presbyteri et Confessóris, qui Ordinis Clericórum Regulárium 
 Páuperum Matris Dei Scholárum Piárum éxstitit Fundátor, atque octávo 
 Kaléndas Septémbris obdormívit in Dómino.
\switchcolumn
\selectlanguage{english}
\lettrine[lines=2]{S}{t.} Joseph Calasanctius, priest and 
 confessor, who founded the Order of Poor Clerics Regular of the Mother of 
 God of the Christian Schools. He fell asleep in the Lord on the 25th 
 of August.
\switchcolumn*
\selectlanguage{latin}
Poténtiæ, in Lucánia, 
 pássio sanctórum Aróntii, Honoráti, Fortunáti et Sabiniáni; qui, sanctórum 
 Bonifátii et Theclæ fílii, a Valeriáno Júdice, sub Maximiáno Imperatóre, 
 jussi sunt capitálem subíre senténtiam. Eórum tamen ac reliquórum ex 
 duódecim frátribus festum Kaléndis Septémbris celebrátur.
\switchcolumn
\selectlanguage{english}
At Potenza in Lucania, the passion 
 of Saints Arontius, Honoratus, Fortunatus, and Sabinian. They were the 
 sons of Saints Boniface and Thecla, and were condemned to death by the judge 
 Valerian in the reign of Emperor Maximian. Their feast, together with 
 that of the other twelve holy brethren, is celebrated on the first of 
 September.
\switchcolumn*
\selectlanguage{latin}
Bérgomi sancti Narni, 
 qui, a beáto Bárnaba baptizátus, primus ab ipso ejúsdem civitátis Epíscopus 
 ordinátus est.
\switchcolumn
\selectlanguage{english}
At Bergamo, St. Narnus, who was 
 baptized by blessed Barnabas and consecrated by him first bishop of that 
 city.
\switchcolumn*
\selectlanguage{latin}
Cápuæ natális sancti 
 Rufi, Epíscopi et Mártyris; qui, cum esset patríciæ dignitátis, a beáto 
 Apollináre, sancti Petri discípulo, cum univérsa família baptizátus est.
\switchcolumn
\selectlanguage{english}
At Capua, the birthday of St. Rufus, 
 bishop and martyr, a patrician, who was baptized with all his family by 
 blessed Apollinaris, disciple of St. Peter.
\switchcolumn*
\selectlanguage{latin}
Ibídem sanctórum 
 Mártyrum Rufi et Carpóphori, qui sub Diocletiáno et Maximiáno passi sunt.
\switchcolumn
\selectlanguage{english}
In the same place, the holy martyrs 
 Rufus and Carpophorus, who suffered under Diocletian and Maximian.
\switchcolumn*
\selectlanguage{latin}
Tomis, in Ponto, 
 sanctórum Mártyrum Marcellíni Tribúni, et uxóris Mannéæ, ac filiórum Joánnis, 
 Serapiónis et Petri.
\switchcolumn
\selectlanguage{english}
At Tomis in Pontus, the holy martyrs 
 Marcellinus, a tribune, and Mannea, his wife, and his sons John, Serapion, 
 and Peter.
\switchcolumn*
\selectlanguage{latin}
Apud Leonínos, in 
 Sicília, sanctæ Eutháliæ Vírginis, quæ, cum esset Christiána, ad cæléstem 
 Sponsum, a fratre suo Sermiliáno cæsa gládio, migrávit.
\switchcolumn
\selectlanguage{english}
At Lentini in Sicily, St. Euthalia, 
 virgin. Because she was a Christian she was put to the sword by her 
 brother Sermilian, and went to her Spouse.
\switchcolumn*
\selectlanguage{latin}
Eódem die pássio sanctæ 
 Anthúsæ junióris, quæ ob Christi fidem, in púteum mersa, martyrium sumpsit.
\switchcolumn
\selectlanguage{english}
The same day, the martyrdom of St. 
 Anthusa the Younger, who was made a martyr by being cast into a well for the 
 faith of Christ.
\switchcolumn*
\selectlanguage{latin}
Areláte, in Gállia, 
 sancti Cæsárii Epíscopi, miræ sanctitátis et pietátis viri.
\switchcolumn
\selectlanguage{english}
At Arles in France, the holy bishop 
 Caesarius, a man of great sanctity and piety.
\switchcolumn*
\selectlanguage{latin}
Augustodúni sancti 
 Syágrii, Epíscopi et Confessóris.
\switchcolumn
\selectlanguage{english}
At Autun, St. Syagrius, bishop and 
 confessor.
\switchcolumn*
\selectlanguage{latin}
Papíæ sancti Joánnis 
 Epíscopi.
\switchcolumn
\selectlanguage{english}
At Pavia, St. John, bishop.
\switchcolumn*
\selectlanguage{latin}
Ilérdæ, in Hispánia 
 Tarraconénsi, sancti Licérii Epíscopi.
\switchcolumn
\selectlanguage{english}
At Lerida in Spain, St. Licerius, 
 bishop.
\switchcolumn*
\selectlanguage{latin}
In Thebáide sancti 
 Pœmenis Anachorétæ.
\switchcolumn
\selectlanguage{english}
In Thebais, St. Poemen, abbot.
\switchcolumn*
\selectlanguage{latin}
Apud Septempedános, in 
 Picéno, sanctæ Margarítæ Víduæ.
\switchcolumn
\selectlanguage{english}
At San Severino, in Piceno, St. 
 Margaret, widow.
\switchcolumn*
\selectlanguage{latin}
\end{paracol}


% ---- martyrology/mart08/mart0828.htm
\needspace{10\baselineskip}
\begin{paracol}{2}
\selectlanguage{latin}
\begin{center}{\color{gregoriocolor} Quinto Kaléndas Septémbris. 
 Luna\dots\ }\end{center}
\switchcolumn
\selectlanguage{english}
\begin{center}{\color{gregoriocolor} The 
 28\textsuperscript{th} Day of 
 August. The\dots\ Day of the Moon.}\end{center}
\end{paracol}

\noindent\begin{tabularx}{\linewidth}{*{19}{>{\centering\arraybackslash}X}}
 \textcolor{gregoriocolor}{a} & \textcolor{gregoriocolor}{b} & \textcolor{gregoriocolor}{c} & \textcolor{gregoriocolor}{d} & \textcolor{gregoriocolor}{e} & \textcolor{gregoriocolor}{f} & \textcolor{gregoriocolor}{g} & \textcolor{gregoriocolor}{h} & \textcolor{gregoriocolor}{i} & \textcolor{gregoriocolor}{k} & \textcolor{gregoriocolor}{l} & \textcolor{gregoriocolor}{m} & \textcolor{gregoriocolor}{n} & \textcolor{gregoriocolor}{p} & \textcolor{gregoriocolor}{q} & \textcolor{gregoriocolor}{r} & \textcolor{gregoriocolor}{s} & \textcolor{gregoriocolor}{t} & \textcolor{gregoriocolor}{u} \\
 5 & 6 & 7 & 8 & 9 & 10 & 11 & 12 & 13 & 14 & 15 & 16 & 17 & 18 & 19 & 20 & 21 & 22 & 23 \\
\end{tabularx}
\vspace{0.5\baselineskip}
\noindent\begin{tabularx}{\linewidth}{*{12}{>{\centering\arraybackslash}X}}
 \textcolor{gregoriocolor}{A} & \textcolor{gregoriocolor}{B} & \textcolor{gregoriocolor}{C} & \textcolor{gregoriocolor}{D} & \textcolor{gregoriocolor}{E} & F & \textcolor{gregoriocolor}{F} & \textcolor{gregoriocolor}{G} & \textcolor{gregoriocolor}{H} & \textcolor{gregoriocolor}{M} & \textcolor{gregoriocolor}{N} & \textcolor{gregoriocolor}{P} \\
 24 & 25 & 26 & 27 & 28 & 29 & 28 & 29 & 1 & 2 & 3 & 4 \\
\end{tabularx}

\begin{paracol}{2}
\selectlanguage{latin}
\lettrine[lines=2]{H}{ippóne} Régio, in 
 Africa, natális sancti Augustíni Epíscopi, Confessóris et Ecclésiæ Doctóris 
 exímii, qui, beáti Ambrósii Epíscopi ópera ad cathólicam fidem convérsus et 
 baptizátus, eam advérsus Manichæos aliósque hæréticos acérrimus propugnátor 
 deféndit, multísque áliis pro Ecclésia Dei perfúnctus labóribus, ad præmia 
 migrávit in cælum. Ejus relíquiæ, primo de sua civitáte propter 
 bárbaros in Sardíniam advéctæ, et póstea a Rege Longobardórum Luitprándo 
 Papíam translátæ, ibi honorífice cónditæ sunt.
\switchcolumn
\selectlanguage{english}
\lettrine[lines=2]{A}{t} Hippo in Africa, the birthday of 
 St. Augustine, bishop and famous doctor of the Church. Converted and 
 baptized by the blessed bishop Ambrose, he defended the Catholic faith with 
 the greatest zeal against the Manicheans and other heretics, and after 
 having sustained many other labours for the Church of God, he went to his 
 reward in heaven. His relics, owing to the invasion of barbarians, 
 were first brought from his own city into Sardinia, and afterwards taken by Luitprand, king of the Lombards, to Pavia, where they were deposited with 
 due honours.
\switchcolumn*
\selectlanguage{latin}
Romæ item natális 
 sancti Hermétis, viri illústris, qui (ut in Actis beáti Alexándri Papæ 
 légitur), prius carceráli custódiæ mancipátus, deínde, cum áliis plúrimis, 
 gládio cædénte, martyrium complévit, sub Aureliáno Júdice.
\switchcolumn
\selectlanguage{english}
At Rome, the birthday of St. Hermes, 
 an illustrious man, who, as we read in the Acts of blessed Pope Alexander, 
 was first confined in prison, and afterwards fulfilled his martyrdom by the 
 sword, at the time of the judge Aurelian.
\switchcolumn*
\selectlanguage{latin}
Venúsiæ, in Apúlia, pássio sanctórum Septimíni, Januárii et Felícis, qui, sanctórum Bonifátii et 
 Theclæ fílii, a Valeriáno Júdice, sub Maximiáno Imperatóre, jussi sunt 
 decollári. Ipsórum tamen ac reliquórum ex duódecim frátribus 
 festívitas ágitur Kaléndis Septémbris.
\switchcolumn
\selectlanguage{english}
At Venosa in Apulia, the passion of 
 Saints Septiminus, Januarius, and Felix. During the reign of Emperor 
 Maximian, the judge Valerian ordered these sons of Saints Boniface and 
 Thecla to be beheaded. Their feast, however, is observed with that of 
 the other Twelve Holy Brethren on the first of September.
\switchcolumn*
\selectlanguage{latin}
Briváte, apud Arvérnos, 
 item pássio sancti Juliáni Mártyris, qui, cum esset beáti Ferreóli Tribúni 
 comes et in hábitu militári occúlte Christo servíret, in persecutióne 
 Diocletiáni, a milítibus tentus est, et, desécto gútture, morte horríbili 
 necátus.
\switchcolumn
\selectlanguage{english}
At Prinde in Auvergne, St. Julian, 
 martyr, during the persecution of Diocletian. He was the companion of 
 the blessed tribune Ferreol, and under a military garb he secretly served 
 Christ until arrested by the soldiers, and killed in a barbarous manner by 
 having his throat cut.
\switchcolumn*
\selectlanguage{latin}
Constántiæ, in Germánia, sancti Pelágii Mártyris, qui sub Numeriáno Imperatóre et Evilásio Júdice, 
 cápite amputátus, martyrii corónam accépit.
\switchcolumn
\selectlanguage{english}
At Constance, in Germany, St. 
 Pelagius, martyr, who was beheaded and received the crown of martyrdom under Emperor Numerian 
 and the judge Evilasius.
\switchcolumn*
\selectlanguage{latin}
Salérni sanctórum 
 Mártyrum Fortunáti, Caji et Anthis; qui, sub Diocletiáno Imperatóre et 
 Leóntio Procónsule, decolláti sunt.
\switchcolumn
\selectlanguage{english}
At Salerno, the holy martyrs 
 Fortunatus, Caius, and Anthes, beheaded under Emperor Diocletian and the 
 proconsul Leontius.
\switchcolumn*
\selectlanguage{latin}
Constantinópoli sancti 
 Alexándri Epíscopi, gloriósi senis; ob cujus oratiónem Arius, divíno judício 
 damnátus, crépuit médius, et effúsa sunt víscera ejus.
\switchcolumn
\selectlanguage{english}
At Constantinople, the holy bishop 
 Alexander, an aged and celebrated man, through whose efficacious prayers 
 Arius, by the judgment of God, burst asunder and his bowels were poured 
 out.
\switchcolumn*
\selectlanguage{latin}
Apud Sántonas, in 
 Gállia, sancti Viviáni, Epíscopi et Confessóris.
\switchcolumn
\selectlanguage{english}
At Saintes, St. Vivian, bishop and 
 confessor.
\switchcolumn*
\selectlanguage{latin}
Item sancti Móysis 
 Æthíopis, qui, ex insígni latróne insígnis Anachoréta, multos latrónes 
 convértit et secum duxit ad monastérium.
\switchcolumn
\selectlanguage{english}
Also, St. Moses the Ethiopian, who gave up a life of 
 robbery and became a renowned anchoret. He converted many robbers, and 
 led them to a monastery.
\switchcolumn*
\selectlanguage{latin}
\end{paracol}


% ---- martyrology/mart08/mart0829.htm
\needspace{10\baselineskip}
\begin{paracol}{2}
\selectlanguage{latin}
\begin{center}{\color{gregoriocolor} Quarto Kaléndas Septémbris. 
 Luna\dots\ }\end{center}
\switchcolumn
\selectlanguage{english}
\begin{center}{\color{gregoriocolor} The 
 29\textsuperscript{th} Day of 
 August. The\dots\ Day of the Moon.}\end{center}
\end{paracol}

\noindent\begin{tabularx}{\linewidth}{*{19}{>{\centering\arraybackslash}X}}
 \textcolor{gregoriocolor}{a} & \textcolor{gregoriocolor}{b} & \textcolor{gregoriocolor}{c} & \textcolor{gregoriocolor}{d} & \textcolor{gregoriocolor}{e} & \textcolor{gregoriocolor}{f} & \textcolor{gregoriocolor}{g} & \textcolor{gregoriocolor}{h} & \textcolor{gregoriocolor}{i} & \textcolor{gregoriocolor}{k} & \textcolor{gregoriocolor}{l} & \textcolor{gregoriocolor}{m} & \textcolor{gregoriocolor}{n} & \textcolor{gregoriocolor}{p} & \textcolor{gregoriocolor}{q} & \textcolor{gregoriocolor}{r} & \textcolor{gregoriocolor}{s} & \textcolor{gregoriocolor}{t} & \textcolor{gregoriocolor}{u} \\
 6 & 7 & 8 & 9 & 10 & 11 & 12 & 13 & 14 & 15 & 16 & 17 & 18 & 19 & 20 & 21 & 22 & 23 & 24 \\
\end{tabularx}
\vspace{0.5\baselineskip}
\noindent\begin{tabularx}{\linewidth}{*{12}{>{\centering\arraybackslash}X}}
 \textcolor{gregoriocolor}{A} & \textcolor{gregoriocolor}{B} & \textcolor{gregoriocolor}{C} & \textcolor{gregoriocolor}{D} & \textcolor{gregoriocolor}{E} & F & \textcolor{gregoriocolor}{F} & \textcolor{gregoriocolor}{G} & \textcolor{gregoriocolor}{H} & \textcolor{gregoriocolor}{M} & \textcolor{gregoriocolor}{N} & \textcolor{gregoriocolor}{P} \\
 25 & 26 & 27 & 28 & 29 & 30 & 29 & 1 & 2 & 3 & 4 & 5 \\
\end{tabularx}

\begin{paracol}{2}
\selectlanguage{latin}
\lettrine[lines=2]{D}{ecollátio} sancti 
 Joánnis Baptístæ, quem Heródes circa festum Paschæ decollári præcépit. 
 Ipsíus tamen memória solémniter hac die cólitur, qua venerándum ejus caput 
 secúndo invéntum fuit; quod, póstea Romam translátum, in Ecclésia sancti 
 Silvéstri, ad Campum Mártium, summa pópuli devotióne asservátur.
\switchcolumn
\selectlanguage{english}
\lettrine[lines=2]{T}{he} beheading of St. John Baptist, 
 who was put to death by Herod about the feast of Easter. However, his 
 solemn commemoration takes place today, when his venerable head was found 
 for the second time. It was afterwards solemnly carried to Rome, where 
 it is kept in the church of St. Sylvester, near the Campus Martius, and 
 honoured by the people with the greatest devotion.
\switchcolumn*
\selectlanguage{latin}
Romæ, in monte Aventíno, 
 natális sanctæ Sabínæ Mártyris, quæ sub Hadriáno Imperatóre, gládio percússa, 
 martyrii palmam adépta est.
\switchcolumn
\selectlanguage{english}
At Rome, on Mount Aventine, the 
 birthday of St. Sabina, martyr. Under Emperor Hadrian, she was struck 
 with the sword, and thus obtained the palm of martyrdom.
\switchcolumn*
\selectlanguage{latin}
Veliniáni, in 
 confínibus Apúliæ, pássio sanctórum Vitális, Satóris, et Repósiti; qui, 
 sanctórum Bonifátii et Theclæ fílii, a Valeriáno Júdice, sub Maximiáno 
 Imperatóre, capitálem senténtiam pertulérunt. Eórum tamen ac ceterórum 
 ex duódecim frátribus memória Kaléndis Septémbris recólitur.
\switchcolumn
\selectlanguage{english}
At Valiniano in Apulia, the passion 
 of Saints Vitalis, Sator, and Repositus. They were the sons of Saints 
 Boniface and Thecla, and were condemned to death by the judge Valerian in 
 the reign of Emperor Maximian. Their feast along with that of the 
 other Twelve Holy Brethren is observed on the first of September.
\switchcolumn*
\selectlanguage{latin}
Romæ sanctæ Cándidæ, 
 Vírginis et Mártyris; cujus corpus beátus Paschális Primus Papa in Ecclésiam 
 sanctæ Praxédis tránstulit.
\switchcolumn
\selectlanguage{english}
At Rome, St. Candida, virgin and 
 martyr, whose body was transferred to the Church of St. Praxedes by Pope 
 Paschal I.
\switchcolumn*
\selectlanguage{latin}
Constantinópoli 
 sanctórum Mártyrum Hypátii, Asiáni Epíscopi, et Andréæ Presbyteri; qui ambo, 
 ob cultum sanctárum Imáginum, sub Leóne Isáurico, barba pice íllita atque 
 incénsa, et cute cápitis extrácta, juguláti sunt.
\switchcolumn
\selectlanguage{english}
At Constantinople, the holy martyrs 
 Hypatius, an Asiatic bishop, and Andrew, a priest, who for the veneration of 
 holy images, under Leo the Isaurian had their beards besmirched with pitch 
 and set on fire, the skin of the heads torn off, and were beheaded.
\switchcolumn*
\selectlanguage{latin}
Antiochíæ natális 
 sanctórum Mártyrum Nicǽæ et Pauli.
\switchcolumn
\selectlanguage{english}
At Antioch, the birthday of the holy 
 martyrs Nicaeas and Paul.
\switchcolumn*
\selectlanguage{latin}
Metis, in Gállia, 
 sancti Adélphi, Epíscopi et Confessóris.
\switchcolumn
\selectlanguage{english}
At Metz in France, St. Adelphus, 
 bishop and confessor.
\switchcolumn*
\selectlanguage{latin}
Lutétiæ Parisiórum 
 deposítio sancti Mederíci Presbyteri.
\switchcolumn
\selectlanguage{english}
At Paris, the death of St. Merry, 
 priest.
\switchcolumn*
\selectlanguage{latin}
Perúsiæ sancti Euthymii 
 Románi, qui, cum uxóre et Crescéntio fílio persecutiónem Diocletiáni fúgiens, 
 ad eam urbem secéssit, et ibi póstmodum quiévit in Dómino.
\switchcolumn
\selectlanguage{english}
At Perugia, St. Euthymius, a Roman, 
 who fled from the persecution of Diocletian with this wife and his son 
 Crescentius, and there rested in the Lord.
\switchcolumn*
\selectlanguage{latin}
In Anglia sancti Sebbi 
 Regis.
\switchcolumn
\selectlanguage{english}
In England, St. Sebbe, king.
\switchcolumn*
\selectlanguage{latin}
Apud Sírmium natális 
 sanctæ Basíllæ Vírginis.
\switchcolumn
\selectlanguage{english}
At Smyrna, the birthday of St. 
 Basilla, virgin.
\switchcolumn*
\selectlanguage{latin}
In pago Tricassíno 
 sanctæ Sabínæ Vírginis, virtútibus et miráculis gloriósæ.
\switchcolumn
\selectlanguage{english}
In the vicinity of Troyes, St. 
 Sabina, a virgin, celebrated for virtues and miracles.
\switchcolumn*
\selectlanguage{latin}
\end{paracol}


% ---- martyrology/mart08/mart0830.htm
\needspace{10\baselineskip}
\begin{paracol}{2}
\selectlanguage{latin}
\begin{center}{\color{gregoriocolor} Tértio Kaléndas Septémbris. 
 Luna\dots\ }\end{center}
\switchcolumn
\selectlanguage{english}
\begin{center}{\color{gregoriocolor} The 
 30\textsuperscript{th} Day of 
 August. The\dots\ Day of the Moon.}\end{center}
\end{paracol}

\noindent\begin{tabularx}{\linewidth}{*{19}{>{\centering\arraybackslash}X}}
 \textcolor{gregoriocolor}{a} & \textcolor{gregoriocolor}{b} & \textcolor{gregoriocolor}{c} & \textcolor{gregoriocolor}{d} & \textcolor{gregoriocolor}{e} & \textcolor{gregoriocolor}{f} & \textcolor{gregoriocolor}{g} & \textcolor{gregoriocolor}{h} & \textcolor{gregoriocolor}{i} & \textcolor{gregoriocolor}{k} & \textcolor{gregoriocolor}{l} & \textcolor{gregoriocolor}{m} & \textcolor{gregoriocolor}{n} & \textcolor{gregoriocolor}{p} & \textcolor{gregoriocolor}{q} & \textcolor{gregoriocolor}{r} & \textcolor{gregoriocolor}{s} & \textcolor{gregoriocolor}{t} & \textcolor{gregoriocolor}{u} \\
 7 & 8 & 9 & 10 & 11 & 12 & 13 & 14 & 15 & 16 & 17 & 18 & 19 & 20 & 21 & 22 & 23 & 24 & 25 \\
\end{tabularx}
\vspace{0.5\baselineskip}
\noindent\begin{tabularx}{\linewidth}{*{12}{>{\centering\arraybackslash}X}}
 \textcolor{gregoriocolor}{A} & \textcolor{gregoriocolor}{B} & \textcolor{gregoriocolor}{C} & \textcolor{gregoriocolor}{D} & \textcolor{gregoriocolor}{E} & F & \textcolor{gregoriocolor}{F} & \textcolor{gregoriocolor}{G} & \textcolor{gregoriocolor}{H} & \textcolor{gregoriocolor}{M} & \textcolor{gregoriocolor}{N} & \textcolor{gregoriocolor}{P} \\
 26 & 27 & 28 & 29 & 30 & 1 & 1 & 2 & 3 & 4 & 5 & 6 \\
\end{tabularx}

\begin{paracol}{2}
\selectlanguage{latin}
\lettrine[lines=2]{S}{anctæ} Rosæ a Sancta María, e tértio Ordine sancti Domínici, Vírginis; cujus dies natális nono 
 Kaléndas Septémbris recensétur.
\switchcolumn
\selectlanguage{english}
\lettrine[lines=2]{T}{he} feast of St. Rose of St. Mary, 
 virgin of the Third Order of St. Dominic, whose birthday is recalled on the 
 24th of August.
\switchcolumn*
\selectlanguage{latin}
Romæ, via Ostiénsi, 
 pássio beáti Felícis Presbyteri, sub Diocletiáno et Maximiáno Imperatóribus. 
 Hic, post equúlei vexatiónem, data senténtia, cum ducerétur ad decollándum, 
 óbvius ei fuit quidam Christiánus, qui, dum se Christiánum esse sponte 
 profiterétur, mox cum eódem páriter decollátus est; cujus nomen ignorántes, 
 Christiáni Adáuctum eum appellavérunt, eo quod sancto Felíci auctus sit ad 
 corónam.
\switchcolumn
\selectlanguage{english}
At Rome, on the Ostian Way, the 
 martyrdom of the blessed priest Felix, under Emperors Diocletian and 
 Maximian. After being racked he was sentenced to death, and as they 
 led him to execution, he met a man who spontaneously declared himself a 
 Christian, and was forthwith beheaded with him. The Christians, not 
 knowing his name, called him Adauctus, because he was added to St. Felix and 
 shared his crown.
\switchcolumn*
\selectlanguage{latin}
Item Romæ sanctæ 
 Gaudéntiæ, Vírginis et Mártyris, cum áliis tribus.
\switchcolumn
\selectlanguage{english}
Also at Rome, St. Gaudentia, virgin 
 and martyr, with three others.
\switchcolumn*
\selectlanguage{latin}
Colóniæ Suffetulánæ, in 
 Africa, beatórum sexagínta Mártyrum, qui furóre Gentílium cæsi sunt.
\switchcolumn
\selectlanguage{english}
At Colonia Suffetulana in Africa, 
 sixty blessed martyrs, who were murdered by the furious heathen.
\switchcolumn*
\selectlanguage{latin}
Bonóniæ sancti Bonónii 
 Abbátis.
\switchcolumn
\selectlanguage{english}
At Bologna, St. Bononius, abbot.
\switchcolumn*
\selectlanguage{latin}
Romæ sancti Pammáchii 
 Presbyteri, qui fuit doctrína et sanctitáte conspícuus.
\switchcolumn
\selectlanguage{english}
At Rome, St Pammachius, priest, who 
 was noteworthy for learning and sanctity.
\switchcolumn*
\selectlanguage{latin}
Adruméti, in Africa, 
 sanctórum Bonifátii et Theclæ, qui beatórum duódecim filiórum Mártyrum 
 paréntes fuérunt.
\switchcolumn
\selectlanguage{english}
At Adrumetum, also in Africa, the 
 Saints Boniface and Thecla, who were the parents of twelve blessed sons, all 
 martyrs.
\switchcolumn*
\selectlanguage{latin}
Thessalonícæ sancti 
 Fantíni Confessóris, qui, multa a Saracénis perpéssus, atque e monastério, 
 in quo abstinéntia víxerat admirábili, expúlsus, demum, cum plúrimos ad viam 
 salútis perduxísset, in senectúte bona quiévit.
\switchcolumn
\selectlanguage{english}
At Thessalonica, St. Fantinus, 
 confessor, who suffered much from the Saracens, and was driven from his 
 monastery, in which he had lived in great abstinence. After having 
 brought many to the way of salvation, he rested at last at an advanced age.
\switchcolumn*
\selectlanguage{latin}
In território Meldénsi 
 sancti Fiácrii Confessóris.
\switchcolumn
\selectlanguage{english}
In the diocese of Meaux, St. Fiacre, 
 confessor.
\switchcolumn*
\selectlanguage{latin}
Trebis, in Látio, 
 sancti Petri Confessóris, qui, multis clarus virtútibus et miráculis, ibídem 
 migrávit ad Dóminum, et honorífice cólitur.
\switchcolumn
\selectlanguage{english}
At Trevi in Lazio, St. Peter, 
 confessor, who was distinguished for many virtues and miracles. He is 
 honoured in that place from which he departed for heaven.
\switchcolumn*
\selectlanguage{latin}
\end{paracol}


% ---- martyrology/mart08/mart0831.htm
\needspace{10\baselineskip}
\begin{paracol}{2}
\selectlanguage{latin}
\begin{center}{\color{gregoriocolor} Prídie Kaléndas Septémbris. 
 Luna\dots\ }\end{center}
\switchcolumn
\selectlanguage{english}
\begin{center}{\color{gregoriocolor} The 
 31\textsuperscript{st} Day of 
 August. The\dots\ Day of the Moon.}\end{center}
\end{paracol}

\noindent\begin{tabularx}{\linewidth}{*{19}{>{\centering\arraybackslash}X}}
 \textcolor{gregoriocolor}{a} & \textcolor{gregoriocolor}{b} & \textcolor{gregoriocolor}{c} & \textcolor{gregoriocolor}{d} & \textcolor{gregoriocolor}{e} & \textcolor{gregoriocolor}{f} & \textcolor{gregoriocolor}{g} & \textcolor{gregoriocolor}{h} & \textcolor{gregoriocolor}{i} & \textcolor{gregoriocolor}{k} & \textcolor{gregoriocolor}{l} & \textcolor{gregoriocolor}{m} & \textcolor{gregoriocolor}{n} & \textcolor{gregoriocolor}{p} & \textcolor{gregoriocolor}{q} & \textcolor{gregoriocolor}{r} & \textcolor{gregoriocolor}{s} & \textcolor{gregoriocolor}{t} & \textcolor{gregoriocolor}{u} \\
 8 & 9 & 10 & 11 & 12 & 13 & 14 & 15 & 16 & 17 & 18 & 19 & 20 & 21 & 22 & 23 & 24 & 25 & 26 \\
\end{tabularx}
\vspace{0.5\baselineskip}
\noindent\begin{tabularx}{\linewidth}{*{12}{>{\centering\arraybackslash}X}}
 \textcolor{gregoriocolor}{A} & \textcolor{gregoriocolor}{B} & \textcolor{gregoriocolor}{C} & \textcolor{gregoriocolor}{D} & \textcolor{gregoriocolor}{E} & F & \textcolor{gregoriocolor}{F} & \textcolor{gregoriocolor}{G} & \textcolor{gregoriocolor}{H} & \textcolor{gregoriocolor}{M} & \textcolor{gregoriocolor}{N} & \textcolor{gregoriocolor}{P} \\
 27 & 28 & 29 & 30 & 1 & 2 & 2 & 3 & 4 & 5 & 6 & 7 \\
\end{tabularx}

\begin{paracol}{2}
\selectlanguage{latin}
\lettrine[lines=2]{S}{ancti} Raymúndi Nonnáti, 
 ex Ordine beátæ Maríæ de Mercéde redemptiónis captivórum, Cardinális et 
 Confessóris; cujus dies natális séptimo Kaléndas Septémbris recólitur.
\switchcolumn
\selectlanguage{english}
\lettrine[lines=2]{S}{t.} Raymund Nonnatus, cardinal and 
 confessor, of the Order of our Lady of Ransom for the Redemption of 
 Captives. His birthday is commemorated on the 26th of August.
\switchcolumn*
\selectlanguage{latin}
Apud montem Senárium, 
 in Etrúria, natális sancti Bonajúnctæ Confessóris, e septem Fundatóribus 
 Ordinis Servórum beátæ Maríæ Vírginis; qui, de Passióne Domínica ad fratres 
 verba fáciens, in manus Dómini trádidit spíritum. Ipsíus autem ac 
 Sociórum festum prídie Idus Februárii celebrátur.
\switchcolumn
\selectlanguage{english}
In Tuscany, on Mount Senario, the 
 birthday of St. Bonajuncta, confessor, one of the seven founders of the 
 Order of Servites of the Blessed Virgin Mary, who gave up his soul into the 
 hands of the Lord while he was preaching to his brethren on the Passion of 
 our Saviour. his feast is kept with that of his companions on the 12th 
 of February.
\switchcolumn*
\selectlanguage{latin}
Tréviris item natális 
 sancti Paulíni Epíscopi, qui, témpore Ariánæ infestatiónis, ab Ariáno 
 Imperatóre Constántio ob cathólicam fidem relegátus exsílio, et extra 
 Christiánum nomen usque ad mortem mutándo exsília fatigátus, tandem, apud 
 Phrygiam defúnctus, beátæ passiónis corónam percépit a Dómino.
\switchcolumn
\selectlanguage{english}
At Treves, the birthday of St. 
 Paulinus, a bishop, who was exiled for the Catholic faith by the Arian 
 emperor Constantius, in the time of the Arian persecution. By having 
 to change the place of his exile, which was beyond the limits of 
 Christendom, he became wearied unto death, and finally, dying in Phrygia, 
 received a crown from the Lord for his blessed martyrdom.
\switchcolumn*
\selectlanguage{latin}
Tránsaquis, ad lacum 
 Fúcinum, in Marsis, natális quoque sanctórum Mártyrum Cæsídii Presbyteri, et 
 Sociórum; qui martyrio coronáti sunt in persecutióne Maximíni.
\switchcolumn
\selectlanguage{english}
At Transaco, in the Marches near 
 Lake Fucino, the birthday of the holy martyrs Caesidius, priest, and his 
 companions, who were crowned with martyrdom in the persecution of Maximinus.
\switchcolumn*
\selectlanguage{latin}
Item sanctórum Mártyrum 
 Robustiáni et Marci.
\switchcolumn
\selectlanguage{english}
Also, the holy martyrs Robustian and 
 Mark.
\switchcolumn*
\selectlanguage{latin}
Cæsaréæ, in Cappadócia, sanctórum Theódoti, Rufínæ et Ammiæ; quorum duo primi paréntes fuérunt 
 sancti Mamántis Mártyris, quem Rufína in cárcere péperit, et Ammia educávit.
\switchcolumn
\selectlanguage{english}
At Caesarea in Cappadocia, the 
 Saints Theodotus, Rufina, and Ammia. The first two were the parents of 
 the martyr St. Mamas, who was born in prison, and whom Ammia brought up.
\switchcolumn*
\selectlanguage{latin}
Antisiodóri sancti 
 Optáti, Epíscopi et Confessóris.
\switchcolumn
\selectlanguage{english}
At Auxerre, St. Optatus, bishop and 
 confessor.
\switchcolumn*
\selectlanguage{latin}
In Anglia sancti Aidáni, 
 Epíscopi Lindisfarnénsis; cujus ánimam, cum sanctus Cuthbértus, cujus 
 memória tertiodécimo Kaléndas Aprilis cólitur, tunc óvium pastor, in cælum 
 ferri vidísset, relíctis ovibus, factus est Mónachus.
\switchcolumn
\selectlanguage{english}
In England, St. Aidan, bishop of 
 Lindisfarne. When St. Cuthbert, then a shepherd, saw his soul going up 
 to heaven, he left his sheep and became a monk. Mention is made of St. 
 Cuthbert on the 20th of March.
\switchcolumn*
\selectlanguage{latin}
Apud Nuscum sancti 
 Amáti Epíscopi.
\switchcolumn
\selectlanguage{english}
At Nosco, St. Amatus, bishop.
\switchcolumn*
\selectlanguage{latin}
Athénis sancti 
 Aristídis, fide et sapiéntia claríssimi, qui Hadriáno Príncipi egrégium de 
 religióne Christiána volúmen óbtulit, nostri dógmatis cóntinens ratiónem; et 
 quod Christus Jesus solus esset Deus, præsénte ipso Imperatóre, 
 luculentíssime perorávit.
\switchcolumn
\selectlanguage{english}
At Athens, St. Aristides, most 
 celebrated for his faith and wisdom, who presented to Emperor Hadrian a 
 treatise on the Christian religion, containing the exposition of our 
 doctrine. In the presence of the emperor, he also delivered a 
 discourse in which he clearly demonstrated that Jesus Christ is the only God.
\switchcolumn*
\selectlanguage{latin}
\end{paracol}
